% concatenated from Slater_6_3.tex
\documentclass[reqno,12pt]{amsart}

\usepackage{amsmath, epsfig, cite}
\usepackage{amssymb}
\usepackage{comment}
\usepackage{amsfonts}
\usepackage{latexsym}
\usepackage{amsthm}
\usepackage[colorlinks, linkcolor=blue]{hyperref}
\usepackage{tabularx}
\usepackage{booktabs}
\usepackage{amsmath, epsfig, cite}
\usepackage{amssymb}
\usepackage{amsfonts}
\usepackage{latexsym}
\usepackage{extarrows}
\newtheorem{theorem}{Theorem}[section]%[section]
\newtheorem{proposition}[theorem]{Proposition}
\newtheorem{corollary}[theorem]{Corollary}
%\newtheorem{conjecture}[theorem]{Conjecture}
\newtheorem{conjecture}{Conjecture}
%\newtheorem{lemma}[theorem]{Lemma}
\newtheorem{lemma}[theorem]{Lemma}
\newtheorem{problem}[theorem]{Problem}

\theoremstyle{remark}
\newtheorem{remark}{Remark}
\newtheorem{definition}{Definition}


\setlength{\textwidth}{160.0mm} \setlength{\oddsidemargin}{0mm}
\setlength{\evensidemargin}{0mm} \addtolength{\topmargin}{-.1cm}
\addtolength{\textheight}{.2cm}

\numberwithin{equation}{section}

\newcommand{\ta}{\theta}
\newcommand{\la}{\lambda}
\newcommand{\ba}{\beta}
\newcommand{\ga}{\gamma}
\newcommand{\da}{\delta}
\newcommand{\pa}{\partial}
\newcommand{\rd}{\,\mathrm d}
\newcommand{\C}{\mathbb C}
\newcommand{\N}{\mathbb N}
\newcommand{\Z}{\mathbb Z}

\allowdisplaybreaks
\author{Chang Xu$^*$}
\author{Dunkun Yang}
\address{School of Mathematical Sciences, East China Normal University,
Shanghai 200241, P. R. China}
\email{52275500007@stu.ecnu.edu.cn, Yangdunkun@foxmail.com}
\thanks{$^*$ Corresponding author}





\title{Some $q$-transformation formulas and Rogers-Ramanujan type identities}
\subjclass[2010]{Primary 33D15; Secondary 11P83}

\keywords{basic hypergeometric series; transformation formula; summation formula; Rogers-Ramanujan type identity}

\begin{document}
\maketitle
\begin{abstract}
In this paper, we explore the role that Liu's transformation formula can play in discovering Rogers-Ramanujan type identities.  
Specifically, we combine Liu's transformation formula with other $q$-series summations to derive a series of parameterized identities.
Thereafter, through careful selection of these  parameters,  we obtain  $75$ Rogers-Ramanujan type identities from Slater's list, and uncover several new Rogers-Ramanujan type identities.%, a new Appell-Lerch series identity and a theta-function identity that was proposed by Srivastava, Chaudhary and  Wakene. 
% 
%
%Slater applied  different  Bailey pairs to find $130$ Rogers-Ramanujan type identities. In this paper, without using Bailey pairs, we give a new  proof of Slater's identities.  employ an elegant transformation formula for $q$-series, developed by Liu, to derive a series of parameterized identities. Through careful selection of these  parameters, we uncover numerous Rogers-Ramanujan type identities, including $72$ identities listed by Slater. 
\end{abstract}
%
\section{Introduction}
Throughout the paper, we assume that $|q|<1$ for convergence and define the $q$-shifted factorials as
\begin{align*}
\left(a;q\right)_{0}=1,
\quad\left(a;q\right)_{n}=\prod_{k=0}^{n-1}\left(1-aq^k\right)
,\quad
n=1,2,\cdots, \mathrm{or}\, {\infty}.\end{align*}
Also,  following the definition of Slater \cite{Slater}, we define
$\left(q;q\right)_{-1}=1$.
For brevity, we abbreviate a collection of several  $q$-shifted factorials on the same base $q$ as
\begin{equation*}
\left(a_1,a_2,\cdots,a_m;q\right)_n
=\left(a_1;q\right)_{n}\left(a_2;q\right)_{n}\cdots\left(a_m;q\right)_{n}.
\end{equation*} 

Following Gasper and Rahman's definition \cite{Gasper_and_Rahman}, the $_r\phi_s(\cdot)$ basic hypergeometric series is defined as
\begin{equation*}
{_r\phi_s}\left(
\begin{array}{cccc}
a_1,\,a_2,\,\cdots,a_r\\
b_1,\,b_2,\,\cdots,b_s
\end{array}
;q,\,z
\right)
=
\sum_{n=0}^{\infty}
\frac{\left(a_1,\,a_2,\,\cdots,a_r;q\right)_n}
{\left(q,\,b_1,\,b_2,\,\cdots,b_s;q\right)_n}
\left(\left(-1\right)^nq^{n(n-1)/2}\right)^{1+s-r}z^n
\end{equation*}
where $q\neq0$ when $r>s+1$.

Many fantastic sum-product $q$-sereis were uncovered by Rogers \cite{Rogers}, including
\begin{align}
&\sum_{n=0}^{\infty}\frac{q^{n^2}}{\left(q;q\right)_n}
=\frac{1}{\left(q,q^4;q^5\right)_{\infty}},\label{q1}\\
&
\sum_{n=0}^{\infty}\frac{q^{n(n+1)}}{\left(q;q\right)_n}
=\frac{1}{\left(q^2,q^3;q^5\right)_{\infty}}.\label{q2}
\end{align}
However, Rogers's work didn't attract people's attention for some time until Ramanujan independently rediscovered  the identities \eqref{q1}-\eqref{q2}. Thereafter, the identities \eqref{q1}-\eqref{q2} were usually named as the Rogers-Ramanujan identities. Later, MacMahon \cite{MacMahon} gave combinatorial interpretations.  Schur \cite{Schur} independently rediscovered \eqref{q1} and \eqref{q2} along with their combinatorial interpretation.

Rogers-Ramanujan identities play crucial roles in various branches of mathematics and physics. For example, these identities have a close relation to combinatorics, vertex algebras, knot theory, and statistical mechanics.
Therefore, discovering identities similar to \eqref{q1}-\eqref{q2}, commonly  referred to as Rogers-Ramanujan type identities, has attracted considerable attention from experts and led to significant progress. One of the celebrated  works concerning this topic is Slater's list \cite{Slater} which includes $130$ of such identities, such as
\cite[Eqs. (58), (59)]{Slater}
\begin{align*}
\sum_{n=0}^{\infty}\frac{\left(q^2;q^2\right)_{n-1}}
{\left(q;q\right)_{n-1}\left(q;q\right)_n\left(q;q^2\right)_n}q^{n^2}
&{=}\frac{\left(-q^5,-q^7,q^{12};q^{12}\right)_{\infty}}{\left(q;q\right)_{\infty}},
\quad{\tag{S.58}}
\\
\sum_{n=0}^{\infty}\frac{q^{n^2+2n}}
{\left(q;q\right)_n\left(q;q^2\right)_{n+1}}
&{=}\frac{\left(q^2,q^{12},q^{14};q^{14}\right)_{\infty}}
{\left(q;q\right)_{\infty}}.\quad {\tag{S.59}}
\end{align*}
Here and forth, we use the label $\left(\mathrm{S}.n\right)$ to denote the equation $(n)$ in Slater's list \cite{Slater}.

Slater used the Bailey lemma to  discovery these identities. Also, more Rogers-Ramanujan type identities  were uncovered by applying the Bailey lemma. The reader may be referred to \cite{Andrews2,Bowman,McLaughlin-2008, McLaughlin_2008}.   In this paper,  we shall show how to obtain Rogers-Ramanujan type identities  without using the Bailey lemma. 
The key tool in our proof is the following elegant transformation formula for $q$-series introduced by Liu \cite{Liu2}.
 \begin{theorem}\label{theorem1}
 For a complex sequence $\left\{A_n\right\}$, under suitable convergence conditions, there holds
 \begin{align}
 &\frac{\left(\alpha q,\alpha ab/q;q\right)_{\infty}}
 {\left(\alpha a,\alpha b;q\right)_{\infty}}
 \sum_{n=0}^{\infty}A_n\left(q/a;q\right)_n\left(\alpha a\right)^n\nonumber\\
 &\quad=\sum_{n=0}^{\infty}
 \frac{\left(1-\alpha q^{2n}\right)\left(\alpha,q/a,q/b;q\right)_n}{\left(1-\alpha\right)\left(q,\alpha a,\alpha b;q\right)_n}\left(-\alpha ab\right)^nq^{n(n-3)/2}\nonumber\\
 &\qquad\times\sum_{k=0}^{n}\frac{\left(q^{-n},\alpha q^n;q\right)_k}{\left(q/b;q\right)_k}\left(q^2/b\right)^kA_k.
 \label{Formula3}
  \end{align}
 \end{theorem}
Specifically,
we take
\begin{equation*}
A_n=\frac{\left(q/b,\beta,\gamma,\lambda;q\right)_n}
{\left(q, c, d,e,h;q\right)_n}\left(b/q\right)^n
\end{equation*}
into Theorem \ref{theorem1} to obtain the following  transformation formula.
\begin{theorem}\label{thm1}
For $\max\left\{|c|,|d|,|e|,|h|,|\alpha a|,|\alpha b|,|\alpha ab/q| \right\}<1$, we have
\begin{align}
&\frac{\left(\alpha q,\alpha ab/q;q\right)_{\infty}}{\left(\alpha a,\alpha b;q\right)_{\infty}}
{_5\phi_4}\left(\begin{array}{ccccc}
q/a,\,q/b,\,\beta,\,\gamma,\,\lambda\\
c,\,d,\,e,\,h
\end{array}
;q,{\alpha ab}/{q}\right)\nonumber\\
&\quad=
\sum_{n=0}^{\infty}
\frac{\left(1-\alpha q^{2n}\right)\left(\alpha,q/a,q/b;q\right)_n}
{\left(1-\alpha\right)\left(q,\alpha a,\alpha b;q\right)_n}\left(-\alpha ab\right)^nq^{n\left(n-3\right)/2}\nonumber\\
&\qquad\times{_5\phi_4}\left(\begin{array}{ccccc}
q^{-n},\,\alpha q^n,\,\beta,\,\gamma,\,\lambda\\
c,\,d,\,e,\,h
\end{array}
;q,q\right).\label{Formula4}
\end{align}
\end{theorem}
By combining  \eqref{Formula4} with other well-known $q$-series summation formulas and  the Jacobi triple product identity  \cite[Eq. (1.6.1)]{Gasper_and_Rahman}
\begin{equation}\label{Jacobi}
\sum_{n=-\infty}^{\infty}q^{\frac{n^2-n}{2}}z^n
=\left(-z,-q/z,q;q\right)_{\infty},
\end{equation}
we derive a number of  parameter-dependent identities in Sections \ref{sec2}-\ref{sec7}.  In particular, by making specific choices for the parameters, we obtain numerous  Rogers-Ramanujan type identities, including $75$ identities listed by Slater and  the following ones that appear to be new,
\begin{align*}
%\sum_{n=0}^{\infty}\frac{\left(-q^2;q^2\right)_n}{\left(q;q\right)_{2n}}q^{n^2-2n}&\!=\!\frac{\left(-1/q;q^2\right)_{\infty}}{\left(q;q^2\right)_{\infty}},\\
\sum_{n=0}^{\infty}\frac{1}{\left(q;q\right)_{2n+1}}q^{n^2}&=
\frac{\left(q^8,q^{12},q^{20};q^{20}\right)_{\infty}+q\left(q^4,q^{16},q^{20};q^{20}\right)_{\infty}}
{\left(q;q\right)_{\infty}\left(-q^2;q^2\right)_{\infty}},
\\
%\sum_{n=0}^{\infty}\frac{\left(-q,-q;q^2\right)_n}{\left(q;q\right)_{2n+1}\left(-q^2;q^2\right)_{n}}q^{n\left(n+2\right)}
%&\!=\!\left(-q;q\right)_{\infty}\left(-q^4,-q^4;q^4\right)_{\infty},\\
%\sum_{n=0}^{\infty}\frac{\left(q;q\right)_{3n+1}\left(-q^3;q^3\right)_n}{\left(q^3;q^3\right)_n\left(q^3;q^3\right)_{2n+1}}q^{\frac{{3n^2+3n}}{2}}
%&\!=\!\frac{\left(-q^3;q^3\right)_{\infty}
%\left(q,q^5,q^6;q^6\right)_{\infty}}{\left(q^3;q^3\right)_{\infty}},\\\sum_{n=0}^{\infty}\frac{\left(-1\right)^n\left(q;q^2\right)_{2n}}{\left(-q^2;q^4\right)_{n}\left(q^8;q^8\right)_{n}}q^{2n^2}&\!
%=\!\frac{\left(q^2;q^4\right)_{\infty}\left(-q^3,-q^5,q^{8};q^{8}\right)_{\infty}}{\left(q^4;q^4\right)_{\infty}},\\
\sum_{n=0}^{\infty}\frac{\left(q;q\right)_{3n}\left(-q^3;q^3\right)_n}
{\left(q^3;q^3\right)_n\left(q^3;q^3\right)_{2n}}q^{\frac{3n^2-3n}{2}}
&\!=\!\frac{\left(-q^3;q^3\right)_{\infty}}{\left(q^3;q^3\right)_{\infty}}
\left(\left(q,q^5,q^6;q^6\right)_{\infty}+\left(q^2,q^4,q^6;q^6\right)_{\infty}\right),\\
\sum_{n=0}^{\infty}\frac{1}{\left(-q;q^2\right)_{n}\left(q;q\right)_{2n+1}}q^{2n^2+2n}
&\!=\!\frac{\left(q^8,q^{20},q^{28};q^{28}\right)_{\infty}+
q
\left(q^4,q^{24},q^{28};q^{28}\right)_{\infty}
}{\left(q^2;q^2\right)_{\infty}}.
\end{align*}
\begin{comment}
We also give a proof of the theta-function identity which was given by \cite{Hari} without proving:
\begin{align*}
\sum_{n=0}^{\infty}\frac{\left(-1\right)^n
q^{\frac{n\left(n+3\right)}{2}}}{\left(q;q\right)_n
\left(1+q^{n+1}\right)^2}
=
\frac{\varphi(-q)}
{q}\sum_{n=1}^{\infty}\frac{q^n}{1+q^n},
\end{align*}
 where the Ramanujan's theta function $\varphi(q)$ satisfies
\begin{align*}
\varphi(q)=\sum_{n=-\infty}^{\infty}q^{n^2}
=\frac{\left(q;q\right)_{\infty}}{\left(-q;q\right)_{\infty}}.
\end{align*} 
In addition, our approach yields a new Appell-Lerch series identity as follows,
\begin{align*}
\sum_{n=0}^{\infty}\frac{\left(-1\right)^n\left(q^2;q^4\right)_nq^{n\left(n+2\right)}}
{\left(q;q\right)_{2n+1}\left(-q^2;q^2\right)_{n}}%=\frac{\left(q;q^2\right)_{\infty}}
%{\left(q^2;q^2\right)_{\infty}}
&=\frac{1}{\psi(q)}\sum_{n=-\infty}^{\infty}
\left(-1\right)^n\frac{q^{2n^2+2n}}{1-q^{2n+1}},
\end{align*}
where  the  Ramanujan's theta fuction $\psi(q)$ is defined as
\begin{align*}
\psi(q)=\sum_{n=0}^{\infty}{q^{n(n+1)/2}}=\frac{\left(q^2;q^2\right)_{\infty}}
{\left(q;q^2\right)_{\infty}}.
\end{align*}
The
Appell-Lerch series are series of the form
$$\sum_{n=-\infty}^{\infty}\frac{\left(-1\right)^{ln}q^{ln\left(n+1\right)/2}b^n}
{1-aq^n}.$$
Obviously, the summation of the right hand-side of the above identity is an Appell-Lerch series with $\left(q,a,b,l\right)\to\left(q^2,q,-1,2\right)$.  

%\begin{align*}
%\sum_{n=0}^{\infty}\frac{q^{2n^2+2n}}{\left(-q;q^2\right)_{n+1}\left(q;q\right)_{2n}}
%=\frac{q
%\left(q^4,q^{24},q^{28};q^{28}\right)_{\infty}-\left(q^8,q^{20},q^{28};q^{28}\right)_{\infty}
%}{\left(q^2;q^2\right)_{\infty}}.
%\end{align*}
\end{comment}

Our method is constructive, which means that we are able to  derive the product side of a sum-product $q$-series identity from its given sum side, without requiring prior knowledge of product side.  The same methodology has been applied  to the  study of identities on false theta functions, Hecke-type series and Appell-Lerch series, as documented in \cite{Wang_Chen,WangChun,Wang,Wang_Yee}.


%Our method is constructive, which means that we can deduce the product part from the given sum part of a sum-product $q$-series identity.
%The paper is organized as follows. In Section \ref{sec2}, we apply \eqref{Formula4} with three famous $_2\phi_1$ transformation formulas  including $q$-Chu-Vandermonde summation formula.
 %Then, in Section \ref{sec3}, we combing \eqref{Formula4} with $q$-Pfaff-Saalsch\"utz summation}


%\section{Preliminaries}
%In this section, we collect some identities on $q$-series in the literature for later use.



\section{Identities from three ${_2\phi_1}$ summations}\label{sec2}
When  $\lambda=\gamma=h=e$, Theorem \ref{thm1} reduces to  the following result which was first proposed by Liu in \cite{Liu_1} and later applied in \cite{Wang, Wang_Chen,Wang_Yee}: for $|\alpha ab/q|<1$,
\begin{align}
&\frac{\left(\alpha q,\alpha ab/q;q\right)_{\infty}}{\left(\alpha a,\alpha b;q\right)_{\infty}}
{_3\phi_2}\left(\begin{array}{ccc}
q/a,\,q/b,\,\beta\\
c,\,d
\end{array}
;q,{\alpha ab}/{q}\right)\nonumber\\
&\quad=
\sum_{n=0}^{\infty}
\frac{\left(1-\alpha q^{2n}\right)\left(\alpha,q/a,q/b;q\right)_n}
{\left(1-\alpha\right)\left(q,\alpha a,\alpha b;q\right)_n}
\left(-\alpha ab\right)^nq^{{n\left(n-3\right)}/{2}}
{_3\phi_2}\left(\begin{array}{ccc}
q^{-n},\,\alpha q^n,\,\beta\\
c,\,d
\end{array}
;q,q\right).\label{Formula1}
\end{align}
Now, we use the $q$-Chu-Vandermonde summation formula and  the transformation \eqref{Formula1} to derive the following transformation formula.
\begin{theorem}\label{thm9}
For $\max\left\{|\alpha a|,|\alpha b|,|c|,|\alpha ab/q^2|\right\}<1$, we have
\begin{align*}
&\frac{\left(\alpha q^2,\alpha ab/q^2;q^2\right)_{\infty}}
{\left(\alpha a,\alpha b;q^2\right)_{\infty}}
\sum_{n=0}^{\infty}\frac{\left(q^2/a,q^2/b;q^2\right)_n}
{\left(q^2,c;q^2\right)_n}\left(\alpha ab/q^2\right)^n\\
&\quad=\sum_{n=0}^{\infty}\frac{1-\alpha q^{4n}}{1-\alpha}
\frac{\left(\alpha,q^2/a,q^2/b,\alpha q^2/c;q^2\right)_n}
{\left(q^2,\alpha a,\alpha b,c;q^2\right)_n}\alpha^n a^nb^nc^nq^{2n^2-4n}.
\end{align*}
\end{theorem}
\begin{proof}
We begin by recalling  the $q$-Chu-Vandermonde summation formula \cite[Eq. (1.5.3)]{Gasper_and_Rahman}
\begin{equation}\label{q-Chu}
_2\phi_1\left(
\begin{array}{cc}
q^{-n},\,a\\
c
\end{array};q,\,q\right)=\frac{\left(c/a;q\right)_{n}}{\left(c;q\right)_n}a^n.
\end{equation}
Replacing $q$ by $q^2$ and setting $\beta=d$ in \eqref{Formula1}, and invoking \eqref{q-Chu} with $a=\alpha q^{2n}$, we arrive at Theorem \ref{thm9}.
\end{proof}

\begin{corollary}
We have
\begin{align}
\sum_{n=0}^{\infty}\frac{\left(-1\right)^nq^{n^2}}{\left(q^2;q^2\right)_n}
&{=}\left(q;q^2\right)_{\infty},\tag{S.3 and S.23}\label{S.3}\\
%&=
%\frac{\left(q;q^2\right)_{\infty}\left(q^2,q^4,q^6;q^6\right)_{\infty}}
%{\left(q^2;q^2\right)_{\infty}},\tag{S.23}\label{S.23}\\
\sum_{n=0}^{\infty}\frac{\left(-1\right)^n\left(-q;q^2\right)_n}
{\left(q^4;q^4\right)_n}q^{n^2}
&{=}\left(q;q^2\right)_{\infty}
\left(q^2;q^4\right)_{\infty}.\tag{S.4}\label{S.4}
%\sum_{n=0}^{\infty}\frac{\left(-1\right)^nq^{2n^2}}{\left(-1,q^2;q^2\right)_n}
%&=
%\frac{\left(-q^2,-q^6,q^8;q^8\right)_{\infty}+1}
%{\left(-1;q^2\right)_{\infty}},\label{e8}\\
%\sum_{n=0}^{\infty}\frac{q^{2n^2+n}}{\left(q;q\right)_{2n}}
%&=\frac{\left(q,q,q^2;q^2\right)_{\infty}-1}{2\left(q;q^2\right)_{\infty}}.\label{e7}
\end{align}
\end{corollary}
\begin{proof}
When $\alpha=-1, c=-q^2$, Theorem \ref{thm9} reduces to
\begin{align}
\sum_{n=0}^{\infty}\frac{\left(q^2/a,q^2/b;q^2\right)_{n}}{\left(q^2,-q^2;q^2\right)_{n}}
\left(-ab/q^2\right)^n=\frac{\left(-a,-b;q^2\right)_{\infty}}{\left(-q^2,-ab/q^2;q^2\right)_{\infty}}, \label{A14}
\end{align}
where $\max\{|a|,|b|,|ab/q^2|\}<1$. 

By setting $\left(a,b\right)\to\left(0,0\right)$ in \eqref{A14}, we obtain 
\begin{align*}
\sum_{n=0}^{\infty}\frac{\left(-1\right)^nq^{2n^2}}{\left(q^4;q^4\right)_n}=\frac{1}{\left(-q^2;q^2\right)_{\infty}}.
\end{align*}
Then, replacing $q^2$ by $q$ in the above equation and using the relation
\begin{equation*}
\left(q^2;q^2\right)_{\infty}=\left(q^2,q^4,q^6;q^6\right)_{\infty},
\end{equation*}
we obtain \eqref{S.3}.

 Also, by taking $\left(a,b\right)\to\left(0,-q\right)$ in \eqref{A14},  we get \eqref{S.4}.
 \end{proof}
 % Also, by taking $\left(a,b,c\right)\to\left(0,0,-1\right)$ into \eqref{A1} and applying the Jacobi triple product identity \eqref{Jacobi}
%with $\left(q,z\right)\to\left(q^8,q^2\right)$,
%we arrive at  \eqref{e8}.
\begin{corollary}
We have
\begin{align}
\sum_{n=0}^{\infty}\frac{\left(-1\right)^nq^{n\left(2n+1\right)}}
{\left(q^2;q^2\right)_n\left(-q;q^2\right)_{n+1}}
&=\left(q,-q^2;q^2\right)_{\infty}
,\tag{S.5}\label{S.5}\\
\sum_{n=0}^{\infty}\frac{q^{n\left(2n+1\right)}}{\left(q;q\right)_{2n+1}}
&=\frac{\left(-q,-q^3,q^4;q^4\right)_{\infty}}{\left(q^2;q^2\right)_{\infty}}
\tag{S.9}\label{S.9}\\
&=\frac{\left(q^2,q^6,q^8;q^8\right)_{\infty}
\left(q^4,q^{12};q^{16}\right)_{\infty}}{\left(q;q\right)_{\infty}}
,\tag{S.84}\label{S.84}\\
\sum_{n=0}^{\infty}\frac{\left(-q;q^2\right)_n}
{\left(q;q\right)_{2n+1}}q^{n\left(n+1\right)}&=
\frac{\left(q^4,q^8,q^{12};q^{12}\right)_{\infty}}{\left(q;q\right)_{\infty}}, \tag{S.51}\label{S.51}
\\
\sum_{n=0}^{\infty}\frac{\left(q^2;q^4\right)_n\left(-q^2;q^2\right)_n
q^{n\left(n+1\right)}}{\left(q^4;q^4\right)_n\left(q;q^2\right)_n
\left(q;q^2\right)_{n+1}}
&=
\frac{\left(-q^2;q^2\right)_{\infty}\left(-q,-q^3,q^4;q^4\right)_{\infty}}
{\left(q^2;q^2\right)_{\infty}},\tag{S.64}\label{S.64}\\
\sum_{n=0}^{\infty}\frac{\left(-q^2;q^2\right)_n}
{\left(q;q\right)_{2n+1}}q^{n^2}
&=\frac{\left(-q;q^2\right)_{\infty}}{\left(q;q^2\right)_{\infty}}.\tag{\cite[Entry 1.7.13]{Ramanujan}}\label{N1}
\end{align}
\end{corollary}
\begin{proof}
Substituting $\alpha=-q,c=-q^3$ into Theorem \ref{thm9} results in
\begin{align}
\sum_{n=0}^{\infty}\frac{\left(q^2/a,q^2/b;q^2\right)_{n}}
{\left(q^2,-q^3;q^2\right)_{n}}
\left(-ab/q\right)^n=\frac{\left(-aq,-bq;q^2\right)_{\infty}}{\left(-q^3,-ab/q;q^2\right)_{\infty}},\label{A2}
\end{align}
where $\max\{|aq|,|bq|,|ab/q|\}<1$. 

After inserting $\left(a,b\right)\to\left(0,0\right)$ into \eqref{A2}, 
we obtain \eqref{S.5}.  Furthermore, replacing $q$ by $-q$ in \eqref{S.5} and utilizing the following relation
\begin{equation*}
\frac{1}{\left(q;q^2\right)_{\infty}}
=\frac{\left(-q,-q^3,q^4;q^4\right)_{\infty}}{\left(q^2;q^2\right)_{\infty}}
=\frac{\left(q^2,q^6,q^8;q^8\right)_{\infty}
\left(q^4,q^{12};q^{16}\right)_{\infty}}{\left(q;q\right)_{\infty}},
\end{equation*}
we  obtain \eqref{S.9} and \eqref{S.84}.
  
  
In addition, setting $\left(a,b\right)\to\left(0,q\right)$ into \eqref{A2}, we get
\begin{align}
\sum_{n=0}^{\infty}\frac{\left(q;q^2\right)_n}{\left(q^2,-q^3;q^2\right)_n}q^{n(n+1)}
=\frac{\left(-q^2;q^2\right)_{\infty}}{\left(-q^3;q^2\right)_{\infty}}.\label{C7}
\end{align}
Replacing $q$ by $-q$ in \eqref{C7}, we have
\begin{align*}
\sum_{n=0}^{\infty}\frac{\left(-q;q^2\right)_n}{\left(q;q\right)_{2n+1}}q^{n(n+1)}
=\frac{\left(-q^2;q^2\right)_{\infty}}{\left(q;q^2\right)_{\infty}}.
\end{align*}
Then, combining the above equation with the following two relations
\begin{align*}
&\sum_{n=0}^{\infty}\frac{\left(q^2;q^4\right)_n\left(-q^2;q^2\right)_n}{\left(q^4;q^4\right)_n\left(q;q^2\right)_n
\left(q;q^2\right)_{n+1}}q^{n\left(n+1\right)}
=\sum_{n=0}^{\infty}\frac{\left(-q;q^2\right)_n
}{\left(q;q\right)_{2n+1}}q^{n\left(n+1\right)},\\
&\frac{1}{\left(q;q^2\right)_{\infty}}=\frac{\left(q^2;q^2\right)_{\infty}}{\left(q;q\right)_{\infty}}
=\frac{\left(-q,-q^3,q^4;q^4\right)_{\infty}}{\left(q^2;q^2\right)_{\infty}},
%=\frac{\left(q^2,q^6,q^8;q^8\right)_{\infty}
%\left(q^4,q^{12};q^{16}\right)_{\infty}}{\left(q;q\right)_{\infty}},
\end{align*}
we arrive at \eqref{S.51} and \eqref{S.64}.

Finally, replacing $q$ by $-q$ in \eqref{A2} and taking $\left(a,b\right)\to\left(0,-1\right)$, we get \ref{N1}.
%with $\left(q,z\right)\to\left(q^2,q\right)$, we get
%\begin{align*}
%\left(-q;q^2\right)_{\infty}\sum_{n=0}^{\infty}
%\frac{\left(-1\right)^nq^{2n^2+n}}{\left(-q,q^2;q^2\right)_n}
%=\sum_{n=0}^{\infty}\left(1+q^{4n+1}\right)q^{4n^2}=\sum_{n=0}^{\infty}q^{n^2}=
%\frac{\left(-q,-q,q^2;q^2\right)_{\infty}-1}{2}.
%\end{align*}
%Then, replacing $q$ by $-q$, we arrive at \eqref{e7}.
\end{proof}
%






%Similarly, setting $\left(a,b\right)\to\left(0,-q\right)$ in \eqref{q5} and applying the following two relations
%\begin{align*}
%\sum_{n=0}^{\infty}\frac{\left(q^2;q^4\right)_n\left(-q^2;q^2\right)_n
%}{\left(q^4;q^4\right)_n\left(q;q^2\right)_n
%\left(q;q^2\right)_{n+1}}q^{n\left(n+1\right)}
%&=\sum_{n=0}^{\infty}\frac{\left(-q;q^2\right)_n
%}{\left(q;q\right)_{2n+1}}q^{n\left(n+1\right)},\\
%\frac{\left(-q^2;q^2\right)_{\infty}\left(-q,-q^3,q^4;q^4\right)_{\infty}}
%{\left(q^2;q^2\right)_{\infty}}
%=\left(-q,-q^2,-q^2;q^2\right)_{\infty}
%=\frac{\left(q^4;q^4\right)_{\infty}}{\left(q;q\right)_{\infty}}
%&=
%\frac{\left(q^4,q^8,q^{12};q^{12}\right)_{\infty}}{\left(q;q\right)_{\infty}}
%,
%\end{align*}
%we arrive at \eqref{S.51} and \eqref{S.64}.
%Additionally, letting $\left(a,b\right)\to\left(0,-1\right)$ in \eqref{q5}, we obtain \eqref{e6}.

\begin{corollary}
We have
\begin{align}
\sum_{n=0}^{\infty}\frac{\left(-1;q\right)_{n}}
{\left(q;q\right)_n}q^{\frac{1}{2}n\left(n+1\right)}
&=\frac{\left(-q;q\right)_{\infty}\left(q^2,q^2,q^4;q^4\right)_{\infty}}
{\left(q;q\right)_{\infty}},\tag{S.12}\label{S.12}\\
\sum_{n=0}^{\infty}\frac{\left(-q;q\right)_{n}}
{\left(q;q\right)_n}q^{\frac{1}{2}n\left(n-1\right)}
&=\frac{\left(-q;q\right)_{\infty}}
{\left(q;q\right)_{\infty}}\left(\left(q,q^3,q^4;q^4\right)_{\infty}
+\left(q^2,q^2,q^4;q^4\right)_{\infty}\right),\tag{S.13}\label{S.13}\\
\sum_{n=0}^{\infty}\frac{q^{n^2}}
{\left(q;q\right)_n}
&=\frac{\left(q^2,q^3,q^5;q^5\right)_{\infty}}
{\left(q;q\right)_{\infty}},\tag{S.18}\label{S.18}\\
\sum_{n=0}^{\infty}\frac{\left(-1;q\right)_{2n}}
{\left(q^2;q^2\right)_n}q^{n}
&=\frac{\left(-q;q\right)_{\infty}\left(q^3,q^3,q^6;q^6\right)_{\infty}}
{\left(q;q\right)_{\infty}},\tag{S.24}\label{S.24}\\
\sum_{n=0}^{\infty}\frac{\left(-1\right)^n\left(-1,q;q^2\right)_n}
{\left(q^2;q^2\right)_n}q^{n}
&=\frac{\left(q,-q^2;q^2\right)_{\infty}\left(-q^3,-q^3,q^6;q^6\right)_{\infty}}
{\left(-q,q^2;q^2\right)_{\infty}},\tag{S.30}\label{S.30}\\
\sum_{n=0}^{\infty}\frac{\left(-q;q^2\right)_{n}}
{\left(q^2;q^2\right)_{n}}q^{n^2}
&=\frac{\left(-q^3;q^2\right)_{\infty}\left(q^3,q^5,q^8;q^8\right)_{\infty}}
{\left(q^2;q^2\right)_{\infty}}.\tag{S.36}\label{S.36}
\end{align}
\end{corollary}

\begin{proof}
Taking $\alpha=1,c=0$ into Theorem \ref{thm9}, we deduce that
\begin{align}
&\frac{\left(q^2,ab/q^2;q^2\right)_{\infty}}
{\left(a,b;q^2\right)_{\infty}}\sum_{n=0}^{\infty}
\frac{\left(q^2/a,q^2/b;q^2\right)_{n}}
{\left(q^2;q^2\right)_n}\left(ab/q^2\right)^n\nonumber\\
&\quad=1+\sum_{n=1}^{\infty}\left(-1\right)^n\left(1+q^{2n}\right)
\frac{\left(q^2/a,q^2/b;q^2\right)_{n}}{\left(a,b;q^2\right)_n}
a^nb^nq^{3n^2-3n},\label{q6}
\end{align}
where $\max\{|a|,|b|,|ab/q^2|\}<1$.

When taking $\left(a,b\right)\to\left(0,-q^2\right)$ into \eqref{q6} and employing the Jacobi triple product identity \eqref{Jacobi}
with $\left(q,z\right)\to\left(q^8,-q^4\right)$, we deduce that
\begin{equation}\label{equ8}
\frac{\left(q^2;q^2\right)_{\infty}}{\left(-q^2;q^2\right)_{\infty}}
\sum_{n=0}^{\infty}\frac{\left(-1;q^2\right)_{n}}{\left(q^2;q^2\right)_{n}}
q^{n\left(n+1\right)}
=\left(q^4,q^4,q^8;q^8\right)_{\infty}.
\end{equation}
Then, we complete the proof of \eqref{S.12} by letting $q^2\to q$ in \eqref{equ8}.

Moreover, we take $\left(a,b\right)\to\left(0,-1\right)$ into  \eqref{q6} to get
\begin{align}
&\frac{\left(q^2;q^2\right)_{\infty}}{\left(-q^2;q^2\right)_{\infty}}
\sum_{n=0}^{\infty}\frac{\left(-q^2;q^2\right)_{n}}{\left(q^2;q^2\right)_{n}}
q^{n\left(n-1\right)}\nonumber\\
&\quad=\sum_{n=-\infty}^{\infty}\left(-1\right)^nq^{4n^2-2n}
+\sum_{n=-\infty}^{\infty}\left(-1\right)^nq^{4n^2}\nonumber\\
&\quad=\left(q^2,q^6,q^8;q^8\right)_{\infty}
+\left(q^4,q^4,q^8;q^8\right)_{\infty},\label{equ9}
\end{align}
where the last identity follows from the Jacobi triple product identity \eqref{Jacobi} with $\left(q,z\right)\to\left(q^8,-q^2\right)$ and $\left(q,z\right)\to\left(q^8,-q^4\right)$, respectively.
Then, we arrive at \eqref{S.13} by letting $q^2\to q$ in \eqref{equ9}.

Taking $\left(a,b\right)\to\left(0,0\right)$ into \eqref{q6} and employing the Jacobi triple product identity \eqref{Jacobi} with $\left(q,z\right)\to\left(q^{10},-q^4\right)$, we obtain
\begin{equation}\label{Equ8}
\left(q^2;q^2\right)_{\infty}
\sum_{n=0}^{\infty}\frac{q^{2n^2}}{\left(q^2;q^2\right)_{n}}
=\left(q^4,q^6,q^{10};q^{10}\right)_{\infty}.
\end{equation}
Then, we derive \eqref{S.18} by letting $q^2\to q$ in \eqref{Equ8}.

In addition, taking $\left(a,b\right)\to\left(-q,-q^2\right)$, $\left(a,b\right)\to\left(q,-q^2\right)$, $\left(a,b\right)\to\left(-q,0\right)$, respectively, into  \eqref{q6} and employing the Jacobi triple product identity \eqref{Jacobi}, we obtain \eqref{S.24}, \eqref{S.30} and  \eqref{S.36}.
\end{proof}
\begin{corollary}
We have
\begin{align}
\sum_{n=0}^{\infty}\frac{\left(-q;q^2\right)_n}{\left(q;q\right)_{2n}}q^{n^2}
&=
\frac{\left(-q;q^2\right)_{\infty}\left(-q^2,-q^4,q^6;q^6\right)_{\infty}}
{\left(q^2;q^2\right)_{\infty}},\tag{S.29}\label{S.29}\\
\sum_{n=0}^{\infty}\frac{q^{2n^2}}{\left(q;q\right)_{2n}}
&=
\frac{\left(-q^3,-q^5,q^8;q^8\right)_{\infty}}
{\left(q^2;q^2\right)_{\infty}}.\tag{S.39}\label{S.39}
\end{align}
\end{corollary}
\begin{proof}
Taking $\alpha=1,c=q$ into Theorem \ref{thm9}, we deduce that
\begin{align}
&\frac{\left(q^2,ab/q^2;q^2\right)_{\infty}}
{\left(a,b;q^2\right)_{\infty}}\sum_{n=0}^{\infty}
\frac{\left(q^2/a,q^2/b;q^2\right)_{n}}
{\left(q,q^2;q^2\right)_n}\left(ab/q^2\right)^n\nonumber\\
&\quad=1+\sum_{n=1}^{\infty}\left(1+q^{2n}\right)
\frac{\left(q^2/a,q^2/b;q^2\right)_{n}}{\left(a,b;q^2\right)_n}
a^nb^nq^{2n^2-3n},\label{A15}
\end{align}
where $\max\{|a|,|b|,|ab/q^2|\}<1$.

Substituting $\left(a,b\right)\to\left(0,-q\right)$ into \eqref{A15} and using the Jacobi triple product with $\left(q,z\right)\to\left(q^6,q^2\right)$, we derive \eqref{S.29}.



Additionally, taking $\left(a,b\right)\to\left(0,0\right)$ into  \eqref{A15} and using Jacobi triple product with $\left(q,z\right)\to\left(q^8,q^3\right)$, we obtain \eqref{S.39}.
\end{proof}
\begin{corollary}
We have
\begin{align}
\sum_{n=0}^{\infty}\frac{\left(-q;q^2\right)_{n}}
{\left(q^4;q^4\right)_n}q^{n^2}
&=\frac{\left(-q;q^2\right)_{\infty}\left(q^3,q^3,q^6;q^6\right)_{\infty}}
{\left(q^2;q^2\right)_{\infty}}.\tag{S.25}\label{S.25}
%&\sum_{n=0}^{\infty}\frac{\left(-1;q^2\right)_nq^{n\left(n+1\right)}}{\left(q;q\right)_{2n}}
%=\frac{\left(-q^2;q^2\right)_{\infty}\left(-q^3,-q^3,q^6;q^6\right)_{\infty}}
%{\left(q^2;q^2\right)_{\infty}},\quad \mathrm{\left(S. 48\right)}\label{S.48}\\
%&\sum_{n=0}^{\infty}\frac{\left(q;q^2\right)_{n}\left(-1\right)^nq^{n^2}}
%{\left(q^2;q^2\right)_{n}}
%=\frac{\left(q;q^2\right)_{\infty}\left(-q^3,-q^5,q^8;q^8\right)_{\infty}}
%{\left(q^2;q^2\right)_{\infty}}.\label{new1}
\end{align}
\end{corollary}
\begin{proof}
Taking $\alpha=1,c=-q^2$ into Theorem \ref{thm9}, we deduce that
\begin{align}
&\frac{\left(q^2,ab/q^2;q^2\right)_{\infty}}
{\left(a,b;q^2\right)_{\infty}}\sum_{n=0}^{\infty}
\frac{\left(q^2/a,q^2/b;q^2\right)_{n}}
{\left(q^2,-q^2;q^2\right)_n}\left(ab/q^2\right)^n\nonumber\\
&\quad=1+2\sum_{n=1}^{\infty}\left(-1\right)^n
\frac{\left(q^2/a,q^2/b;q^2\right)_{n}}{\left(a,b;q^2\right)_n}
a^nb^nq^{2n^2-2n},\label{q61}
\end{align}
where $\max\{|a|,|b|,|ab/q^2|\}<1$.

Inserting $\left(a,b\right)\to\left(0,-q\right)$ into \eqref{q61} and employing the Jacobi  triple product identity \eqref{Jacobi} with $\left(q,z\right)\to\left(q^6,-q^3\right)$, we get \eqref{S.25}.
%Additionally, taking $\left(a,b,c\right)\to\left(0,-q^2,q\right)$ into  \eqref{q6} and using  Jacobi triple product with $\left(q,z\right)\to\left(q^6, q^3\right)$, we obtain \eqref{S.48}.
%Finally, making $\left(a,b,c\right)\to\left(q,0,0\right)$ into \eqref{q6} and employing the Jacobi triple product identity \eqref{Jacobi} with $\left(q,z\right)\to\left(q^8,q^3\right)$, we obtain \eqref{new1} which seems to be new.
\end{proof}
\begin{corollary}
We have
\begin{align}
\sum_{n=0}^{\infty}\frac{\left(-q;q\right)_n
}{\left(q;q\right)_n}q^{\frac{1}{2}n\left(n+1\right)}
&=
\frac{\left(-q;q\right)_{\infty}\left(q,q^3,q^4;q^4\right)_{\infty}}
{\left(q;q\right)_{\infty}},\tag{S.8}\label{S.8}\\
\sum_{n=0}^{\infty}\frac{q^{n\left(n+1\right)}}
{\left(q;q\right)_n}
&=\frac{\left(q,q^4,q^5;q^5\right)_{\infty}}
{\left(q;q\right)_{\infty}},\tag{S.14}\label{S.14}\\
\sum_{n=0}^{\infty}\frac{\left(-q;q^2\right)_n}
{\left(q^2;q^2\right)_n}q^{n\left(n+2\right)}
&=\frac{\left(-q;q^2\right)_{\infty}\left(q,q^7,q^8;q^8\right)_{\infty}}
{\left(q^2;q^2\right)_{\infty}}.\tag{S.34} \label{S.34}
\end{align}
\end{corollary}
\begin{proof}
By taking $\alpha=q^2,c=0$ into Theorem \ref{thm9}, we derive 
\begin{align}
&\frac{\left(q^2,ab;q^2\right)_{\infty}}
{\left(aq^2,bq^2;q^2\right)_{\infty}}\sum_{n=0}^{\infty}
\frac{\left(q^2/a,q^2/b;q^2\right)_{n}}
{\left(q^2;q^2\right)_{n}}a^nb^n\nonumber\\
&\quad=\sum_{n=0}^{\infty}\left(1-q^{4n+2}\right)
\frac{\left(q^2/a,q^2/b;q^2\right)_{n}}
{\left(aq^2,bq^2;q^2\right)_n}
a^nb^nq^{3n^2+n},\label{A17}
\end{align}
where $\max\{|ab|,|aq^2|,|bq^2|\}<1$.

Taking $\left(a,b\right)\to\left(0,-1\right)$ into \eqref{A17} and employing the Jacobi triple product identity \eqref{Jacobi} with $\left(q,z\right)\to\left(q^8,-q^2\right)$, we obtain
\begin{equation}\label{equa8}
\frac{\left(q^2;q^2\right)_{\infty}}{\left(-q^2;q^2\right)_{\infty}}
\sum_{n=0}^{\infty}\frac{\left(-q^2;q^2\right)_{n}}{\left(q^2;q^2\right)_{n}}
q^{n\left(n+1\right)}
=\left(q^2,q^6,q^8;q^8\right)_{\infty}.
\end{equation}
Then, we derive \eqref{S.8} by letting $q^2\to q$ in \eqref{equa8}.

In addition, taking $\left(a,b\right)\to\left(0,0\right)$ into  \eqref{A17} and applying the Jacobi triple product identity \eqref{Jacobi} with $\left(q,z\right)\to\left(q^{10},-q^2\right)$, we obtain
\begin{equation}\label{equa7}
\left(q^2;q^2\right)_{\infty}
\sum_{n=0}^{\infty}\frac{q^{2n\left(n+1\right)}}{\left(q^2;q^2\right)_{n}}
=\left(q^2,q^8,q^{10};q^{10}\right)_{\infty}.
\end{equation}
Then, we complete the proof of \eqref{S.14} by letting $q^2\to q$ in \eqref{equa7}.

Taking $\left(a,b\right)\to\left(0,-q\right)$ into \eqref{A17} and utilizing the Jacobi triple product identity \eqref{Jacobi} with $\left(q,z\right)\to\left(q^8,-q\right)$, we immediately obtain \eqref{S.34}.
\end{proof}
\begin{corollary}
We have
\begin{align}
\sum_{n=0}^{\infty}\frac{q^{\frac{n\left(n+1\right)}{2}}}
{\left(q;q\right)_{n}}
&=\left(-q;q\right)_{\infty},\tag{S.2}\label{S.2}\\
\sum_{n=0}^{\infty}\frac{
q^{n\left(n+1\right)}}{\left(q^2;q^2\right)_n}
&=
\frac{\left(q,q^3,q^4;q^4\right)_{\infty}}
{\left(q;q\right)_{\infty}},\tag{S.7}\label{S.7}
\\
\sum_{n=0}^{\infty}\frac{\left(-q;q^2\right)_n}
{\left(q^4;q^4\right)_n}q^{n\left(n+2\right)}
&=\frac{\left(-q;q^2\right)_{\infty}\left(q,q^5,q^6;q^6\right)_{\infty}}
{\left(q^2;q^2\right)_{\infty}}.\tag{\cite[Entry 5.3.7]{Ramanujan}}\label{new2}
\end{align}
\end{corollary}
\begin{proof}
By inserting $\alpha=q^2,c=-q^2$ into Theorem \ref{thm9}, we find that 
\begin{align}
&\frac{\left(q^2,ab;q^2\right)_{\infty}}
{\left(aq^2,bq^2;q^2\right)_{\infty}}\sum_{n=0}^{\infty}
\frac{\left(q^2/a,q^2/b;q^2\right)_{n}}
{\left(q^2,-q^2;q^2\right)_{n}}a^nb^n\nonumber\\
&\quad=\sum_{n=0}^{\infty}\left(-1\right)^n\left(1-q^{4n+2}\right)
\frac{\left(q^2/a,q^2/b;q^2\right)_{n}}
{\left(aq^2,bq^2;q^2\right)_n}
a^nb^nq^{2n^2},\label{A16}
\end{align}
where $\max\{|ab|,|aq^2|,|bq^2|\}<1$.

Letting $\left(a,b\right)\to\left(0,0\right)$ in \eqref{A16} and utilizing the Jacobi triple product identity \eqref{Jacobi} with $\left(q,z\right)\to\left(q^8,-q^2\right)$, we deduce that
\begin{equation*}
\sum_{n=0}^{\infty}\frac{q^{2n\left(n+1\right)}}{\left(q^4;q^4\right)_n}
=\frac{\left(q^2,q^6,q^8;q^8\right)_{\infty}}{\left(q^2;q^2\right)_{\infty}}=\left(-q^4;q^4\right)_{\infty},
\end{equation*}
which can be reduced to \eqref{S.2} and  \eqref{S.7} by taking $q^4\to q$ and $q^2\to q$, respectively.

In addition, taking $\left(a,b\right)\to\left(0,-q\right)$ into  \eqref{A16}, we deduce that
\begin{align*}
\frac{\left(q^2;q^2\right)_{\infty}}{\left(-q;q^2\right)_{\infty}}
\sum_{n=0}^{\infty}
\frac{\left(-q;q^2\right)_{n}}
{\left(q^2,-q^2;q^2\right)_{n}}q^{n\left(n+2\right)}
=\sum_{n=0}^{\infty}\left(-1\right)^n
\left(1-q^{2n+1}\right)q^{3n^2+2n}
=\left(q,q^5,q^6;q^6\right)_{\infty},
\end{align*}
where the last equation follows from  the Jacobi triple product identity \eqref{Jacobi} with $\left(q,z\right)\to\left(q^6,-q\right)$. Hence, we complete the proof of \ref{new2}.
\end{proof}
\begin{corollary}
We have
\begin{align}
\sum_{n=0}^{\infty}\frac{\left(-q^2;q^2\right)_n}
{\left(q;q\right)_{2n+1}}q^{n\left(n+1\right)}
&=\frac{\left(-q^2;q^2\right)_{\infty}\left(-q,-q^5,q^6;q^6\right)_{\infty}}
{\left(q^2;q^2\right)_{\infty}},\tag{S.28}\label{S.28}\\
\sum_{n=0}^{\infty}\frac{q^{2n\left(n+1\right)}}
{\left(q;q\right)_{2n+1}}
&=\frac{\left(-q,-q^7,q^8;q^8\right)_{\infty}}
{\left(q^2;q^2\right)_{\infty}}\tag{S.38} \label{S.38}\\
&=\frac{\left(q^3,q^5,q^8;q^8\right)_{\infty}
\left(q^2,q^{14};q^{16}\right)_{\infty}}{\left(q;q\right)_{\infty}}, \tag{S.86}\label{S.86}\\
\sum_{n=0}^{\infty}\frac{\left(-q;q^2\right)_n
}{\left(q;q\right)_{2n+1}}q^{n\left(n+2\right)}
&=
\frac{\left(q^2,q^{10},q^{12};q^{12}\right)_{\infty}}
{\left(q;q\right)_{\infty}}.\tag{S.50}\label{S.50}
\end{align}
\end{corollary}
\begin{proof}
By taking $\alpha=q^2,c=q^3$ into Theorem \ref{thm9}, we derive 
\begin{align}
&\frac{\left(q^2,ab;q^2\right)_{\infty}}
{\left(aq^2,bq^2;q^2\right)_{\infty}}\sum_{n=0}^{\infty}
\frac{\left(q^2/a,q^2/b;q^2\right)_{n}}
{\left(q;q\right)_{2n+1}}a^nb^n\nonumber\\
&\quad=\sum_{n=0}^{\infty}\left(1+q^{2n+1}\right)
\frac{\left(q^2/a,q^2/b;q^2\right)_{n}}
{\left(aq^2,bq^2;q^2\right)_n}
a^nb^nq^{2n^2+n},\label{q7}
\end{align}
where $\max\{|ab|,|aq^2|,|bq^2|\}<1$.



Inserting $\left(a,b\right)\to\left(0,-1\right)$ into \eqref{q7} and employing the Jacobi triple product identity \eqref{Jacobi} with $\left(q,z\right)\to\left(q^6,q\right)$, we obtain \eqref{S.28}.
%\begin{equation*}
%\frac{\left(q^2;q^2\right)_{\infty}}
%{\left(-q^2;q^2\right)_{\infty}}
%\sum_{n=0}^{\infty}\frac{\left(-q^2;q^2\right)_{n}q^{n\left(n+1\right)}}
%{\left(q;q\right)_{2n+1}}
%=\sum_{n=0}^{\infty}\left(1+q^{2n+1}\right)q^{3n^2+2n}
%=\left(-q,-q^5,q^6;q^6\right),
%\end{equation*}
%which emplies \eqref{S.28} is true.


Taking $\left(a,b\right)\to\left(0,0\right)$ into \eqref{q7} and employing the Jacobi triple product identity \eqref{Jacobi} with $\left(q,z\right)\to\left(q^8,q\right)$, we obtain
\begin{equation*}
\left(q^2;q^2\right)_{\infty}
\sum_{n=0}^{\infty}\frac{q^{n\left(n+2\right)}}{\left(q;q\right)_{2n+1}}
%=\sum_{n=0}^{\infty}\left(1+q^{2n+1}\right)q^{4n^2+3n}
=\left(-q,-q^7,q^8;q^8\right)_{\infty}
=\frac{\left(q^3,q^5,q^8;q^8\right)_{\infty}
\left(q^2,q^{14};q^{16}\right)_{\infty}}
{\left(q;q^2\right)_{\infty}},
\end{equation*}
which means \eqref{S.38} and \eqref{S.86} are true.


Taking $\left(a,b\right)\to\left(0,-q\right)$ into \eqref{q7} and employing the Jacobi triple product identity \eqref{Jacobi} with $\left(q,z\right)\to\left(q^6,q^6\right)$, we obtain
\begin{equation*}
\frac{\left(q^2;q^2\right)_{\infty}}{\left(-q;q^2\right)_{\infty}}
\sum_{n=0}^{\infty}\frac{\left(-q;q^2\right)_{n}}{\left(q;q\right)_{2n+1}}
q^{n\left(n+2\right)}
=\sum_{n=0}^{\infty}q^{3n^2+3n}
=\left(-q^6,-q^6,q^{6};q^{6}\right)_{\infty},
\end{equation*}
which implies that \eqref{S.50} is true.
\end{proof}
%
%\subsection{Proofs of Slater's (47) and (52)}
\begin{corollary}
We have
\begin{align}
\sum_{n=0}^{\infty}\frac{\left(-1;q^2\right)_n}
{\left(q;q\right)_{2n}}q^{n^2}&=
\frac{\left(-q;q^2\right)_{\infty}\left(q^2,q^6;q^{8}\right)_{\infty}
\left(q^4;q^4\right)_{\infty}}
{\left(q;q\right)_{\infty}},
\tag{S.47}\label{S.47}\\
\sum_{n=0}^{\infty}\frac{\left(-q;q^2\right)_n}
{\left(q^2;q^4\right)_{n}\left(q^2;q^2\right)_{n}}q^{n\left(2n-1\right)}
&=
\frac{\left(-q;q^2\right)_{\infty}\left(q^4,q^8,q^{12};q^{12}\right)_{\infty}}
{\left(q^2;q^2\right)_{\infty}}.
\tag{S.52}\label{S.52}
%\sum_{n=0}^{\infty}\frac{\left(-q^2;q^2\right)_n}{\left(q;q\right)_{2n}}q^{n^2-2n}&=\frac{\left(-1/q;q^2\right)_{\infty}}{\left(q;q^2\right)_{\infty}}.\label{N2}
\end{align}
\end{corollary}
\begin{proof}
We make the substitution $\alpha=1/q, c=q$ into Theorem \ref{thm9}
to obtain
\begin{align}
\sum_{n=0}^{\infty}\frac{\left(q^2/a,q^2/b;q^2\right)_{n}}
{\left(q,q^2;q^2\right)_n}\left(ab/q^3\right)^n
=\frac{\left(a/q,b/q;q^2\right)_{\infty}}
{\left(q,ab/q^3;q^2\right)_{\infty}},\label{q8}
\end{align}
where $\max\{|a/q|, |b/q|, |ab/q^3|\}<1$.

We observe that
\begin{equation*}
\frac{\left(-q;q^2\right)_{\infty}}{\left(q;q^2\right)_{\infty}}=
\frac{\left(-q;q^2\right)_{\infty}\left(q^2,q^6;q^{8}\right)_{\infty}
\left(q^4;q^4\right)_{\infty}}
{\left(q;q\right)_{\infty}}.
\end{equation*}
Then, by substituting $\left(a,b\right)\to\left(0,-q^2\right)$ into \eqref{q8}, we get \eqref{S.47}.

Similarly, in \eqref{q8},  when we set $\left(a,b\right)\to\left(0,0\right)$,  the result simplifies to  \eqref{S.52}.
\end{proof}
%\subsection{Proofs of Slater's (45) and (60)}
Next, with the help of a ${_2\phi_1}$ summation formula found in \cite[Eq. (6.125)]{Harsh}:
\begin{align}
&{_2\phi_1}\left(
\begin{array}{cc}
a,\,b\\
aq^2/b
\end{array};q,\,-q/b\right)=\frac{\left(-q;q\right)_{\infty}\left(
\left(aq,aq^2/b^2;q^2\right)_{\infty}-q/b
\left(a,aq^3/b^2;q^2\right)_{\infty}\right)}{\left(1-q/b\right)
\left(aq^2/b,-q/b;q\right)_{\infty}},\label{C1}
\end{align} 
we establish the following transformation formula. 
\begin{theorem}\label{T7}
For $\max\{|a^2|,|aq^2|,|bq^2|,|ab|\}<1$, we have
\begin{align*}
&\frac{\left(q^4,ab;q^2\right)_{\infty}}
{\left(-q^2,a^2,bq^2;q^2\right)_{\infty}\left(q^2;q^4\right)_{\infty}}
\sum_{n=0}^{\infty}\frac{\left(q^2/a,q^2/b;q^2\right)_n}
{\left(q^2,q^3,-q^3;q^2\right)_n}a^nb^n\\
&\quad =\sum_{n=0}^{\infty}\left(-1\right)^n
\frac{\left(q^2/a,q^2/b;q^2\right)_{2n}}{\left(aq^2,bq^2;q^2\right)_{2n}}
a^{2n}b^{2n}q^{6n^2+2n}\\
&\qquad-
\sum_{n=0}^{\infty}\left(-1\right)^n
\frac{\left(q^2/a,q^2/b;q^2\right)_{2n+1}}{\left(aq^2,bq^2;q^2\right)_{2n+1}}
a^{2n+1}b^{2n+1}q^{6n^2+10n+4}.
\end{align*}
\end{theorem}
\begin{proof}
First, replacing $q$ by $q^2$ and substituting $\left(a,b\right)\to\left(q^{-2n},q^{1-2n}\right)$ into \eqref{C1}, we have
\begin{align}
&{_2\phi_1}\left(
\begin{array}{cc}
q^{-2n},\,q^{1-2n}\\
q^3
\end{array};q^2,\,-q^{2n+1}\right)\nonumber\\
&\quad=\frac{\left(-q^2;q^2\right)_{\infty}\left(
\left(q^{2-2n},q^{2+2n};q^4\right)_{\infty}-q^{1+2n}
\left(q^{-2n},q^{4+2n};q^4\right)_{\infty}\right)}{\left(1-q^{2n+1}\right)
\left(q^3,-q^{2n+1};q^2\right)_{\infty}}.\label{eq1}
\end{align}
On the other hand, from \cite[p.71]{Gasper_and_Rahman}, we find 
\begin{equation}\label{C2}
_3\phi_2\left(
\begin{array}{ccc}
q^{-n},\,b,\,0\\
d,\,e
\end{array};q,\,q\right)
=\frac{\left(-e\right)^nq^{{\frac{n\left(n-1\right)}{2}}}}{\left(e;q\right)_n}
{_2\phi_1}\left(
\begin{array}{cc}
q^{-n},\,d/b\\
d
\end{array};q,\,bq/e\right).
\end{equation}
Thereafter, setting $\left(q,b,d,e\right)\to\left(q^2,q^{2n+2},q^3,-q^3\right)$ in \eqref{C2} and applying \eqref{eq1}, we get
\begin{align}
&{_3\phi_2}\left(
\begin{array}{ccc}
q^{-2n},\,q^{2n+2},\,0\\
q^3,-q^3
\end{array};q^2,\,q^2\right)\nonumber\\
&\quad=\frac{q^{n^2+2n}}{\left(-q^3;q^2\right)_n}{_2\phi_1}\left(
\begin{array}{cc}
q^{-2n},\,q^{1-2n}\\
q^3
\end{array};q^2,\,-q^{2n+1}\right)\nonumber\\
&\quad=\frac{\left(1-q^2\right)\left(-q^2;q^2\right)_{\infty}}
{\left(1-q^{4n+2}\right)\left(q^2;q^4\right)_{\infty}}q^{n^2+2n}\left(
\left(q^{2-2n},q^{2+2n};q^4\right)_{\infty}-q^{1+2n}
\left(q^{-2n},q^{4+2n};q^4\right)_{\infty}\right).\label{eq2}
\end{align}
Then, by inserting $\left(q,\alpha,\beta,c,d\right)\to\left(q^2,q^2,0,q^3,-q^3\right)$ into \eqref{Formula1}, we deduce that
\begin{align*}
&\frac{\left(q^2;q^4\right)_{\infty}\left(q^4,ab;q^2\right)_{\infty}}
{\left(-q^2,a^2,bq^2;q^2\right)_{\infty}}
\sum_{n=0}^{\infty}\frac{\left(q^2/a,q^2/b;q^2\right)_n}
{\left(q^2,q^3,-q^3;q^2\right)_n}a^nb^n\\
&\quad=\sum_{n=0}^{\infty}\left(-1\right)^n
\frac{\left(q^2/a,q^2/b;q^2\right)_n}{\left(aq^2,bq^2;q^2\right)_n}a^nb^nq^{2n^2+n}
\left(
\left(q^{2-2n},q^{2+2n};q^4\right)_{\infty}-q^{1+2n}
\left(q^{-2n},q^{4+2n};q^4\right)_{\infty}\right)\\
&\quad=\left(q^2;q^4\right)^2_{\infty}\left(\sum_{n=0}^{\infty}\left(-1\right)^n
\frac{\left(q^2/a,q^2/b;q^2\right)_{2n}}{\left(aq^2,bq^2;q^2\right)_{2n}}
a^{2n}b^{2n}q^{6n^2+2n}\right.\\
&\qquad\left.-
\sum_{n=0}^{\infty}\left(-1\right)^n
\frac{\left(q^2/a,q^2/b;q^2\right)_{2n+1}}{\left(aq^2,bq^2;q^2\right)_{2n+1}}
a^{2n+1}b^{2n+1}q^{6n^2+10n+4}\right),
\end{align*}
which is what we want to show.
\end{proof}
\begin{corollary}
We have
\begin{align}
\sum_{n=0}^{\infty}\frac{\left(-q;q\right)_nq^{\frac{n^2+n}{2}}}{\left(q;q^2\right)_{n+1}
\left(q;q\right)_n}
&=\frac{\left(-q;q\right)_{\infty}\left(q^3,q^{7},q^{10};q^{10}\right)_{\infty}}{\left(q;q\right)_{\infty}}
,\tag{S.45}\label{S.45}\\
\sum_{n=0}^{\infty}\frac{q^{n^2+n}}{\left(q;q^2\right)_{n+1}\left(q;q\right)_n}
&=\frac{\left(q^4,q^{10},q^{14};q^{14}\right)_{\infty}}{\left(q;q\right)_{\infty}}
,\tag{S.60}\label{S.60}\\
\sum_{n=0}^{\infty}\frac{q^{n^2}}{\left(q;q\right)_{2n+1}}&=
\frac{\left(q^8,q^{12},q^{20};q^{20}\right)_{\infty}+q\left(q^4,q^{16},q^{20};q^{20}\right)_{\infty}}{\left(q;q\right)_{\infty}\left(-q^2;q^2\right)_{\infty}}.\label{A21}
\end{align}
\end{corollary}
\begin{proof}
Taking $\left(a,b\right)\to\left(0,-1\right)$ into Theorem \ref{T7} and using the Jacobi triple product identity \eqref{Jacobi} with $\left(q,z\right)\to\left(q^{20},-q^6\right)$, we derive
\begin{align*}
\frac{\left(q^4;q^2\right)_{\infty}}
{\left(q^2;q^4\right)_{\infty}\left(-q^2;q^2\right)^2_{\infty}}
\sum_{n=0}^{\infty}\frac{\left(-q^2;q^2\right)_nq^{n^2+n}}
{\left(q^2,q^3,-q^3;q^2\right)_{n}}
=\sum_{n=-\infty}^{\infty}\left(-1\right)^nq^{10n^2+4n}
=\left(q^6,q^{14},q^{20};q^{20}\right)_{\infty}.
\end{align*}
Then, replacing $q^2$ by $q$, we arrive at \eqref{S.45}.

Inserting $\left(a,b\right)\to\left(0,0\right)$ into Theorem \ref{T7} and applying the Jacobi triple product identity \eqref{Jacobi} with $\left(q,z\right)\to\left(q^{28},-q^8\right)$, we obtain
\begin{align*}
\frac{\left(q^4;q^2\right)_{\infty}}
{\left(q^2;q^4\right)_{\infty}\left(-q^2;q^2\right)_{\infty}}
\sum_{n=0}^{\infty}\frac{q^{2n^2+2n}}{\left(q^2,q^3,-q^3;q^2\right)_{n}}
=\sum_{n=-\infty}^{\infty}\left(-1\right)^nq^{14n^2+6n}
=\left(q^8,q^{20},q^{28};q^{28}\right)_{\infty}.
\end{align*}
Thereby, we prove \eqref{S.60} is true by replacing $q^2$ by $q$ in the above identity.

In addition, letting $\left(a,b\right)\to\left(0,-1/q\right)$ into Theorem \ref{T7}, we get
\begin{align*}
&\frac{\left(q^4;q^2\right)_{\infty}}
{\left(-q^2,-q^3;q^4\right)_{\infty}\left(q^2;q^4\right)_{\infty}}
\sum_{n=0}^{\infty}\frac{q^{n^2}}{\left(q^2,q^3;q^2\right)_{n}}\\
&\quad=\sum_{n=-\infty}^{\infty}\left(-1\right)^nq^{10n^2+2n}+q\sum_{n=-\infty}^{\infty}\left(-1\right)^nq^{10n^2+6n}\\
&\quad=\left(q^8,q^{12},q^{20};q^{20}\right)_{\infty}+q\left(q^4,q^{16},q^{20};q^{20}\right)_{\infty},
\end{align*}
which reduces to \eqref{A21} after simplification.
\end{proof}
%\subsection{Proof of Slater's (96)}
Another interesting  $_2\phi_1$ summation formula can also be  found in \cite[Eq. (6.110)]{Harsh} as follows,
\begin{align}
&{_2\phi_1}\left(
\begin{array}{cc}
a,\,bq\\
aq/b
\end{array};q,\,-q/b\right)=\frac{b\left(-q;q\right)_{\infty}\left(
\left(a,aq/b^2;q^2\right)_{\infty}-
\left(aq,a/b^2;q^2\right)_{\infty}\right)}{a\left(1-b\right)
\left(aq/b,-1/b;q\right)_{\infty}}.\label{C3}
\end{align}
Now we employ \eqref{C3} to deduce the following transformation formula.
\begin{theorem}\label{T9}
For $\max\left\{|aq^2|,|bq^2|,|ab|\right\}<1$, we have
\begin{align*}
&\frac{\left(q^2,ab;q^2\right)_{\infty}}{\left(aq^2,bq^2;q^2\right)_{\infty}}
\sum_{n=0}^{\infty}\frac{\left(q^2/a,q^2/b;q^2\right)_n}
{\left(-q,q^2;q^2\right)_n\left(q;q^2\right)_{n+1}}a^nb^n\\
&\quad=\sum_{n=0}^{\infty}\left(-1\right)^n\left(1+q^{4n+1}\right)
\frac{\left(q^2/a,q^2/b;q^2\right)_{2n}}{\left(aq^2,bq^2;q^2\right)_{2n}}
\left(ab\right)^{2n}q^{6n^2+2n}\\
&\qquad -\sum_{n=0}^{\infty}\left(-1\right)^n\left(1+q^{4n+3}\right)
\frac{\left(q^2/a,q^2/b;q^2\right)_{2n+1}}{\left(aq^2,bq^2;q^2\right)_{2n+1}}
\left(ab\right)^{2n+1}q^{6n^2+6n+1}
.
\end{align*}
\end{theorem}
\begin{proof}
Replacing $q$ by $q^2$ and setting $a=q^{-2n}, b=q^{-1-2n}$ in \eqref{C3}, we have
\begin{align*}
&{_2\phi_1}\left(\begin{array}{cc}
q^{-2n},\,q^{1-2n}\\
q^3
\end{array};q^2,\,-q^{2n+3}
\right)\\
&\quad=\frac{\left(-q^2;q^2\right)_{\infty}q^{2n}}{\left(1-q^{2n+1}\right)
\left(q^3,-q^{1+2n};q^2\right)_{\infty}}
\left(\left(q^{2n+2},q^{2-2n};q^4\right)_{\infty}
-\left(q^{-2n},q^{4+2n};q^4\right)_{\infty}\right).
\end{align*}
Based on the above identity, we substitute $\left(q,b,d,e\right)\to\left(q^2,q^{2n+2},q^3,-q\right)$ into \eqref{C2} to obtain
\begin{align*}
&{_3\phi_2}\left(\begin{array}{ccc}
q^{-2n},\,q^{2n+2},\,0\\
q^3,\,-q
\end{array};q^2,\,q^2
\right)\\
&\quad=\frac{q^{n^2}}{\left(-q;q^2\right)_n}
{_2\phi_1}\left(\begin{array}{cc}
q^{-2n},\,q^{1-2n}\\
q^3
\end{array};q^2,\,-q^{2n+3}
\right)\\
&\quad=\frac{\left(-q^2;q^2\right)_{\infty}q^{n^2+2n}}{\left(1-q^{2n+1}\right)
\left(q^3,-q;q^2\right)_{\infty}}
\left(\left(q^{2n+2},q^{2-2n};q^4\right)_{\infty}
-\left(q^{-2n},q^{4+2n};q^4\right)_{\infty}\right).
\end{align*}
Therefore, taking $\left(q,\alpha,\beta,c,d\right)\to\left(q^2,q^2,0,q^3,-q\right)$ into \eqref{Formula1}, we have
\begin{align*}
&\frac{\left(q^2,ab;q^2\right)_{\infty}}{\left(aq^2,bq^2;q^2\right)_{\infty}}
\sum_{n=0}^{\infty}\frac{\left(q^2/a,q^2/b;q^2\right)_n}
{\left(-q,q^2,q^3;q^2\right)_n}a^nb^n\\
&\quad=
\sum_{n=0}^{\infty}\left(-1\right)^n\left(1-q^{4n+2}\right)
\frac{\left(q^2/a,q^2/b;q^2\right)_n}{\left(aq^2,bq^2;q^2\right)_n}a^nb^nq^{n^2-n}
{_3\phi_2}\left(\begin{array}{ccc}
q^{-2n},\,q^{2n+2},\,0\\
q^3,\,-q
\end{array};q^2,\,q^2
\right)\\
&\quad=\frac{\left(-q^2;q^2\right)_{\infty}}{\left(q^3,-q;q^2\right)_{\infty}}
\sum_{n=0}^{\infty}\left(-1\right)^n\left(1+q^{2n+1}\right)
\frac{\left(q^2/a,q^2/b;q^2\right)_n}{\left(aq^2,bq^2;q^2\right)_n}
a^nb^nq^{2n^2+n}\\
&\qquad\times\left(\left(q^{2n+2},q^{2-2n};q^4\right)_{\infty}
-\left(q^{-2n},q^{4+2n};q^4\right)_{\infty}\right)\\
&\quad=\left(1-q\right)\left(-q^2;q^2\right)_{\infty}\left(q^2;q^4\right)_{\infty}
\left(
\sum_{n=0}^{\infty}\left(-1\right)^n\left(1+q^{4n+1}\right)
\frac{\left(q^2/a,q^2/b;q^2\right)_{2n}}{\left(aq^2,bq^2;q^2\right)_{2n}}
\left(ab\right)^{2n}q^{6n^2+2n}\right.\\
&\qquad\left. -
\sum_{n=0}^{\infty}\left(-1\right)^n\left(1+q^{4n+3}\right)
\frac{\left(q^2/a,q^2/b;q^2\right)_{2n+1}}{\left(aq^2,bq^2;q^2\right)_{2n+1}}
\left(ab\right)^{2n+1}q^{6n^2+6n+1}
\right)
\end{align*}
as desired.
\end{proof}
\begin{corollary}
We have
\begin{align}
\sum_{n=0}^{\infty}\frac{q^{n^2+2n}}{\left(q;q\right)_{2n+1}}
&=\frac{\left(q^4,q^6,q^{10};q^{10}\right)_{\infty}
\left(q^2,q^{18};q^{20}\right)_{\infty}}{\left(q;q\right)_{\infty}}, \tag{S.96}\label{S.96}\\
\sum_{n=0}^{\infty}\frac{q^{2n^2+2n}}{\left(-q;q^2\right)_{n}\left(q;q\right)_{2n+1}}
&=\frac{\left(q^8,q^{20},q^{28};q^{28}\right)_{\infty}+
q
\left(q^4,q^{24},q^{28};q^{28}\right)_{\infty}
}{\left(q^2;q^2\right)_{\infty}}
.\label{new7}
\end{align}
\end{corollary}
\begin{proof}
By taking $\left(a,b\right)\to\left(0,-q\right)$ into Theorem \ref{T9} and using the Jacobi triple product identity \eqref{Jacobi} with $\left(q,z\right)\to\left(q^{20},-q^4\right)$, we can prove \eqref{S.96} as follows,
\begin{align*}
\sum_{n=0}^{\infty}\frac{q^{n^2+2n}}{\left(q;q\right)_{2n+1}}
=\frac{\left(q^2;q^4\right)_{\infty}\left(q^4,q^{16},q^{20};q^{20}
\right)_{\infty}}{\left(q;q\right)_{\infty}}
=
\frac{\left(q^4,q^6,q^{10};q^{10}\right)_{\infty}
\left(q^2,q^{18};q^{20}\right)_{\infty}}{\left(q;q\right)_{\infty}}.
\end{align*}

Taking $\left(a,b\right)\to\left(0,0\right)$ into Theorem \ref{T9}, we have
\begin{align*}
&\left(q^2;q^2\right)_{\infty}\sum_{n=0}^{\infty}\frac{q^{2n^2+2n}}
{\left(-q;q^2\right)_n\left(q;q\right)_{2n+1}}\nonumber\\
&\quad=
-\sum_{n=0}^{\infty}\left(-1\right)^n\left(1+q^{4n+3}\right)q^{14n^2+18n+5}
+\sum_{n=0}^{\infty}\left(-1\right)^n\left(1+q^{4n+1}\right)q^{14n^2+6n}
\nonumber\\
&\quad=
\left(q^8,q^{20},q^{28};q^{28}\right)_{\infty}
+q\left(q^4,q^{24},q^{28};q^{28}\right)_{\infty}
,
\end{align*}
which implies \eqref{new7} is true.
 \end{proof}



\section{Identities from the $q$-Pfaff-Saalsch\"utz summation}\label{sec3}
Recall the $q$-Pfaff-Saalsch\"utz summation formula \cite[Eq. (1.7.2)]{Gasper_and_Rahman}
\begin{equation}\label{Pfaff}
_3\phi_2\left(
\begin{array}{ccc}
q^{-n},\,\alpha q^n,\,\alpha bc/q\\
\alpha b,\,\alpha c
\end{array};q,\,q\right)=\frac{\left(q/b,q/c;q\right)_{n}}
{\left(\alpha b,\alpha c;q\right)_n}\left(\alpha bc/{q}\right)^n.
\end{equation}
%\subsection{Proofs of Slater's (11), (27) and (87)}
In \cite{Liu2}, Liu applied the $q$-Pfaff-Saalsch\"utz summation formula to sum the $_3\phi_2$ series (with $\beta=\alpha cd/q$)  on the right-hand side of \eqref{Formula1} to obtain
\begin{equation*}
  \frac{\left(q/c,q/d;q\right)_n}{\left(\alpha c,\alpha d;q\right)_n}\left(\alpha cd/q\right)^n,
\end{equation*}
and made \eqref{Formula1} reduce to the following beautiful result.
\begin{theorem}\label{thm2}
For $|\alpha ab/q|<1$, we have
\begin{align}
&\frac{\left(\alpha q,\alpha ab/q;q\right)_{\infty}}{\left(\alpha a,\alpha b;q\right)_{\infty}}{_3\phi_2}\left(
\begin{array}{ccc}
q/a,q/b,\alpha cd/q\\
\alpha c,\alpha d
\end{array};q,{\alpha ab}/{q}\right)\nonumber\\
&\quad=
\sum_{n=0}^{\infty}\frac{\left(1-\alpha q^{2n}\right)
\left(\alpha,q/a,q/b,q/c,q/d;q\right)_n}{\left(1-\alpha\right)\left(q,\alpha a,\alpha b,\alpha c,\alpha d;q\right)_n}
\left(-\alpha^2abcd\right)^n
q^{n\left(n-5\right)/2}
.\label{B4}
\end{align}
\end{theorem}
In this section, by specializing the choices of $\left(q,\alpha,a,b,c,d\right)$ in Theorem \ref{thm2}, we  
find some Rogers-Ramanujan type identities.% and an Appell-Lerch series identity. 
\begin{corollary}
We have
\begin{align}
\sum_{n=0}^{\infty}\frac{\left(-q;q\right)_{2n}q^{n\left(n+1\right)}}
{\left(q;q^2\right)_{n+1}\left(q^4;q^4\right)_{n}}
&=\frac{\left(-q^2;q^2\right)_{\infty}\left(-q,-q^3,q^4;q^4\right)_{\infty}}
{\left(q^2;q^2\right)_{\infty}},\tag{S.11}\label{S.11}\\
\sum_{n=0}^{\infty}\frac{\left(q^2;q^4\right)_nq^{2n\left(n+1\right)}}
{\left(q^4;q^4\right)_{n}\left(q;q^2\right)_{n}\left(q;q^2\right)_{n+1}}
&=\frac{\left(-q,-q^5,q^6;q^6\right)_{\infty}}
{\left(q^2;q^2\right)_{\infty}},\tag{S.27\, and\, S.87}\label{S.87}\\
\sum_{n=0}^{\infty}\frac{\left(-q,-q;q^2\right)_n}
{\left(q;q\right)_{2n+1}\left(-q^2;q^2\right)_{n}}q^{n\left(n+2\right)}
&=\left(-q;q\right)_{\infty}\left(-q^4,-q^4;q^4\right)_{\infty}
.\label{new3}
%\sum_{n=0}^{\infty}\frac{\left(-1\right)^n\left(q^2;q^4\right)_n}
%{\left(q;q\right)_{2n+1}\left(-q^2;q^2\right)_{n}}q^{n\left(n+2\right)}
%=\frac{\left(q;q^2\right)_{\infty}}
%{\left(q^2;q^2\right)_{\infty}}
%&=\frac{1}{\psi(q)}\sum_{n=-\infty}^{\infty}
%\left(-1\right)^n\frac{q^{2n^2+2n}}{1-q^{2n+1}}.\label{new4}
\end{align}
\end{corollary}
\begin{remark}
\eqref{new3} is a special case of Andrews' $q$-analogue of Gauss's second theorem \cite[Eq. (1.9)]{Andrews_1973} 
\begin{align*}
\sum_{k=0}^{\infty}\frac{\left(b,q/b;q\right)_k}{\left(q^2;q^2\right)_k\left(c;q\right)_k}c^kq^{\frac{k^2-k}{2}}=
\frac{\left(bc,qc/b;q^2\right)_{\infty}}{\left(c;q\right)_{\infty}}
\end{align*}
with $q\to q^2, b=-q$ and $c=q^3$. 
\end{remark}
\begin{proof}
By substituting $\left(q,\alpha,c,d\right)\to\left(q^2,q^2,-1,q\right)$ into \eqref{B4}, we obtain
\begin{align}
&\frac{\left(q^2,ab;q^2\right)_{\infty}}
{\left(aq^2,bq^2;q^2\right)_{\infty}}\sum_{n=0}^{\infty}
\frac{\left(-q,q^2/a,q^2/b;q^2\right)_{n}}
{\left(q;q\right)_{2n+1}\left(-q^2;q^2\right)_{n}}a^nb^n\nonumber\\
&\quad=\sum_{n=0}^{\infty}\left(1+q^{2n+1}\right)
\frac{\left(q^2/a,q^2/b;q^2\right)_{n}}
{\left(aq^2,bq^2;q^2\right)_n}
a^nb^nq^{n^2},\label{thm5}
\end{align}
where $\max\{|ab|, |aq^2|,|bq^2|\}<1$.

Note that
\begin{equation*}
\sum_{n=0}^{\infty}\frac{\left(-q;q\right)_{2n}}
{\left(q;q^2\right)_{n+1}\left(q^4;q^4\right)_{n}}q^{n\left(n+1\right)}
=\sum_{n=0}^{\infty}\frac{\left(-q;q^2\right)_{n}}
{\left(q;q\right)_{2n+1}}q^{n\left(n+1\right)}.
\end{equation*}
Hence, we take $\left(a,b\right)\to\left(0,-1\right)$ into \eqref{thm5} and use the Jacobi triple product identity \eqref{Jacobi} with $\left(q,z\right)\to\left(q^4,q\right)$ to obtain \eqref{S.11}.

Also, we observe that
\begin{equation*}
\sum_{n=0}^{\infty}\frac{\left(q^2;q^4\right)_nq^{2n\left(n+1\right)}}
{\left(q^4;q^4\right)_{n}\left(q;q^2\right)_{n}\left(q;q^2\right)_{n+1}}
=\sum_{n=0}^{\infty}\frac{\left(-q;q^2\right)_nq^{2n\left(n+1\right)}}
{\left(q;q\right)_{2n+1}\left(-q^2;q^2\right)_{n}}.
\end{equation*}
Then, we take $\left(a,b\right)\to\left(0,0\right)$ into \eqref{thm5} and use the Jacobi triple product identity \eqref{Jacobi} with $\left(q,z\right)\to\left(q^6,q\right)$ to obtain \eqref{S.87}.

Moreover, by inserting $\left(a,b\right)\to\left(0,-q\right)$ into \eqref{thm5}, we arrive at \eqref{new3}. % $\left(a,b\right)\to\left(0,q\right)$ \eqref{new4}.
\end{proof}

\begin{corollary}
We have
\begin{align}
\sum_{n=0}^{\infty}\frac{\left(-q;q\right)_nq^{n^2+n}}
{\left(q;q\right)_n\left(q;q^2\right)_{n+1}}
&=
\frac{\left(-q;q\right)_{\infty}\left(q,q^5,q^6;q^6\right)_{\infty}}
{\left(q;q\right)_{\infty}},\tag{S.22}\label{S.22}\\
\sum_{n=0}^{\infty}\frac{\left(-q^2;q^2\right)_n}
{\left(q;q\right)_{2n+1}}q^{n^2}
&=\frac{\left(-q;q^2\right)_{\infty}\left(-q,-q^3,q^4;q^4\right)_{\infty}}{\left(q;q^2\right)_{\infty}
}.\tag{\rm{\cite[Entry 1.7.13]{Ramanujan}}}\label{new6}
\end{align}
%which can be found in {\rm{\cite[Entry 1.7.13]{Ramanujan}}}.
\end{corollary}
\begin{proof}
We take $\left(q,\alpha,c,d\right)\to\left(q^2,q^2,q,-q\right)$ into Theorem \ref{thm2} to arrive at
\begin{align}
\frac{\left(q^2,ab;q^2\right)_{\infty}}{\left(aq^2,bq^2;q^2\right)_{\infty}}
\sum_{n=0}^{\infty}\frac{\left(q^2/a,q^2/b,-q^2;q^2\right)_n}
{\left(q^2;q^2\right)_n\left(q^2;q^4\right)_{n+1}}a^nb^n
=
\sum_{n=0}^{\infty}\frac{\left(q^2/a,q^2/b;q^2\right)_n}
{\left(aq^2,bq^2;q^2\right)_n}a^nb^nq^{n^2+n},\label{thm3}
\end{align}
where $\max\left\{|ab|,|aq^2|,|bq^2|\right\}<1$.

Taking $\left(a,b\right)\to\left(0,0\right)$ into \eqref{thm3} and using the Jacobi triple product identity \eqref{Jacobi} with $\left(q,z\right)\to\left(q^6,1\right)$, we see that
\begin{align}
\left(q^2;q^2\right)_{\infty}
\sum_{n=0}^{\infty}\frac{\left(-q^2;q^2\right)_nq^{2n^2+2n}}
{\left(q^2;q^2\right)_n\left(q^2;q^4\right)_{n+1}}
=\left(-q^6,-q^6,q^6;q^6\right)_{\infty}.\label{e2}
\end{align}
After replacing $q^2$ by $q$, \eqref{e2} reduces to \eqref{S.22}.

Taking $\left(a,b\right)\to\left(0,-1/q\right)$ into \eqref{thm3} and using the Jacobi triple product identity \eqref{Jacobi} with $\left(q,z\right)\to\left(q^4,q\right)$, we arrive at the identity \ref{new6}.
\end{proof}
\begin{corollary}
We have
\begin{align}
\sum_{n=0}^{\infty}\frac{\left(q^2;q^2\right)_{n-1}q^{n^2}}
{\left(q;q\right)_n\left(q;q\right)_{n-1}\left(q;q^2\right)_n}
&=\frac{\left(-q^5,-q^7,q^{12};q^{12}\right)_{\infty}}{\left(q;q\right)_{\infty}}
,\tag{S.58}\label{S.58}\\
\sum_{n=0}^{\infty}\frac{\left(-q;q^2\right)_{n}\left(q^4;q^4\right)_{n-1}q^{n^2}}
{\left(q^2;q^2\right)_n\left(q^2;q^2\right)_{n-1}\left(q^2;q^4\right)_n}
&=\frac{\left(-q^6,-q^{10},q^{16};q^{16}\right)_{\infty}
\left(-q;q^2\right)_{\infty}}
{\left(q^2;q^2\right)_{\infty}}
.\tag{S.72}\label{S.72}
\end{align}
\end{corollary}
\begin{proof}
Taking $\left(\alpha,c,d\right)\to\left(1,q^{1/2},-q^{1/2}\right)$ into Theorem \ref{thm2} , 
\begin{comment}
 we have
\begin{equation*}
_3\phi_2\left(
\begin{array}{ccc}
q^{-n},q^{n},-1\\
q^{{1}/{2}},-q^{{1}/{2}}
\end{array};q,\,q\right)
=\left\{\begin{array}{ll}
1,&\quad n\equiv0\pmod2,\\
-1,&\quad n\equiv1\pmod2.
\end{array}
\right.
\end{equation*}
\end{comment}
 we have
\begin{align}
&\frac{\left(q,ab/q;q\right)_{\infty}}{\left(a,b;q\right)_{\infty}}
\sum_{n=0}^{\infty}\frac{\left(q/a,q/b,-1;q\right)_{n}}
{\left(q;q\right)_{n}\left(q;q^2\right)_{n}}
\left(ab/q\right)^n\nonumber\\
&\quad=1+\sum_{n=1}^{\infty}\left(1+q^{n}\right)
\frac{\left(q/a,q/b;q\right)_{n}}{\left(a,b;q\right)_{n}}
a^{n}b^{n}q^{\frac{n^2-3n}{2}},\label{C16}
%&\quad=1+\sum_{n=1}^{\infty}\left(1+q^{2n}\right)
%\frac{\left(q/a,q/b;q\right)_{2n}}{\left(a,b;q\right)_{2n}}
%a^{2n}b^{2n}q^{2n^2-3n}\nonumber\\
%&\qquad+q^{-1}\sum_{n=0}^{\infty}\left(1+q^{2n+1}\right)
%\frac{\left(q/a,q/b;q\right)_{2n+1}}{\left(a,b;q\right)_{2n+1}}
%a^{2n+1}b^{2n+1}q^{2n^2-n},
\end{align}
where $\max\{|a|,|b|,|ab/q|\}<1$.

Taking $\left(a,b\right)\to\left(0,0\right)$ into \eqref{C16}, we have
\begin{align}
&\left(q;q\right)_{\infty}\sum_{n=0}^{\infty}
\frac{\left(-1;q\right)_nq^{n^2}}
{\left(q;q\right)_n\left(q;q^2\right)_n}\nonumber\\
&\quad=
\sum_{n=-\infty}^{\infty}q^{6n^2+n}+q\sum_{n=-\infty}^{\infty}q^{6n^2+5n}\nonumber\\
&\quad=
\left(-q^5,-q^7,q^{12};q^{12}\right)_{\infty}
+q\left(-q,-q^{11},q^{12};q^{12}\right)_{\infty}.
\label{equ5}
\end{align}
Consequently, we find that
\begin{align*}
&\sum_{n=0}^{\infty}\frac{\left(q^2;q^2\right)_{n-1}q^{n^2}}
{\left(q;q\right)_n\left(q;q\right)_{n-1}\left(q;q^2\right)_n}\nonumber\\
&\quad=\frac{1}{2}+\frac{1}{2}\sum_{n=0}^{\infty}
\frac{\left(-1;q\right)_nq^{n^2}}{\left(q;q\right)_n
\left(q;q^2\right)_n}\nonumber\\
&\quad=\frac{1}{2}+\frac{\left(-q^5,-q^7,q^{12};q^{12}\right)_{\infty}
+q\left(-q,-q^{11},q^{12};q^{12}\right)_{\infty}}{2\left(q;q\right)_{\infty}}\\
&\quad=\left(-q^5,-q^7,q^{12};q^{12}\right)_{\infty},
\end{align*}
where we have used
\begin{align*}
\left(q;q\right)_{\infty}
%&=\left(q,q^2,q^3;q^3\right)_{\infty}\nonumber\\
&=\sum_{n=-\infty}^{\infty}\left(-1\right)^nq^{\frac{3n^2-n}{2}}\nonumber\\
&=\sum_{n=-\infty}^{\infty}q^{6n^2-n}-q\sum_{n=-\infty}^{\infty}q^{6n^2+5n}\nonumber\\
&=\left(-q^5,-q^7,q^{12};q^{12}\right)_{\infty}
-q\left(-q,-q^{11},q^{12};q^{12}\right)_{\infty}.
\end{align*}
Therefore, the proof of \eqref{S.58} is completed.

Taking $\left(a,b\right)\to\left(0,-q^{1/2}\right)$ into \eqref{C16}, we have
\begin{align}
&\frac{\left(q;q\right)_{\infty}}{\left(-q^{1/2};q\right)_{\infty}}
\sum_{n=0}^{\infty}
\frac{\left(-1;q\right)_n}
{\left(q;q\right)_n\left(q^{1/2};q\right)_n}q^{{n^2}/{4}}\nonumber\\
&\quad=
\sum_{n=-\infty}^{\infty}q^{4n^2-n}
+q^{{1}/{2}}\sum_{n=-\infty}^{\infty}q^{4n^2+3n}\nonumber\\
&\quad
=\left(-q^3,-q^{5},q^{8};q^{8}\right)_{\infty}
+q^{{1}/{2}}\left(-q,-q^{7},q^{8};q^{8}\right)_{\infty}.
\label{equ6}
\end{align}
Therefore, we can prove \eqref{S.72} as follows,
\begin{align*}
&\sum_{n=0}^{\infty}\frac{\left(-q;q^2\right)_{n}\left(q^4;q^4\right)_{n-1}}
{\left(q^2;q^2\right)_n\left(q^2;q^2\right)_{n-1}\left(q^2;q^4\right)_n}q^{n^2}\\
&\quad=\frac{1}{2}+\frac{1}{2}
\sum_{n=0}^{\infty}\frac{\left(-1;q^2\right)_n}{\left(q,q^2;q^2\right)_n}q^{n^2}\\
&\quad=\frac{1}{2}+\frac{
\left(-q;q^2\right)_{\infty}
}
{2\left(q^2;q^2\right)_{\infty}}\left(\left(-q^6,-q^{10},q^{16};q^{16}\right)_{\infty}
+q\left(-q^2,-q^{14},q^{16};q^{16}\right)_{\infty}\right)\\
&\quad=\frac{\left(-q;q^2\right)_{\infty}\left(-q^6,-q^{10},q^{16};q^{16}\right)_{\infty}
}
{\left(q^2;q^2\right)_{\infty}},
\end{align*}
where we have used \eqref{equ6} by replacing $q$ by $q^2$ and the following result:
\begin{align*}
\frac{\left(q^2;q^2\right)_{\infty}}{\left(-q;q^2\right)_{\infty}}
&=\sum_{n=-\infty}^{\infty}\left(-1\right)^nq^{2n^2-n}\\
&=\sum_{n=-\infty}^{\infty}q^{8n^2-2n}-
\sum_{n=-\infty}^{\infty}q^{8n^2+6n}\\
&=\left(-q^6,-q^{10},q^{16};q^{16}\right)_{\infty}
-q\left(-q^2,-q^{14},q^{16};q^{16}\right)_{\infty}.
\end{align*}
\end{proof}
%\subsection{An identity on Ramanujan's theta function $\varphi(q)$ }

\begin{comment}
\begin{corollary}
We have
\begin{align}
\sum_{n=0}^{\infty}\frac{\left(-1\right)^n
q^{\frac{n\left(n+3\right)}{2}}}{\left(q;q\right)_n
\left(1+q^{n+1}\right)^2}
=
\frac{\varphi(-q)}
{q}\sum_{n=1}^{\infty}\frac{q^n}{1+q^n}.\label{e1}
\end{align}
%which is given by \cite{Hari} without proving. Here $\varphi(q)$ is another Ramanujan's theta function and satisfies
%\begin{align*}
%\varphi(q)=\sum_{n=-\infty}^{\infty}q^{n^2}
%=\frac{\left(q;q\right)_{\infty}}{\left(-q;q\right)_{\infty}}.
%\end{align*}
\end{corollary}
\begin{proof}
Substituting $\left(\alpha,c,d\right)\to\left(-q^2,-q^2,-q^2\right)$ into Theorem \ref{thm2}, we obtain 
\begin{align}
&\frac{\left(-q^2,-abq;q\right)_{\infty}}
{\left(-aq^2,-bq^2;q\right)_{\infty}}\sum_{n=0}^{\infty}
\frac{\left(q/a,q/b,-q;q\right)_{n}}
{\left(q,-q^2,-q^2;q\right)_n}\left(-abq\right)^n\nonumber\\
&\quad=\sum_{n=0}^{\infty}\left(-1\right)^n\left(1+q^{2n+2}\right)
\frac{\left(q,q/a,q/b;q\right)_n}{\left(-q^2,-aq^2,-bq^2;q\right)_n}
a^nb^nq^{\frac{n^2+3n}{2}},\label{thm8}
\end{align}
where $\max\{|aq^2|, |bq^2|, |abq|\}<1$.

Taking $\left(a,b\right)\to\left(0,-1\right)$ into \eqref{thm8}, we have
\begin{align}
&\sum_{n=0}^{\infty}
\left(-1\right)^n\frac{\left(-q,-q;q\right)_n
}{\left(q,-q^2,-q^2;q\right)_n}q^{\frac{n\left(n+3\right)}{2}}\nonumber\\
&\quad=\frac{\left(1+q\right)\left(q;q\right)_{\infty}}
{\left(-q^2;q\right)_{\infty}}\sum_{n=0}^{\infty}\left(-1\right)^n
\frac{\left(1+q^{2n+2}\right)}{\left(1-q^{2n+2}\right)}q^{n^2+2n}\nonumber\\
&\quad=\frac{\left(1+q\right)\left(q;q\right)_{\infty}}
{q\left(-q^2;q\right)_{\infty}}\sum_{n=1}^{\infty}
\frac{q^{n}}{1+q^n},
\end{align}
where the last equation follows from Liu's result \cite{Liu2}
\begin{equation*}
\sum_{n=1}^{\infty}\frac{q^n}{1+q^n}
=-\sum_{n=1}^{\infty}\left(-1\right)^n\frac{1+q^{2n}}{1-q^{2n}}q^{n^2}.
\end{equation*}
Then, we obtain
\begin{equation*}
\sum_{n=0}^{\infty}\frac{\left(-1\right)^n
q^{\frac{n\left(n+3\right)}{2}}}{\left(q;q\right)_n
\left(1+q^{n+1}\right)^2}
=\frac{1}{\left(1+q\right)^2}\sum_{n=0}^{\infty}
\left(-1\right)^n\frac{\left(-q,-q;q\right)_n
}{\left(q,-q^2,-q^2;q\right)_n}q^{\frac{n\left(n+3\right)}{2}}
=\frac{\varphi(-q)}
{q}\sum_{n=1}^{\infty}\frac{q^n}{1+q^n}
.
\end{equation*}
\end{proof}
\end{comment}

\section{Identities from a terminating $q$-analogue of Whipple's ${_3F_2}$ sum}
%\subsection{Proofs of Slater's (2), (17), (20), (31), (32), (94) and (99)}
In this section,  we derive some identities from  an useful summation formula which can be found in \cite{Gasper_and_Rahman},
\begin{equation}\label{Whipple}
_4\phi_3\left(
\begin{array}{cccc}
q^{-n},\,q^{n+1},\,c,\,-c\\
e,\,c^2q/e,\,-q
\end{array};q,\,q\right)
=\frac{\left(eq^{-n},eq^{1+n},c^2q^{1-n}/e,c^2q^{n+2}/e;q^2\right)_{\infty}}
{\left(e,c^2q/e;q\right)_{\infty}}q^{n\left(n+1\right)/2}.
\end{equation}
%From \eqref{Whipple}, we can obtain the following summations,


\begin{theorem}\label{thm11}
For $\max\{|ab|,|aq^2|,|bq^2|\}<1$, we have
\begin{align*}
&\frac{\left(q^2,ab;q^2\right)_{\infty}}{\left(aq^2,bq^2;q^2\right)_{\infty}}
\sum_{n=0}^{\infty}
\frac{\left(q^2/a,q^2/b;q^2\right)_n}{\left(-q^2,q^2,-q;q^2\right)_n}a^nb^n\\
&\quad=\sum_{n=0}^{\infty}\left(-1\right)^n\left(1-q^{4n+2}\right)
\frac{\left(q^2/a,q^2/b;q^2\right)_n}{\left(aq^2,bq^2;q^2\right)_n}
a^nb^nq^{\frac{3n^2-n}{2}}.
\end{align*}
\end{theorem}
\begin{proof}
Making the substitution $\left(q,c,e\right)\to\left(q^2,0,-q\right)$ into \eqref{Whipple}, we have
\begin{equation}\label{A13}
{_3\phi_2}\left(
\begin{array}{ccc}
q^{-2n},q^{2n+2},0\\
-q^2,-q
\end{array};q^2,q^2\right)
=
\frac{\left(-q^{1-2n},-q^{2n+3};q^4\right)_{\infty}}{\left(-q;q^2\right)_{\infty}}
q^{n\left(n+1\right)}=q^{\frac{n^2+n}{2}}.
\end{equation}
Then, we substitute $\left(q,\alpha,\beta,c,d\right)\to\left(q^2,q^2,0,-q^2,-q\right)$ into \eqref{Formula1} to obtain the desired theorem.
\end{proof}
\begin{corollary}
We have
\begin{align}
\sum_{n=0}^{\infty}\frac{q^{n\left(n+2\right)}}{\left(q^4;q^4\right)_n}
&=\frac{\left(-q;q^2\right)_{\infty}\left(q,q^4,q^5;q^5\right)_{\infty}}
{\left(q^2;q^2\right)_{\infty}}
,\tag{S.16}\label{S.16}\\
\sum_{n=0}^{\infty}\frac{q^{2n\left(n+1\right)}}
{\left(q^2;q^2\right)_n\left(-q;q\right)_{2n}}
&=\frac{\left(q^2,q^5,q^7;q^7\right)_{\infty}}
{\left(q^2;q^2\right)_{\infty}},\tag{S.32}\label{S.32}\\
\sum_{n=0}^{\infty}\frac{\left(-1\right)^n\left(q;q^2\right)_nq^{n^2+2n}}
{\left(-q;q^2\right)_{n}\left(q^4;q^4\right)_n}
&=\frac{\left(q;q^2\right)_{\infty}\left(-q,-q^4,q^5;q^5\right)_{\infty}}
{\left(q^2;q^2\right)_{\infty}}.\tag{\rm{\cite[Eq. (2.23)]{Bowman}}}\label{e6}
\end{align}
\end{corollary}
\begin{proof}
By substituting $\left(a,b\right)\to\left(0,-q\right)$, $\left(a,b\right)\to\left(0,0\right)$ and $\left(a,b\right)\to\left(0,q\right)$ into Theorem \ref{thm11}, respectively, we obtain \eqref{S.16}, \eqref{S.32} and the identity \ref{e6}.
\end{proof}


\begin{theorem}\label{thm10}
For $\max\{|ab|,|aq^2|,|bq^2|\}<1$, we have
\begin{align*}
&\frac{\left(q^2,ab;q^2\right)_{\infty}}{\left(aq^2,bq^2;q^2\right)_{\infty}}
\sum_{n=0}^{\infty}
\frac{\left(q^2/a,q^2/b;q^2\right)_n}{\left(-q^2,q^2;q^2\right)_n\left(-q;q^2\right)_{n+1}}a^nb^n\\
&\quad=\sum_{n=0}^{\infty}\left(-1\right)^n\left(1-q^{2n+1}\right)
\frac{\left(q^2/a,q^2/b;q^2\right)_n}{\left(aq^2,bq^2;q^2\right)_n}
a^nb^nq^{\frac{3n^2+n}{2}}.
\end{align*}
\end{theorem}
\begin{proof}
Making the substitution $\left(q,c,e\right)\to\left(q^2,0,-q^3\right)$ into \eqref{Whipple}, we have
\begin{align}
{_3\phi_2}\left(
\begin{array}{ccc}
q^{-2n},q^{2n+2},0\\
-q^2,-q^3
\end{array};q^2,q^2\right)
%=
%\frac{\left(-q^{3-2n};q^4\right)_{\infty}\left(-q^{5+2n};q^4\right)_{\infty}}
%{\left(-q^{3};q^2\right)_{\infty}}q^{n\left(n+1\right)}
%=\frac{\left(-q^{3-2n};q^4\right)_{n}}
%{\left(-q^{3};q^2\right)_{n}}
=
\frac{1+q}{1+q^{2n+1}}q^{\frac{n^2+3n}{2}},\label{A18}
\end{align}
which can also be found in \cite[Eq. (3.13)]{Wang}.
Based on the above identity,  we substitute $\left(q,\alpha,\beta,c,d\right)\to\left(q^2,q^2,0,-q^2,-q^3\right)$ into \eqref{Formula1} to obtain the desired theorem.
\end{proof}
\begin{corollary}
We have
\begin{align}
\sum_{n=0}^{\infty}\frac{q^{n\left(n+1\right)}}
{\left(q^2;q^2\right)_n\left(-q;q^2\right)_{n+1}}
&=
\frac{\left(-q^2;q^2\right)_{\infty}\left(q,q^4,q^5;q^5\right)_{\infty}}
{\left(q^2;q^2\right)_{\infty}},\tag{S.17}\label{S.17}\\
\sum_{n=0}^{\infty}\frac{\left(-1\right)^n\left(q;q^2\right)_nq^{n^2+2n}}
{\left(-q;q^2\right)_{n+1}\left(q^4;q^4\right)_n}
&=\frac{\left(q;q^2\right)_{\infty}\left(-q^5,-q^5,q^5;q^5\right)_{\infty}}
{\left(q^2;q^2\right)_{\infty}},\tag{\rm{\cite[Eq. (2.5)]{McLaughlin}}}\label{new1}\\
\sum_{n=0}^{\infty}\frac{q^{n^2}}{\left(q^4;q^4\right)_n}
&=\frac{\left(-q;q^2\right)_{\infty}\left(q^2,q^3,q^5;q^5\right)_{\infty}}
{\left(q^2;q^2\right)_{\infty}},\tag{S.20}\label{S.20}\\
\sum_{n=0}^{\infty}\frac{q^{2n\left(n+1\right)}}
{\left(q^2;q^2\right)_n\left(-q;q\right)_{2n+1}}
&=\frac{\left(q,q^6,q^7;q^7\right)_{\infty}}
{\left(q^2;q^2\right)_{\infty}}.\tag{S.31}\label{S.31}
%&\sum_{n=0}^{\infty}\frac{q^{2n\left(n+1\right)}}
%{\left(q;q\right)_{2n}\left(-q^2;q^2\right)_n}
%=\frac{\left(q^{11},q^{17},q^{28};q^{28}\right)_{\infty}
%+q^5\left(q^{-3},q^{31},q^{28};q^{28}\right)_{\infty}}
%{\left(q^2;q^2\right)_{\infty}},\quad \mathrm{\left(S. 32-\right)}\label{e4}\\
\end{align}
\end{corollary}
\begin{proof}
Substituting $\left(a,b\right)\to\left(0,-1\right)$ into Theorem \ref{thm10} and taking $\left(q,z\right)\to\left(q^5, -q\right)$ into the Jacobi triple product identity \eqref{Jacobi},, we deduce that \eqref{S.17} is true.
%\begin{align*}
%&\frac{\left(q^2;q^2\right)_{\infty}}{\left(-q^2;q^2\right)_{\infty}}
%\sum_{n=0}^{\infty}\frac{q^{n\left(n+1\right)}}
%{\left(q^2;q^2\right)_n\left(-q;q^2\right)_{n+1}}
%\\
%&\quad=\frac{\left(q^2;q^2\right)_{\infty}}{\left(-q^2;q^2\right)_{\infty}}
%\sum_{n=0}^{\infty}
%\frac{\left(-q^2;q^2\right)_n}{\left(q^2,-q^2,-q^3;q^2\right)_n}
%q^{n\left(n+1\right)}\\
%&\quad=\sum_{n=0}^{\infty}\left(-1\right)^n\left(1-q^{4n+2}\right)
%\frac{\left(-q^{3-2n};q^4\right)}{\left(-q^3;q^2\right)_n}q^{3n^2+n}\\
%&\quad=\left(1+q\right)\sum_{n=0}^{\infty}\left(-1\right)^n
%\left(1-q^{2n+1}\right)q^{\frac{5}{2}n^2+\frac{3}{2}n}\\
%&\quad=\left(1+q\right)\left(q,q^4,q^5;q^5\right)_{\infty},
%\end{align*}
%which implies .

Moreover, setting $\left(a,b\right)\to\left(0,q\right)$ in Theorem \ref{thm10} and taking $\left(q,z\right)\to\left(q^5,1\right)$ into the Jacobi triple product identity \eqref{Jacobi}, we deduce that \ref{new1} is true.

Likewise, letting $\left(a,b\right)\to\left(0,-1/q\right)$ in Theorem \ref{thm10} and inserting $\left(q,z\right)\to\left(q^5,-q^2\right)$ into the Jacobi triple product identity \eqref{Jacobi}, we obtain \eqref{S.20}.

Similarly, substituting $\left(a,b\right)\to\left(0,0\right)$ into Theorem \ref{thm10} and using the Jacobi triple product identity \eqref{Jacobi} with $\left(q,z\right)\to\left(q^7,-q\right)$, we get \eqref{S.31}.
\end{proof}
%
%
\begin{theorem}\label{thm12}
For $\max\{|ab|,|aq^2|,|bq^2|\}<1$, we have
\begin{align*}
&\frac{\left(q^2,ab;q^2\right)_{\infty}}{\left(aq^2,bq^2;q^2\right)_{\infty}}
\sum_{n=0}^{\infty}
\frac{\left(q^2/a,q^2/b;q^2\right)_n}{\left(-q^2,q^2;q^2\right)_n\left(q;q^2\right)_{n+1}}a^nb^n\\
&\quad=\sum_{n=0}^{\infty}\left(-1\right)^n\left(1+q^{4n+1}\right)
\frac{\left(q^2/a,q^2/b;q^2\right)_{2n}}{\left(aq^2,bq^2;q^2\right)_{2n}}
a^{2n}b^{2n}q^{6n^2+n}\\
&\qquad-
\sum_{n=0}^{\infty}\left(-1\right)^n\left(1+q^{4n+3}\right)
\frac{\left(q^2/a,q^2/b;q^2\right)_{2n+1}}{\left(aq^2,bq^2;q^2\right)_{2n+1}}
a^{2n+1}b^{2n+1}q^{6n^2+7n+2}.
\end{align*}
\end{theorem}
\begin{proof}
Making the substitution $\left(q,c,e\right)\to\left(q^2,0,q^3\right)$ into \eqref{Whipple}, we have
\begin{align}
{_3\phi_2}\left(
\begin{array}{ccc}
q^{-2n},q^{2n+2},0\\
-q^2,q^3
\end{array};q^2,q^2\right)
&=
\frac{\left(q^{3-2n}, q^{5+2n};q^4\right)_{\infty}}
{\left(q^{3};q^2\right)_{\infty}}q^{n\left(n+1\right)}\nonumber\\
&
=\left\{\begin{array}{ll}
\left(-1\right)^mq^{2m^2+3m}\frac{1-q}{1-q^{4m+1}},&\quad {n=2m},\\
\left(-1\right)^{m}q^{2m^2+5m+2}\frac{1-q}{1-q^{4m+3}},&\quad{n=2m+1}.
\end{array}\right.\label{A19}
\end{align}
Based on the above identity,  we substitute $\left(q,\alpha,\beta,c,d\right)\to\left(q^2,q^2,0,-q^2,q^3\right)$ into \eqref{Formula1} to obtain the desired theorem.
\end{proof}
\begin{corollary}
We have
\begin{align}
&\sum_{n=0}^{\infty}\frac{\left(-q;q^2\right)_{n}}
{\left(q;q^2\right)_{n+1}\left(q^4;q^4\right)_n}q^{n\left(n+2\right)}
=\frac{\left(-q^2;q^2\right)_{\infty}
}
{\left(q^2;q^2\right)_{\infty}}\left(q^5,q^{15},q^{20};q^{20}\right)_{\infty},\tag{\cite[Eq (2.7)]{McLaughlin}}\label{N3}\\
%&\sum_{n=0}^{\infty}\frac{q^{2n\left(n+1\right)}}
%{\left(q;q^2\right)_{n+1}\left(q^4;q^4\right)_n}
%=\frac{\left(q^9,q^{19},q^{28};q^{28}\right)_{\infty}
%+q\left(q^5,q^{23},q^{28};q^{28}\right)_{\infty}}
%{\left(q^2;q^2\right)_{\infty}},\label{N4}\\
&\sum_{n=0}^{\infty}\frac{q^{n\left(n+1\right)}}
{\left(q;q\right)_{2n+1}}
=\frac{\left(-q^2;q^2\right)_{\infty}
}
{\left(q^2;q^2\right)_{\infty}}\left(\left(q^7,q^{13},q^{20};q^{20}\right)_{\infty}
+q\left(q^3,q^{17},q^{20};q^{20}\right)_{\infty}\right).\tag{S.94}\label{S.94}
\end{align}
\end{corollary}
\begin{proof}
Taking $\left(a,b\right)\to\left(0,-q\right)$ and  $\left(a,b\right)\to\left(0,-1\right)$ in Theorem \ref{thm12}, respectively,  we arrive at \eqref{N3} and \eqref{S.94}.
%$\left(a,b\right)\to\left(0,0\right)$, we arrive at \eqref{N4}.
\end{proof}
%
%
\begin{theorem}\label{thm13}
For $\max\{|ab|,|aq^2|,|bq^2|\}<1$, we have
\begin{align*}
&\frac{\left(q^2,ab;q^2\right)_{\infty}}{\left(aq^2,bq^2;q^2\right)_{\infty}}
\sum_{n=0}^{\infty}
\frac{\left(q^2/a,q^2/b;q^2\right)_n}{\left(q,q^2,-q^2;q^2\right)_n}a^nb^n\\
&\quad=\sum_{n=0}^{\infty}\left(-1\right)^n\left(1-q^{8n+2}\right)
\frac{\left(q^2/a,q^2/b;q^2\right)_{2n}}{\left(aq^2,bq^2;q^2\right)_{2n}}
a^{2n}b^{2n}q^{6n^2-n}\\
&\qquad+
\sum_{n=0}^{\infty}\left(-1\right)^n\left(1-q^{8n+6}\right)
\frac{\left(q^2/a,q^2/b;q^2\right)_{2n+1}}{\left(aq^2,bq^2;q^2\right)_{2n+1}}
a^{2n+1}b^{2n+1}q^{6n^2+5n+1}.
\end{align*}
\end{theorem}
\begin{proof}
Making the substitution $\left(q,c,e\right)\to\left(q^2,0,q\right)$ into \eqref{Whipple}, we have
\begin{align}
{_3\phi_2}\left(
\begin{array}{ccc}
q^{-2n},q^{2n+2},0\\
-q^2,q
\end{array};q^2,q^2\right)
&=
\frac{\left(q^{1-2n}, q^{3+2n};q^4\right)_{\infty}}
{\left(q;q^2\right)_{\infty}}q^{n\left(n+1\right)}\nonumber\\
&
=\left\{\begin{array}{ll}
\left(-1\right)^mq^{2m^2+m},&\quad {n=2m},\\
\left(-1\right)^{m+1}q^{2m^2+3m+1},&\quad{n=2m+1}.
\end{array}\right.\label{A20}
\end{align}
Based on the above identity,  we substitute $\left(q,\alpha,\beta,c,d\right)\to\left(q^2,q^2,0,-q^2,q\right)$ into \eqref{Formula1} to obtain the desired theorem.
\end{proof}
\begin{corollary}
\begin{align}
%\sum_{n=0}^{\infty}
%\frac{q^{2n\left(n+1\right)}}{\left(q;q^2\right)_n\left(q^4;q^4\right)_n}
%&=\frac{\left(q^{11},q^{17},q^{28};q^{28}\right)_{\infty}
%+q^5\left(q^{-3},q^{31},q^{28};q^{28}\right)_{\infty}}{\left(q^2;q^2\right)_{\infty}},\label{N5}\\
&\sum_{n=0}^{\infty}
\frac{\left(-1\right)^n\left(q;q^2\right)_n}{\left(-q;q^2\right)_n\left(q^4;q^4\right)_n}q^{n\left(n+2\right)}
=\frac{\left(q;q^2\right)_{\infty}\left(-q,-q^4,q^5;q^5\right)_{\infty}}{\left(q^2;q^2\right)_{\infty}},\tag{\rm{\cite[Eq. (2.23)]{Bowman}}}\label{N6}\\
&\sum_{n=0}^{\infty}\frac{q^{n\left(n+1\right)}}
{\left(q;q\right)_{2n}}
=\frac{\left(q^9,q^{11},q^{20};q^{20}\right)_{\infty}
-q^2\left(q,q^{19},q^{20};q^{20}\right)_{\infty}}
{\left(-q;q^2\right)_{\infty}\left(q;q\right)_{\infty}}.\tag{S.99}\label{S.99}
\end{align}
\end{corollary}
\begin{proof}
Inserting %$\left(a,b\right)\to\left(0,0\right)$, 
$\left(a,b\right)\to\left(0,-q\right)$ in Theorem \ref{thm13}, we get
\begin{align*}
&\frac{\left(q^2;q^2\right)_{\infty}}{\left(-q;q^2\right)_{\infty}}
\sum_{n=0}^{\infty}\frac{\left(-q;q^2\right)_n}{\left(q,q^2,-q^2;q^2\right)_n}q^{n^2+2n}\\
&\quad=\sum_{n=0}^{\infty}\left(-1\right)^n\left(1-q^{4n+1}\right)q^{10n^2+3n}
+\sum_{n=0}^{\infty}\left(-1\right)^n\left(1-q^{4n+3}\right)q^{10n^2+13n+4}\\
&\quad=\sum_{n=-\infty}^{\infty}\left(-1\right)^nq^{10n^2+3n}-q\sum_{n=0}^{\infty}q^{10n^2+7n}\\
&\quad=\left(q^7,q^{13},q^{20};q^{20}\right)_{\infty}-q\left(q^3,q^{17},q^{20};q^{20}\right)_{\infty}.
\end{align*}
When replacing $q$ by $-q$,   the above identity is rewritten as
\begin{align*}
&\frac{\left(q^2;q^2\right)_{\infty}}{\left(q;q^2\right)_{\infty}}
\sum_{n=0}^{\infty}\frac{\left(q;q^2\right)_n}{\left(-q,q^2,-q^2;q^2\right)_n}q^{n^2+2n}\\
&\quad=\sum_{n=-\infty}^{\infty}q^{10n^2+3n}+q^4\sum_{n=0}^{\infty}q^{10n^2+13n}\\
&\quad=\sum_{n=-\infty}^{\infty}q^{\frac{5n^2+3n}{2}}\\
&\quad=\left(-q,-q^4,q^5;q^5\right)_{\infty},
\end{align*}
which means \ref{N6} is true.

Similarly, setting $\left(a,b\right)\to\left(0,-1\right)$ in Theorem \ref{thm13}, we obtain \eqref{S.99}.
%\begin{align*}
%&\frac{\left(q^2;q^2\right)_{\infty}}{\left(-q^2;q^2\right)_{\infty}}\sum_{n=0}^{\infty}\frac{q^{n\left(n+1\right)}}{\left(q;q\right)_{2n}}\\
%&\quad=\sum_{n=0}^{\infty}\left(-1\right)^n\left(1-q^{4n+2}\right)q^{3n^2+n}
%\frac{\left(q^{1-2n},q^{3+2n};q^4\right)_{\infty}}{\left(q^2;q^2\right)_{\infty}}\\
%&\quad=\sum_{n=0}^{\infty}\left(-1\right)^n\left(1-q^{8n+2}\right)q^{10n^2+n}
%+\sum_{n=0}^{\infty}\left(-1\right)^n\left(1-q^{8n+6}\right)q^{10n^2+11n+3}\\
%&\quad=\sum_{n=-\infty}^{\infty}\left(-1\right)^nq^{10n^2+n}
%-q^2\sum_{n=-\infty}^{\infty}\left(-1\right)^nq^{10n^2-9n}\\
%&\quad=\left(q^{9},q^{11},q^{20};q^{20}\right)_{\infty}
%-q^2\left(q,q^{19},q^{20};q^{20}\right)_{\infty}
%\end{align*}
%which means \eqref{S.99} is true.
\end{proof}
%\subsection{Proofs of Slater's (65), (67), (70) and (127)}
When $\lambda=e$,
Theorem \ref{thm1} reduces to Liu's
another transformation formula with $z=1$ \cite{Liu_1}: for $|\alpha ab/q|<1$,
\begin{align}
&\frac{\left(\alpha q,\alpha ab/q;q\right)_{\infty}}{\left(\alpha a,\alpha b;q\right)_{\infty}}
{_4\phi_3}\left(\begin{array}{cccc}
q/a,\,q/b,\,\beta,\,\gamma\\
c,\,d,\,h
\end{array}
;q,{\alpha ab}/{q}\right)\nonumber\\
&\quad=
\sum_{n=0}^{\infty}
\frac{\left(1-\alpha q^{2n}\right)\left(\alpha,q/a,q/b;q\right)_n\left(-\alpha ab\right)^nq^{{n\left(n-3\right)}/{2}}}
{\left(1-\alpha\right)\left(q,\alpha a,\alpha b;q\right)_n}
{_4\phi_3}\left(\begin{array}{cccc}
q^{-n},\,\alpha q^n,\,\beta,\,\gamma\\
c,\,d,\,h
\end{array}
;q,q\right).\label{Formula2}
\end{align}
Next, in terms of \eqref{Formula2}, we establish  a generalized form of Theorem \ref{thm11} as follows.
\begin{theorem}\label{T5}
For $\max\{|c|, |ab|,|aq^2|, |bq^2|,  |\beta^2q^2/c|\}<1$, we have
\begin{align*}
&\frac{\left(q^2,ab;q^2\right)_{\infty}}{\left(aq^2,bq^2;q^2\right)_{\infty}}
\sum_{n=0}^{\infty}\frac{\left(q^2/a,q^2/b,\beta,-\beta;q^2\right)_n}
{\left(q^2,-q^2,c,\beta^2q^2/c;q^2\right)_n}a^nb^n\\
&\quad=
\sum_{n=0}^{\infty}
\left(-1\right)^n\left(1-q^{4n+2}\right)
\frac{\left(q^2/a,q^2/b;q^2\right)_n
\left(cq^{-2n},\beta^2q^{2-2n}/c;q^4\right)_n}
{\left(aq^2,bq^2;q^2\right)_n\left(c,\beta^2q^2/c;q^2\right)_n}a^nb^nq^{2n^2}.
\end{align*}
\end{theorem}
\begin{proof}
Setting $\left(q,c,e\right)\to\left(q^2,\beta,c\right)$ in \eqref{Whipple}, we get
\begin{align}
&_4\phi_3\left(\begin{array}{cccc}
q^{-2n},q^{2n+2},\beta,-\beta\\
c,\beta^2q^2/c,-q^2
\end{array};q^2\,q^2\right)\nonumber\\
&\quad=\frac{\left(cq^{-2n},cq^{2n+2},\beta^2q^{2-2n}/c,
\beta^2q^{2n+4}/c;q^4\right)_{\infty}}
{\left(c,\beta^2q^2/c;q^2\right)_{\infty}}q^{n\left(n+1\right)}\nonumber\\
&\quad=\frac{\left(cq^{-2n},\beta^2q^{2-2n}/c;q^4\right)_n}
{\left(c,\beta^2q^2/c;q^2\right)_{n}}q^{n^2+n}.\label{equ2}
\end{align}
Then, taking the substitution $q\to q^2, \alpha=q^2, \beta=-\gamma, cd=\beta^2q^2$ and $h=-q^2$  into \eqref{Formula2}, we obtain
\begin{align*}
&\frac{\left(q^2,ab;q^2\right)_{\infty}}{\left(aq^2,bq^2;q^2\right)_{\infty}}
\sum_{n=0}^{\infty}\frac{\left(q^2/a,q^2/b,\beta,-\beta;q^2\right)_n}
{\left(q^2,-q^2,c,\beta^2q^2/c;q^2\right)_n}a^nb^n\nonumber\\
&\quad=
\sum_{n=0}^{\infty}
\left(-1\right)^n\left(1-q^{4n+2}\right)
\frac{\left(q^2/a,q^2/b;q^2\right)_n}
{\left(aq^2,bq^2;q^2\right)_n}a^nb^nq^{n^2-n}
{_4\phi_3}\left(\begin{array}{cccc}
q^{-2n},q^{2n+2},\beta,-\beta\\
c,\beta^2q^2/c,-q^2
\end{array};q^2,\,q^2\right)\\
&\quad=
\sum_{n=0}^{\infty}
\left(-1\right)^n\left(1-q^{4n+2}\right)
\frac{\left(q^2/a,q^2/b;q^2\right)_n
\left(cq^{-2n},\beta^2q^{2-2n}/c;q^4\right)_n}
{\left(aq^2,bq^2;q^2\right)_n\left(c,\beta^2q^2/c;q^2\right)_n}a^nb^nq^{2n^2}
\end{align*}
as desired.
\end{proof}
\begin{corollary}
We have
\begin{align}
&\sum_{n=0}^{\infty}\frac{\left(-1;q^2\right)_n}{\left(q;q\right)_{2n}}q^{n^2+n}
\nonumber\\
&\quad
=\frac{\left(q^{5},q^{7},q^{12};q^{12}\right)_{\infty}
-q\left(q,q^{11},q^{12};q^{12}\right)_{\infty}}
{\left(q;q\right)_{\infty}},\tag{S.48}\label{S.48}\\
&\sum_{n=0}^{\infty}\frac{\left(-q;q^2\right)_n
\left(-1;q^4\right)_n}{\left(q^2;q^2\right)_{2n}}q^{n\left(n+2\right)}\nonumber\\
&\quad=\frac{\left(-q;q^2\right)_{\infty}}
{\left(q^2;q^2\right)_{\infty}}\left(\left(q^6,q^{10},q^{16};q^{16}\right)_{\infty}
-q\left(q^2,q^{14},q^{16};q^{16}\right)_{\infty}\right)
,\tag{S.67}\label{S.67}\\
&\sum_{n=0}^{\infty}\frac{\left(-q;q^2\right)_n
\left(-q^4;q^4\right)_{n-1}}
{\left(q^2;q^2\right)_{2n}}q^{n\left(n+2\right)}\nonumber\\
&\quad=\frac{\left(-q;q^2\right)_{\infty}
}
{\left(q^2;q^2\right)_{\infty}}\left(\left(-q^{28},-q^{36},q^{64};q^{64}\right)_{\infty}
-q\left(-q^{20},-q^{44},q^{64};q^{64}\right)_{\infty}\right)
.\tag{S.127}\label{S.127}
\end{align}
\end{corollary}
\begin{proof}
In order to prove \eqref{S.48},\eqref{S.67} and \eqref{S.127}, we first
insert $\left(c,\beta\right)=\left(q,\mathrm{i}\right)$ with $\mathrm{i}^2=-1$ in Theorem \ref{T5} to obtain
\begin{align}
&\frac{\left(q^2,ab;q^2\right)_{\infty}}{\left(aq^2,bq^2;q^2\right)_{\infty}}
\sum_{n=0}^{\infty}
\frac{\left(q^2/a,q^2/b;q^2\right)_n\left(-1;q^4\right)_n}{\left(q^2;q^2\right)_{2n}}a^nb^n\notag\\
&\quad=
\sum_{n=0}^{\infty}\left(-1\right)^n\left(1-q^{8n+2}\right)\frac{\left(q^2/a,q^2/b;q^2\right)_{2n}}
{\left(aq^2,bq^2;q^2\right)_{2n}}a^{2n}b^{2n}q^{4n^2-2n}\notag\\
&\qquad+
\sum_{n=0}^{\infty}\left(-1\right)^n\left(1-q^{8n+6}\right)\frac{\left(q^2/a,q^2/b;q^2\right)_{2n+1}}
{\left(aq^2,bq^2;q^2\right)_{2n+1}}a^{2n+1}b^{2n+1}q^{4n^2+2n}.\label{C4}
\end{align}
By taking $\left(a,b\right)\to\left(0,0\right)$  in \eqref{C4}, we deduce that
\begin{align*}
&\left(q^2;q^2\right)_{\infty}\sum_{n=0}^{\infty}
\frac{\left(-1;q^4\right)_n}{\left(q^2;q^2\right)_{2n}}q^{2n^2+2n}\\
&\quad=\sum_{n=-\infty}^{\infty}\left(-1\right)^nq^{12n^2+2n}
-q^2\sum_{n=-\infty}^{\infty}\left(-1\right)^nq^{12n^2+10n}\\
&\quad=\left(q^{10},q^{14},q^{24};q^{24}\right)_{\infty}
-q^2\left(q^{2},q^{22},q^{24};q^{24}\right)_{\infty}.
\end{align*}
Then, by $q^2\to q$, we can arrive at \eqref{S.48} .

Moreover, setting $\left(a,b\right)\to\left(0,-q\right)$ in \eqref{C4} leads to
\begin{align}
&\frac{\left(q^2;q^2\right)_{\infty}}{\left(-q;q^2\right)_{\infty}}
\sum_{n=0}^{\infty}\frac{\left(-q;q^2\right)_{n}
\left(-1;q^4\right)_{n}}
{\left(q^2;q^2\right)_{2n}}q^{n\left(n+2\right)}\nonumber\\
&\quad=\sum_{n=0}^{\infty}\left(-1\right)^n\left(1-q^{4n+1}\right)
q^{8n^2+2n}+\sum_{n=0}^{\infty}\left(-1\right)^n
\left(1-q^{4n+3}\right)q^{8n^2+10n+3}\nonumber\\
&\quad=\left(q^6,q^{10},q^{16};q^{16}\right)_{\infty}
-q\left(q^2,q^{14},q^{16};q^{16}\right)_{\infty},\nonumber
\end{align}
which means  \eqref{S.67} is true.

Similarly, we can prove \eqref{S.127} as follows,
\begin{align*}
&\sum_{n=0}^{\infty}\frac{\left(-q;q^2\right)_n
\left(-q^4;q^4\right)_{n-1}}
{\left(q^2;q^2\right)_{2n}}q^{n\left(n+2\right)}\nonumber\\
&\quad=1+\sum_{n=1}^{\infty}\frac{\left(-q;q^2\right)_n
\left(-q^4;q^4\right)_{n-1}}
{\left(q^2;q^2\right)_{2n}}q^{n\left(n+2\right)}\nonumber\\
&\quad=\frac{1}{2}+\frac{1}{2}\sum_{n=1}^{\infty}\frac{\left(-q;q^2\right)_n
\left(-1;q^4\right)_{n}}
{\left(q^2;q^2\right)_{2n}}q^{n\left(n+2\right)}\nonumber\\
&\quad=\frac{1}{2}+\frac{\left(-q;q^2\right)_{\infty}
}
{2\left(q^2;q^2\right)_{\infty}}\left(\left(q^6,q^{10},q^{16};q^{16}\right)_{\infty}
-q\left(q^2,q^{14},q^{16};q^{16}\right)_{\infty}\right),
\end{align*}
where we have used
\begin{align*}
&\left(q^6,q^{10},q^{16};q^{16}\right)_{\infty}
=\left(-q^{28},-q^{36},q^{64};q^{64}\right)_{\infty}
-q^6\left(-q^{4},-q^{60},q^{64};q^{64}\right)_{\infty},
\\
&\left(q^2,q^{14},q^{16};q^{16}\right)_{\infty}
=\left(-q^{20},-q^{44},q^{64};q^{64}\right)_{\infty}
+q^2\left(-q^{12},-q^{52},q^{64};q^{64}\right)_{\infty}
\end{align*}
and
\begin{align*}
\frac{\left(q^2;q^2\right)_{\infty}}{\left(-q;q^2\right)_{\infty}}
%&=\left(q,q^3,q^4;q^4\right)_{\infty}\\
&=\sum_{n=-\infty}^{\infty}\left(-1\right)^nq^{2n^2-n}\\
&=\sum_{n=-\infty}^{\infty}\left(-1\right)^nq^{32n^2-4n}
-q\sum_{n=-\infty}^{\infty}\left(-1\right)^nq^{32n^2-12n}\\
&\quad-q^3\sum_{n=-\infty}^{\infty}\left(-1\right)^nq^{32n^2+20n}
+q^6\sum_{n=-\infty}^{\infty}\left(-1\right)^nq^{32n^2+28n}\\
&=\left(-q^{28},-q^{36},q^{64};q^{64}\right)_{\infty}
-q\left(-q^{20},-q^{44},q^{64};q^{64}\right)_{\infty}\\
&\quad-q^3\left(-q^{12},-q^{52},q^{64};q^{64}\right)_{\infty}
+q^6\left(-q^{4},-q^{60},q^{64};q^{64}\right)_{\infty}.
\end{align*}
\end{proof}
\begin{corollary}
We have
\begin{align}
&\sum_{n=0}^{\infty}\frac{\left(-q^2;q^2\right)_n
\left(-q^2;q^4\right)_n}
{\left(-1;q^2\right)_{n}\left(q;q^2\right)_{n+1}
\left(q^4;q^4\right)_{n}}q^{n\left(n+1\right)}\nonumber\\
&\quad=\frac{\left(-q^2;q^2\right)_{\infty}
}
{\left(q^2;q^2\right)_{\infty}}\left(\left(q^6,q^{10},q^{16};q^{16}\right)_{\infty}
+q\left(q^2,q^{14},q^{16};q^{16}\right)_{\infty}\right),\tag{S.65}\label{S.65}\\
&\sum_{n=0}^{\infty}\frac{\left(-q;q^2\right)_{n+1}
\left(-q^2;q^4\right)_n}
{\left(q^2;q^2\right)_{2n+1}}q^{n\left(n+2\right)}\nonumber\\
&\quad=\frac{\left(-q;q^2\right)_{\infty}
\left(q^4,q^{12},q^{16};q^{16}\right)_{\infty}}{\left(q^2;q^2\right)_{\infty}
}.\tag{S.70}\label{S.70}
\end{align}
\end{corollary}
\begin{proof}
Letting  $\left(c,\beta\right)=\left(-q,q\mathrm{i}\right)$ in Theorem \ref{T5}, we get
\begin{align}
&\frac{\left(q^2,ab;q^2\right)_{\infty}}{\left(aq^2,bq^2;q^2\right)_{\infty}}
\sum_{n=0}^{\infty}
\frac{\left(q^2/a,q^2/b;q^2\right)_n\left(-q^2;q^4\right)_n}{\left(-q,-q^2,q^2;q^2\right)_{n}\left(q;q^2\right)_{n+1}}a^nb^n\notag\\
&\quad=
\sum_{n=0}^{\infty}\left(-1\right)^n\left(1+q^{4n+1}\right)\frac{\left(q^2/a,q^2/b;q^2\right)_{2n}}
{\left(aq^2,bq^2;q^2\right)_{2n}}a^{2n}b^{2n}q^{4n^2}\notag\\
&\qquad-
\sum_{n=0}^{\infty}\left(-1\right)^n\left(1+q^{4n+3}\right)\frac{\left(q^2/a,q^2/b;q^2\right)_{2n+1}}
{\left(aq^2,bq^2;q^2\right)_{2n+1}}a^{2n+1}b^{2n+1}q^{4n^2+4n+1},\label{C5}
\end{align}
where $|ab|<1$.

By setting $\left(a,b\right)\to\left(0,-1\right)$ in \eqref{C5}, we can  prove \eqref{S.65} as follows,
\begin{align*}
&\frac{\left(q^2;q^2\right)_{\infty}}{\left(-q^2;q^2\right)_{\infty}}\sum_{n=0}^{\infty}\frac{\left(-q^2;q^2\right)_n
\left(-q^2;q^4\right)_nq^{n\left(n+1\right)}}
{\left(-1;q^2\right)_{n}\left(q;q^2\right)_{n+1}
\left(q^4;q^4\right)_{n}}\nonumber\\
%&\quad=\frac{\left(q^2;q^2\right)_{\infty}}{\left(-q^2;q^2\right)_{\infty}}
%\sum_{n=0}^{\infty}\frac{\left(q\mathrm{i},-q\mathrm{i};q^2\right)_n}
%{\left(-q,q^2;q^2\right)_n\left(q;q^2\right)_{n+1}}q^{n\left(n+1\right)}\nonumber\\
&\quad=\sum_{n=0}^{\infty}\left(-1\right)^n\left(1+q^{4n+1}\right)
q^{8n^2+2n}-\sum_{n=0}^{\infty}\left(-1\right)^n
\left(1+q^{4n+3}\right)q^{8n^2+10n+3}\nonumber\\
&\quad={\left(q^6,q^{10},q^{16};q^{16}\right)_{\infty}
+q\left(q^2,q^{14},q^{16};q^{16}\right)_{\infty}}.
\end{align*}

Similarly, setting $\left(a,b\right)\to\left(0,-q\right)$ in \eqref{C5},  we can prove  \eqref{S.70} as follows,
\begin{align*}
&\sum_{n=0}^{\infty}\frac{\left(-q;q^2\right)_{n+1}
\left(-q^2;q^4\right)_n}
{\left(q^2;q^2\right)_{2n+1}}q^{n\left(n+2\right)}\\
&\quad=\frac{\left(-q;q^2\right)_{\infty}}{\left(q^2;q^2\right)_{\infty}}\left(\sum_{n=0}^{\infty}\left(-1\right)^n
q^{8n^2+4n}-\sum_{n=0}^{\infty}\left(-1\right)^n
q^{8n^2+12n+4}\right)\nonumber\\
&\quad=\frac{\left(-q;q^2\right)_{\infty}
\left(q^4,q^{12},q^{16};q^{16}\right)_{\infty}}{\left(q^2;q^2\right)_{\infty}}.
\end{align*}


\end{proof}

\section{Identities from a $_4\phi_3$ summation of Verma and Jain}
%
%\subsection{Proofs of Slater's (21) and (33)}
\begin{theorem}\label{T3}
For $\max\{|a|, |b|, |ab/q^2|\}<1$, we have
\begin{align*}
&\frac{\left(q^2,ab/q^2;q^2\right)_{\infty}}{\left(a,b;q^2\right)_{\infty}}
\sum_{n=0}^{\infty}\frac{\left(q^2/a,q^2/b;q^2\right)_n}
{\left(q^2,-q,-q^2;q^2\right)_n}\left(ab/q^2\right)^n\\
&\quad=1+\sum_{n=0}^{\infty}\left(-1\right)^n\left(1+q^n\right)
\frac{\left(q^2/a,q^2/b;q^2\right)_n}
{\left(a,b;q^2\right)_n}a^nb^nq^{\frac{3n^2-5n}{2}}.
\end{align*}
\end{theorem}
\begin{proof}
Recall an identity due to Verma and Jain \cite[Eq. (5.4)]{Verma}
\begin{equation}\label{Verma1}
_4\phi_3\left(\begin{array}{cccc}
q^{-2n},a^2q^{2n},-c, -cq\\
-aq,-aq^2,c^2
\end{array};q^2\,q^2\right)
=\frac{\left(-q,qa/c;q\right)_n\left(1+a\right)\left(-c\right)^n}
{\left(-a,c;q\right)_n\left(1+aq^{2n}\right)}.
\end{equation}
Letting $\left(a,c\right)\to\left(1,0\right)$ into \eqref{Verma1}, we have
\begin{equation*}
_3\phi_2\left(\begin{array}{ccc}
q^{-2n},q^{2n},0\\
-q,-q^2
\end{array};q^2\,q^2\right)=\frac{1+q^n}{1+q^{2n}}
q^{\frac{n\left(n+1\right)}{2}}.
\end{equation*}
Based on the above summation formula, we substitute $\left(q,\alpha,\beta,c,d\right)\to\left(q^2,1,0,-q,-q^2\right)$ into \eqref{Formula1} to obtain the desired result.
\end{proof}
\begin{corollary}
We have
\begin{align}
&\sum_{n=0}^{\infty}
\frac{\left(q;q^2\right)_n\left(-1\right)^nq^{n^2}}
{\left(-q;q^2\right)_n\left(q^4;q^4\right)_n}
=\frac{\left(-q^2,-q^3,q^5;q^5\right)_{\infty}
\left(q;q^2\right)_{\infty}}{\left(q^2;q^2\right)_{\infty}}, \tag{S.21}\label{S.21}\\
&\sum_{n=0}^{\infty}
\frac{q^{2n^2}}
{\left(q^2;q^2\right)_n\left(-q;q\right)_{2n}}
=\frac{\left(q^3,q^4,q^7;q^7\right)_{\infty}}
{\left(q^2;q^2\right)_{\infty}}.\tag{S.33}\label{S.33}
\end{align}
\end{corollary}
\begin{proof}
Taking $\left(a,b\right)\to\left(0,q\right)$ into Theorem \ref{T3} and using Jacobi triple product identity \eqref{Jacobi} with $\left(q,z\right)\to\left(q^5,q^2\right)$, we  obtain \eqref{S.21}.

Similarly, inserting $\left(a,b\right)\to\left(0,0\right)$ into Theorem \ref{T3} and applying Jacobi triple product identity \eqref{Jacobi} with $\left(q,z\right)\to\left(q^7,-q^3\right)$, we get \eqref{S.33}.
\end{proof}
Next, we generalize  Theorem \ref{T3} into the following form  which was first provided by  Liu in \cite[Proposition 11.4]{Liu_4}.
\begin{theorem}\label{T6}
For $\max\{|\beta|,|\alpha a|, |\alpha b|,  |\alpha^{1/2}q^2|, |\alpha ab/q^4|\}<1$, we have
\begin{align*}
&\frac{\left(\alpha q^4,\alpha ab/q^4;q^4\right)_{\infty}}{\left(\alpha a,\alpha b;q^4\right)_{\infty}}
\sum_{n=0}^{\infty}
\frac{\left(q^4/a,q^4/b,\beta,\beta q^2;q^4\right)_n}
{\left(q^4,-\alpha^{1/2}q^2,-\alpha^{1/2}q^4,\beta^2;q^4\right)_n}\left(\alpha ab/q^4\right)^n\\
&\quad=
\sum_{n=0}^{\infty}
\left(-1\right)^n
\frac{1-\alpha^{1/2}q^{4n}}{1-\alpha^{1/2}}
\frac{\left(\alpha^{1/2},-\alpha^{1/2}q^2/\beta;q^2\right)_n
\left(q^4/a,q^4/b;q^4\right)_n}
{\left(-\beta,q^2;q^2\right)_n\left(\alpha a,\alpha b;q^4\right)_n}\left(\alpha\beta ab\right)^nq^{2n^2-6n}.
\end{align*}
\end{theorem}
\begin{proof}
Replacing $q$ by $q^4$ and substituting $\left(c,d,h,z,\gamma\right)\to
\left(-\alpha^{1/2}q^2,\alpha^{1/2}q^4,\beta^2,1,\beta q^2\right)$ in \eqref{Formula2}, we deduce that
\begin{align*}
&\frac{\left(\alpha q^4,\alpha ab/q^4;q^4\right)_{\infty}}{\left(\alpha a,\alpha b;q^4\right)_{\infty}}
\sum_{n=0}^{\infty}
\frac{\left(q^4/a,q^4/b,\beta,\beta q^2;q^4\right)_n}
{\left(q^4,-\alpha^{1/2}q^2,-\alpha^{1/2}q^4,\beta^2;q^4\right)_n}\left(\alpha ab/q^4\right)^n\\
&\quad=
\sum_{n=0}^{\infty}\left(-1\right)^n\frac{1-\alpha q^{8n}}{1-\alpha}
\frac{\left(\alpha,q^4/a,q^4/b;q^4\right)_n}
{\left(q^4,\alpha a,\alpha b;q^4\right)_n}\left(\alpha ab\right)^nq^{2n^2-6n}\\
&\qquad \times
{_4\phi_3}\left(\begin{array}{cccc}
q^{-4n},\alpha q^{4n},\beta,\beta q^2\\
-\alpha^{1/2}q^2,-\alpha^{1/2}q^4,\beta^2
\end{array};q^4,\,q^4\right)\\
&\quad=\sum_{n=0}^{\infty}
\left(-1\right)^n
\frac{1-\alpha^{1/2}q^{4n}}{1-\alpha^{1/2}}
\frac{\left(\alpha^{1/2},-\alpha^{1/2}q^2/\beta;q^2\right)_n
\left(q^4/a,q^4/b;q^4\right)_n}
{\left(-\beta,q^2;q^2\right)_n\left(\alpha a,\alpha b;q^4\right)_n}\left(\alpha\beta ab\right)^nq^{2n^2-6n},
\end{align*}
where we have used \eqref{Verma1} with the substitution $\left(q,a,c\right)\to\left(q^2,\alpha^{1/2},-\beta\right)$.
\end{proof}
\begin{corollary}
We have
\begin{align}
\sum_{n=0}^{\infty}\frac{\left(q;q^2\right)_{2n}}
{\left(q^4;q^4\right)_{2n}}q^{4n^2}
&=\frac{\left(q^5,q^7,q^{12};q^{12}\right)_{\infty}}
{\left(q^4;q^4\right)_{\infty}},\tag{S.53}
\label{S.53}\\
\sum_{n=0}^{\infty}\frac{\left(-1\right)^n\left(q;q^2\right)_{2n}}
{\left(-q^2;q^4\right)_{n}\left(q^8;q^8\right)_{n}}q^{2n^2}
&=\frac{\left(q^2;q^4\right)_{\infty}\left(-q^3,-q^5,q^{8};q^{8}\right)_{\infty}}
{\left(q^4;q^4\right)_{\infty}},\tag{\cite[Eq. (2.7)]{Bowman}}\label{N7}\\
\sum_{n=0}^{\infty}\frac{\left(-q;q^2\right)_{2n}}
{\left(q^2;q^4\right)_{n}\left(q^8;q^8\right)_{n}}q^{2n^2}
&=\frac{\left(-q^2;q^4\right)_{\infty}\left(-q^3,-q^5,q^{8};q^{8}\right)_{\infty}}
{\left(q^4;q^4\right)_{\infty}}.\tag{\rm{\cite[Eq. (7.24)]{Gessel}}}\label{N8}
\end{align}
\end{corollary}

\begin{proof}
Taking $\left(\alpha,\beta\right)\to\left(1,q\right)$ into Theorem \ref{T6}, we see that
\begin{align}
&\frac{\left(q^4,ab/q^4;q^4\right)_{\infty}}{\left(a,b;q^4\right)_{\infty}}
\sum_{n=0}^{\infty}
\frac{\left(q,q^3,q^4/a,q^4/b;q^4\right)_n}{\left(q^4,-q^4,q^2,-q^2;q^4\right)_n}\left(ab/q^4\right)^n\nonumber\\
&\quad=1+\sum_{n=1}^{\infty}\left(-1\right)^n\left(1+q^{2n}\right)
\frac{\left(q^4/a,q^4/b;q^4\right)_n}{\left(a,b;q^4\right)_n}
a^nb^nq^{2n^2-5n}.\label{C8}
\end{align}
Then, taking $\left(a,b\right)\to\left(0,0\right)$ and $\left(a,b\right)\to\left(0,q^2\right)$ in \eqref{C8}, we arrive at  \eqref{S.53} and \ref{N7}.

In addition, taking $\left(a,b\right)\to\left(0,-q^2\right)$ in \eqref{C8}, we get
\begin{align*}
\sum_{n=0}^{\infty}\frac{\left(q;q^2\right)_{2n}}
{\left(q^2;q^4\right)_{n}\left(q^8;q^8\right)_{n}}q^{2n^2}
&=\frac{\left(-q^2;q^4\right)_{\infty}\left(q^3,q^5,q^{8};q^{8}\right)_{\infty}}
{\left(q^4;q^4\right)_{\infty}}.
\end{align*}
Replacing $q$ by $-q$ in the above identity, we immediately obtain \ref{N8}.
\end{proof}

\begin{corollary}
We have
\begin{align}
%&\sum_{n=0}^{\infty}\frac{\left(-q;q^2\right)_{2n}}
%{\left(q^4;q^4\right)_{2n}}q^{4n^2}
%=\frac{\left(-q^5,-q^7,q^{12};q^{12}\right)_{\infty}}
%{\left(q^4;q^4\right)_{\infty}},\label{new5}\\
\sum_{n=0}^{\infty}\frac{\left(q;q^2\right)_{2n+1}}
{\left(q^4;q^4\right)_{2n+1}}q^{4n^2+4n}
&=\frac{\left(q,q^{11},q^{12};q^{12}\right)_{\infty}}{\left(q^4;q^4\right)_{\infty}},
\tag{S.55}\label{S.55}\\
\sum_{n=0}^{\infty}\frac{\left(-q;q^2\right)_{2n+1}}
{\left(q^4;q^4\right)_{2n+1}}q^{4n^2+4n}
&=\frac{\left(-q,-q^{11},q^{12};q^{12}\right)_{\infty}}{\left(q^4;q^4\right)_{\infty}},
\tag{S.57}\label{S.57}
\\
\sum_{n=0}^{\infty}\frac{\left(q;q^2\right)_{2n+1}\left(-q^4;q^4\right)_n}
{\left(q^4;q^4\right)_{2n+1}}q^{2n^2+2n}
&=\frac{\left(-q^4;q^4\right)_{\infty}\left(q,q^7,q^{8};q^{8}\right)_{\infty}}
{\left(q^4;q^4\right)_{\infty}}.\tag{\rm{\cite[Eq. (7.25)]{Gessel}}}\label{N9}
\end{align}
\end{corollary}

\begin{proof}
Taking $\left(\alpha,\beta\right)\to\left(q^4,q^3\right)$ into Theorem \ref{T6}, we obtain
\begin{align}
&\frac{\left(q^4,ab;q^4\right)_{\infty}}{\left(aq^4,bq^4;q^4\right)_{\infty}}
\sum_{n=0}^{\infty}\frac{\left(q;q^2\right)_{2n+1}\left(q^4/a,q^4/b;q^4\right)_n}{\left(q^4;q^4\right)_{2n+1}}a^nb^n\nonumber\\
&\quad=
\sum_{n=0}^{\infty}\left(-1\right)^n\left(1-q^{2n+1}\right)
\frac{\left(q^4/a,q^4/b;q^4\right)_n}{\left(aq^4,bq^4;q^4\right)_n}a^nb^nq^{2n^2+n}.\label{C9}
\end{align}

Setting $\left(a,b\right)\to\left(0,0\right)$ in \eqref{C9},
 we obtain \eqref{S.55}. When replacing $q$ by $-q$, \eqref{S.55} becomes to \eqref{S.57}.
 
Besides, inserting  $\left(a,b\right)\to\left(0,-1\right)$ in \eqref{C9},  we obtain \ref{N9} .
\end{proof}

%\begin{corollary}
%We have
%\begin{align}
%\sum_{n=0}^{\infty}\frac{\left(-q;q^2\right)_{2n+1}}
%{\left(q^4;q^4\right)_{2n+1}}q^{4n^2+4n}
%&=\frac{\left(-q,-q^{11},q^{12};q^{12}\right)_{\infty}}
%{\left(q^4;q^4\right)_{\infty}}.
%\tag{S.57}\label{S.57}
%\end{align}
%\end{corollary}


%Taking $\left(a,b,\alpha,\beta\right)\to\left(0,0,1,-q\right)$ into Theorem \ref{T6}, we deduce that
%\begin{align*}
%&\left(q^4;q^4\right)_{\infty}
%\sum_{n=0}^{\infty}
%\frac{\left(-q,-q^3;q^4\right)_n}{\left(q^4,-q^4,q^2,-q^2;q^4\right)_n}q^{4n^2}\\
%&=1+\sum_{n=1}^{\infty}\left(1+q^{2n}\right)q^{6n^2-n}\\
%&=\sum_{n=-\infty}^{\infty}q^{6n^2-n}\\
%&=\left(-q^5,-q^7,q^{12};q^{12}\right)_{\infty},
%\end{align*}
%which means
%\begin{align*}
%\sum_{n=0}^{\infty}\frac{\left(-q;q^2\right)_{2n}}
%{\left(q^4;q^4\right)_{2n}}q^{4n^2}
%=\sum_{n=0}^{\infty}
%\frac{\left(-q,-q^3;q^4\right)_n}{\left(q^4,-q^4,q^2,-q^2;q^4\right)_n}q^{4n^2}
%=\frac{\left(-q^5,-q^7,q^{12};q^{12}\right)_{\infty}}{\left(q^4;q^4\right)_{\infty}}.
%\end{align*}
%Therefore, we conclude that \eqref{new5} is true.
%\begin{proof}
%By taking $\left(\alpha,\beta\right)\to\left(q^4,-q^3\right)$ into Theorem \ref{T6}, we deduce that
%\begin{align}
%&\frac{\left(q^4,ab;q^4\right)_{\infty}}{\left(aq^4,bq^4;q^4\right)_{\infty}}
%\sum_{n=0}^{\infty}
%\frac{\left(q^4/a,q^4/b;q^4\right)_n\left(-q;q^2\right)_{2n+1}}{\left(q^4;q^4\right)_{2n+1}}a^nb^n\nonumber\\
%&\quad
%=
%\sum_{n=0}^{\infty}\left(1+q^{2n+1}\right)\frac{\left(q^4/a,q^4/b;q^4\right)_n}{\left(aq^4,bq^4;q^4\right)_n}
%a^nb^nq^{2n^2+n}.\label{C10}
%\end{align}
%Then, substituting $\left(a,b\right)\to\left(0,0\right)$ into \eqref{C10},we obtain \eqref{S.57}.
%\end{proof}
\section{Identities from a ${_4\phi_3}$ summation of Andrews}
%\subsection{Proof of Slater's (35), (54), (68), (69), (71), (102), (104) and (106)}
In  \cite{WangChun}, Wang and Chern established  the following $q$-series transformation formula.
\begin{theorem}\label{T4}
For $\max\{|\alpha a|,|\alpha b|,|\beta^2|,|\sqrt{\alpha q}|,|\alpha ab/q|\}<1$, we have
\begin{align*}
&\frac{\left(\alpha q,\alpha ab/q;q\right)_{\infty}}
{\left(\alpha a,\alpha b;q\right)_{\infty}}\sum_{n=0}^{\infty}
\frac{\left(q/a,q/b,\beta,-\beta;q\right)_{n}}
{\left(q,\sqrt{\alpha q},-\sqrt{\alpha q},\beta^2;q\right)_n}\left(\alpha ab/q\right)^n\nonumber\\
&\quad=\sum_{n=0}^{\infty}
\frac{\left(1-\alpha q^{4n}\right)\left(\alpha,q/a,q/b;q\right)_{2n}\left(q,\alpha q/\beta^2;q^2\right)_n}{\left(1-\alpha\right)
\left(q,\alpha a,\alpha b;q\right)_{2n}\left(\alpha q,\beta^2q;q^2\right)_n}\left(\alpha\beta ab\right)^{2n}q^{2n^2-3n}.
\end{align*}
\end{theorem}
To ensure the completeness of the article's content, we provide a proof of the above result. 
\begin{proof}[Proof of Theorem \ref{T4}]
Recall a summation formula due to Andrew \cite{Andrews1}:
\begin{align}
_4\phi_3\left(\begin{array}{cccc}
q^{-n},a^2q^{n+1},c,-c\\
aq,-aq,c^2
\end{array};q,\,q\right)=\left\{\begin{array}{ll}
0,&\quad n\equiv1\pmod2,\\
\frac{c^n\left(q,a^2q^2/c^2;q^2\right)_{n/2}}
{\left(a^2q^2,c^2q;q^2\right)_{n/2}},&\quad n\equiv0\pmod2.
\end{array}\right.\label{Andrews}
\end{align}
Taking $a=\sqrt{\alpha/q}, c=\beta$ into \eqref{Andrews}, we obtain
\begin{align*}
_4\phi_3\left(\begin{array}{cccc}
q^{-n},\alpha q^n,\beta,-\beta\\
\sqrt{\alpha q},-\sqrt{\alpha q},\beta^2
\end{array};q,\,q\right)
=\left\{\begin{array}{ll}
0,&\quad n\equiv1\pmod2,\\
\frac{\beta^n\left(q,\alpha q/\beta^2;q^2\right)_{n/2}}
{\left(\alpha q,\beta^2q;q^2\right)_{n/2}},& \quad{n\equiv0\pmod2}.
\end{array}\right.
\end{align*}
Based on the above summation formula, we take $c=-d=\sqrt{\alpha q}, \beta=-\gamma, h=\beta^2$ and $z=1$ into \eqref{Formula2} to get the desired result.
\end{proof}

%\subsection{Proof of Slater's (61) and (98)}
% For the reader's convenience, we include a short proof here.
\begin{corollary}
We have
\begin{align}
\sum_{n=0}^{\infty}\frac{q^{n^2}}{\left(q;q\right)_n
\left(q;q^2\right)_{n}}
&=\frac{\left(q^6,q^8,q^{14};q^{14}\right)_{\infty}}
{\left(q;q\right)_{\infty}},\tag{S.61}\label{S.61}\\
\sum_{n=0}^{\infty}\frac{q^{n^2}}{\left(q;q\right)_{2n}}
&=\frac{\left(q^2,q^8,q^{10};q^{10}\right)_{\infty}
\left(q^6,q^{14};q^{20}\right)_{\infty}}
{\left(q;q\right)_{\infty}}.\tag{S.98}\label{S.98}
%&\sum_{n=0}^{\infty}\frac{\left(-1\right)^nq^{n^2}}{\left(-q,q^2;q^2\right)_n}
%=\frac{\left(q^8,q^{12},q^{20};q^{20}\right)_{\infty}}
%{\left(-q;q\right)_{\infty}\left(q^2;q^2\right)_{\infty}}.\label{new9}
\end{align}
\end{corollary}
\begin{proof}
Taking $\left(a,c\right)\to\left(q^{-1/2},0\right)$ into Andrews's formula \eqref{Andrews}, we have
\begin{equation*}
_3\phi_2\left(
\begin{array}{ccc}
q^{-n},q^{n},0\\
q^{1/2},-q^{1/2}
\end{array};q,\,q\right)
=\left\{\begin{array}{ll}
0,&\quad n\equiv1\pmod2,\\
\left(-1\right)^{{n}/{2}}
q^{{n^2}/{4}},&\quad n\equiv0\pmod2.
\end{array}
\right.
\end{equation*}
Based on the above summation, we fix $\alpha=1,\beta=0$ in  Theorem \ref{T4} to obtain
\begin{align}
&\frac{\left(q,ab/q;q\right)_{\infty}}{\left(a,b;q\right)_{\infty}}
\sum_{n=0}^{\infty}\frac{\left(q/a,q/b;q\right)_{n}}
{\left(q;q\right)_{n}\left(q;q^2\right)_{n}}
\left(ab/q\right)^n\nonumber\\
&\quad=1+\sum_{n=1}^{\infty}\left(-1\right)^n\left(1+q^{2n}\right)
\frac{\left(q/a,q/b;q\right)_{2n}}{\left(a,b;q\right)_{2n}}
a^{2n}b^{2n}q^{3n^2-3n}.\label{C17}
\end{align}
where $\max\{|a|,|b|,|ab/q|\}<1$.

Substituting $\left(a,b\right)\to\left(0,0\right)$ into \eqref{C17} and taking $\left(q,z\right)\to\left(q^{14},-q^6\right)$ in Jacobi triple product identity \eqref{Jacobi}, we obtain \eqref{S.61}.

Similarly, taking $\left(a,b\right)\to\left(0,-q^{1/2}\right)$ into \eqref{C17} and taking $\left(q,z\right)\to\left(q^{10},q^4\right)$ in Jacobi triple product identity \eqref{Jacobi}, we obtain
\begin{align}
\frac{\left(q;q\right)_{\infty}}{\left(-q^{1/2};q\right)_{\infty}}
\sum_{n=0}^{\infty}\frac{\left(-q^{1/2};q\right)_n}
{\left(q;q\right)_n\left(q;q^2\right)_{n}}q^{{n^2}/{2}}
=\left(q^4,q^6,q^{10};q^{10}\right)_{\infty}.\label{e5}
\end{align}
We further replace $q$ by $q^2$ in \eqref{e5}, which, after simplification, results in \eqref{S.98}.
%Taking $\left(a,b\right)\to\left(0,q^{1/2}\right)$ into Theorem \ref{thm6} and taking $\left(q,z\right)\to\left(q^{10},q^4\right)$ in Jacobi triple product identity \eqref{Jacobi}, we obtain
%\begin{align}
%\frac{\left(q;q\right)_{\infty}}{\left(q^{1/2};q\right)_{\infty}}
%\sum_{n=0}^{\infty}\frac{\left(-1\right)^nq^{\frac{n^2}{2}}}
%{\left(q;q\right)_n\left(-q^{1/2};q^2\right)_{n}}
%=\left(q^4,q^6,q^{10};q^{10}\right)_{\infty}.\label{eq5}
%\end{align}
%We further replace $q$ by $q^2$ in \eqref{eq5} to arrive at \eqref{new9} after simplification.
\end{proof}
%It is easy to see that Theorem \ref{thm6} is a special case of Liu's result \cite[proposition 11.6]{Liu_4}. Next, we establish a generalization of Theorem \ref{thm6}.

\begin{corollary}
We have
\begin{align}
\sum_{n=0}^{\infty}\frac{
\left(-q;q\right)_{n}\left(-q;q^2\right)_{n}
}{\left(q;q\right)_{2n+1}}q^{\frac{n\left(n+3\right)}{2}}
&=\frac{\left(-q;q\right)_{\infty}\left(q,q^7,q^8;q^8\right)_{\infty}}
{\left(q;q\right)_{\infty}}.\tag{S.35\,and\,S.106} \label{S.35}
\end{align}
\end{corollary}
\begin{proof}
When $\alpha=q^2,\beta=q^{1/2}{\rm i}$, Theorem \ref{T4} reduces to 
\begin{align}
&\frac{\left(q,abq;q\right)_{\infty}}{\left(aq^2,bq^2;q\right)_{\infty}}
\sum_{n=0}^{\infty}\frac{\left(q/a,q/b\right)_n\left(-q;q^2\right)_n}{\left(q;q\right)_{2n+1}}a^nb^nq^n
\nonumber\\
&\quad=
\sum_{n=0}^{\infty}\left(-1\right)^n\left(1-q^{4n+2}\right)\frac{\left(q/a,q/b;q\right)_{2n}}
{\left(aq^2,bq^2;q\right)_{2n}}\left(ab\right)^{2n}q^{2n^2+2n}.\label{C11}
\end{align}
By taking $\left(a,b\right)\to\left(0,-1\right)$ into \eqref{C11}, we obtain \eqref{S.35}.
\end{proof}
\begin{corollary}
We have
\begin{align}
\sum_{n=0}^{\infty}\frac{\left(-q;q\right)_n}{\left(q;q\right)_n
\left(q;q^2\right)_{n+1}}q^{\frac{n\left(n+3\right)}{2}}
&=\frac{\left(-q;q\right)_{\infty}\left(q,q^9,q^{10};q^{10}\right)_{\infty}}
{\left(q;q\right)_{\infty}},\tag{S.43}\label{S.43}
\\
\sum_{n=0}^{\infty}\frac{q^{n\left(n+2\right)}}
{\left(q;q\right)_n\left(q;q^2\right)_{n+1}}
&=\frac{\left(q^2,q^{12},q^{14};q^{14}\right)_{\infty}}{\left(q;q\right)_{\infty}}
,\tag{S.59}\label{S.59}
\end{align}
\end{corollary}
%
\begin{proof}
When $\alpha=q^2,\beta=0$, Theorem \ref{T4} reduces to 
\begin{align}
&\frac{\left(q,abq;q\right)_{\infty}}{\left(aq^2,bq^2;q\right)_{\infty}}
\sum_{n=0}^{\infty}\frac{\left(q/a,q/b\right)_n}{\left(q;q\right)_{n}\left(q;q^2\right)_{n+1}}a^nb^nq^n
\nonumber\\
&\quad=
\sum_{n=0}^{\infty}\left(-1\right)^n\left(1-q^{4n+2}\right)\frac{\left(q/a,q/b;q\right)_{2n}}
{\left(aq^2,bq^2;q\right)_{2n}}\left(ab\right)^{2n}q^{3n^2+3n}.\label{C12}
\end{align}
By taking $\left(a,b\right)\to\left(0,-1\right)$ and $\left(a,b\right)\to\left(0,0\right)$ into \eqref{C12}, respectively,  we obtain \eqref{S.43} and  \eqref{S.59}.
\end{proof}
\begin{corollary}
We have
\begin{align}
&\sum_{n=0}^{\infty}\frac{\left(-q^2;q^2\right)_{n-1}\left(1+q^n\right)}
{\left(q;q\right)_{2n}}q^{n^2}
=\frac{\left(q^5,q^7,q^{12};q^{12}\right)_{\infty}}
{\left(q;q\right)_{\infty}},\tag{S.54} \label{S.54}\\
&\sum_{n=0}^{\infty}\frac{\left(-q^4;q^4\right)_{n-1}\left(-q;q^2\right)_nq^{n^2}}
{\left(q^2;q^4\right)_n\left(q^2;q^2\right)_n\left(-q^2;q^2\right)_{n-1}}
=
\frac{\left(-q;q^2\right)_{\infty}\left(q^6,q^{10},q^{16};q^{16}\right)_{\infty}}
{\left(q^2;q^2\right)_{\infty}},\tag{S.71}\label{S.71}\\
&\sum_{n=0}^{\infty}\frac{
\left(-q;q\right)_{n}\left(-q^2;q^2\right)_{n-1}
}{\left(q;q\right)_{2n}}q^{\frac{n\left(n+1\right)}{2}}
\nonumber\\
&\quad
=\frac{\left(-q;q\right)_{\infty}}
{\left(q;q\right)_{\infty}}\left(\left(-q^{16},-q^{16},q^{32};q^{32}\right)_{\infty}
-q\left(-q^{8},-q^{24},q^{32};q^{32}\right)_{\infty}\right).\tag{S.104} \label{S.104}
\end{align}
\end{corollary}
\begin{proof}
When $\alpha=1,\beta={\rm i}$, Theorem \ref{T4} reduces to 
\begin{align}
&\frac{\left(q,ab/q;q\right)_{\infty}}{\left(a,b;q\right)_{\infty}}
\sum_{n=0}^{\infty}\frac{\left(q/a,q/b;q\right)_n\left(-1;q^2\right)_n}
{\left(q,-1;q\right)_{n}\left(q;q^2\right)_{n}}\left(ab/q\right)^n
\nonumber\\
&\quad=1+
\sum_{n=1}^{\infty}\left(-1\right)^n\left(1+q^{2n}\right)\frac{\left(q/a,q/b;q\right)_{2n}}
{\left(a,b;q\right)_{2n}}\left(ab\right)^{2n}q^{2n^2-3n}.\label{C13}
\end{align}
Then,  by taking $\left(a,b\right)\to\left(0,0\right)$ into \eqref{C13}, we can prove \eqref{S.54} as follows,
\begin{align*}
&\sum_{n=0}^{\infty}\frac{\left(-q^2;q^2\right)_{n-1}\left(1+q^n\right)}
{\left(q;q\right)_{2n}}q^{n^2}=
\sum_{n=0}^{\infty}\frac{\left(-1;q^2\right)_n}
{\left(q,-1;q\right)_{n}\left(q;q^2\right)_{n}}q^{n^2}=\frac{\left(q^5,q^7,q^{12};q^{12}\right)_{\infty}}{\left(q;q\right)_{\infty}}.
\end{align*}

Similarly, replacing $q$ by $q^2$ and inserting $\left(a,b\right)\to\left(0,-q\right)$ into \eqref{C13}, we arrive at \eqref{S.71}.

In addition, taking $\left(a,b\right)\to\left(0,0\right)$ into \eqref{C13}, we obtain 
\begin{align*}
&\sum_{n=0}^{\infty}\frac{
\left(-q;q\right)_{n}\left(-q^2;q^2\right)_{n-1}
}{\left(q;q\right)_{2n}}q^{\frac{n\left(n+1\right)}{2}}\\
&\quad
=
\frac{1}{2}+\frac{1}{2}\sum_{n=0}^{\infty}\frac{\left(-1;q^2\right)_n}
{\left(q;q\right)_n\left(q;q^2\right)_n}q^{\frac{n\left(n+1\right)}{2}}\\
&\quad=\frac{1}{2}+\frac{1}{2}\frac{\left(-q;q\right)_{\infty}\left(q^4,q^4,q^8;q^8\right)_{\infty}}{\left(q;q\right)_{\infty}}\\
&\quad=\frac{\left(-q;q\right)_{\infty}}
{\left(q;q\right)_{\infty}}\left(\left(-q^{16},-q^{16},q^{32};q^{32}\right)_{\infty}
-q\left(-q^{8},-q^{24},q^{32};q^{32}\right)_{\infty}\right)
\end{align*}
where we have used the following relations,
\begin{align*}
\frac{\left(q;q\right)_{\infty}}{\left(-q;q\right)_{\infty}}
&=\sum_{n=-\infty}^{\infty}\left(-1\right)^nq^{n^2}\\
&=\sum_{n=-\infty}^{\infty}q^{16n^2}
-2q\sum_{n=-\infty}^{\infty}q^{16n^2+8n}+q^4
\sum_{n=-\infty}^{\infty}q^{16n^2+16n}
\\
&=\left(-q^{16},-q^{16},q^{32};q^{32}\right)_{\infty}
-2q\left(-q^{8},-q^{24},q^{32};q^{32}\right)_{\infty}
+q^4\left(-1,-q^{32},q^{32};q^{32}\right)_{\infty}
\end{align*}
and
\begin{align*}
\left(q^4,q^4,q^8;q^8\right)_{\infty}=\left(-q^{16},-q^{16},q^{32};q^{32}\right)_{\infty}
-q^4\left(-1,-q^{32},q^{32};q^{32}\right)_{\infty}.
\end{align*}
Therefore, we prove \eqref{S.104} is true.
\end{proof}
%
\begin{corollary}
We have
\begin{align}
\sum_{n=0}^{\infty}\frac{\left(-q^2;q^2\right)_n}{\left(q;q\right)_{2n+1}\left(1+q^{n+1}\right)}q^{n^2+2n}
&=\frac{\left(q,q^{11},q^{12};q^{12}\right)_{\infty}}{\left(q;q\right)_{\infty}},
\tag{S.49}
\label{S.49}\\
\sum_{n=0}^{\infty}\frac{\left(-q^2;q^2\right)_n}{\left(q;q\right)_n\left(q;q^2\right)_{n+1}}q^{\frac{n\left(n+1\right)}{2}}
&=\frac{\left(-q;q\right)_{\infty}\left(q^2,q^6,q^8;q^8\right)_{\infty}}{\left(q;q\right)_{\infty}},\tag{\rm{\cite[Entry. 1.7.5]{Ramanujan}}}\label{N11}\\
\sum_{n=0}^{\infty}\frac{\left(-q;q^2\right)_{n+1}\left(-q^4;q^4\right)_n
q^{n^2+2n}}
{\left(-q^2;q^2\right)_{n+1}\left(q^2;q^2\right)_n\left(q^2;q^4\right)_{n+1}}
&=\frac{\left(-q;q^2\right)_{\infty}\left(q^2,q^{14},q^{16};q^{16}\right)_{\infty}}
{\left(q^2;q^2\right)_{\infty}}.\tag{S.68}\label{S.68}
\end{align}
\end{corollary}
\begin{proof}
When $\alpha=q^2,\beta=q{\rm i}$, Theorem \ref{T4} reduces to 
\begin{align}
&\frac{\left(q,abq;q\right)_{\infty}}{\left(aq^2,bq^2;q\right)_{\infty}}
\sum_{n=0}^{\infty}\frac{\left(q/a,q/b;q\right)_n\left(-q^2;q^2\right)_n}
{\left(q;q\right)_{n}\left(-q;q\right)_{n+1}\left(q;q^2\right)_{n+1}}\left(abq\right)^n
\nonumber\\
&\quad=
\sum_{n=0}^{\infty}\left(-1\right)^n\left(1-q^{2n+1}\right)\frac{\left(q/a,q/b;q\right)_{2n}}
{\left(aq^2,bq^2;q\right)_{2n}}\left(ab\right)^{2n}q^{2n^2+3n}.\label{C14}
\end{align}

Then, we take $\left(a,b\right)\to\left(0,0\right)$ and  $\left(a,b\right)\to\left(0,-1/q\right)$ into \eqref{C14}, respectively,  to obtain \eqref{S.49} and \ref{N11}.

In addition, replacing $q$ by $q^2$ and  setting  $\left(a,b\right)\to\left(0,-1/q\right)$ to \eqref{C14}, we get \eqref{S.68}.
\end{proof}

%
%
\begin{corollary}
We have
\begin{align}
\sum_{n=0}^{\infty}\frac{\left(-q;q^2\right)_{n+1}\left(q^4;q^4\right)_n
q^{n^2+2n}}
{\left(q^2;q^2\right)_{n+1}\left(q^2;q^2\right)_n\left(q^2;q^4\right)_{n+1}}
&=\frac{\left(-q;q^2\right)_{\infty}\left(-q^2,-q^{14},q^{16};q^{16}\right)_{\infty}}
{\left(q^2;q^2\right)_{\infty}}.\tag{S.69}\label{S.69}
%\\
%\sum_{n=0}^{\infty}\frac{
%\left(-q;q\right)_{n+1}\left(-q^2;q^2\right)_{n}
%}{\left(q;q\right)_{2n+2}}q^{\frac{n\left(n+1\right)}{2}}
%&=\frac{\left(-q;q\right)_{\infty}\left(q^2,q^6,q^8;q^8\right)_{\infty}}
%{\left(1+q\right)\left(q;q\right)_{\infty}}.\tag{S.102} \label{S.102}
\end{align}
\end{corollary}
\begin{proof}
\begin{comment}
Substituting $\left(\alpha,\beta,a,b\right)\to\left(1,\mathrm{i},0,0\right)$ into Theorem \ref{T4} and using the Jacobi triple product identity \eqref{Jacobi} with $\left(q,z\right)\to\left(q^{12},-q^5\right)$, we get
\begin{align*}
&\left(q;q\right)_{\infty}
\sum_{n=0}^{\infty}\frac{\left(\mathrm{i},-\mathrm{i};q\right)_n}
{\left(q,-1,q^{1/2},-q^{1/2};q\right)_n}q^{n^2}=1+
\sum_{n=1}^{\infty}\left(-1\right)^n\left(1+q^{2n}\right)
q^{6n^2-n}=\left(q^5,q^7,q^{12};q^{12}\right)_{\infty}.
\end{align*}
Then, we can prove \eqref{S.54} as follows,
\begin{align*}
&1+\sum_{n=1}^{\infty}\left(1+q^n\right)\frac{\left(-q;q^2\right)_{n-1}}
{\left(q;q\right)_{2n}}q^{n^2}\\
&\quad=\sum_{n=0}^{\infty}\frac{\left(-1;q^2\right)_n}
{\left(-1,q;q\right)_n\left(q;q^2\right)_n}q^{n^2}\\
&\quad=\sum_{n=0}^{\infty}\frac{\left(\mathrm{i},-\mathrm{i};q\right)_n}
{\left(q,-1,q^{1/2},-q^{1/2};q\right)_n}q^{n^2}\nonumber\\
&\quad=\frac{\left(q^5,q^7,q^{12};q^{12}\right)_{\infty}}{\left(q;q\right)_{\infty}}.
\end{align*}

Replacing $q$ by $q^2$ and setting $\left(\alpha,\beta,a,b\right)\to\left(q^4,q^2\mathrm{i},0,-1/q\right)$ in Theorem \ref{T4} and using the Jacobi triple product identity \eqref{Jacobi} with $\left(q,z\right)\to\left(q^{16},-q^2\right)$, we get
\begin{align*}
\frac{\left(q^6;q^2\right)_{\infty}}{\left(-q^3;q^2\right)_{\infty}}
\sum_{n=0}^{\infty}\frac{\left(-q^4;q^4\right)_n}
{\left(-q^4;q^2\right)_n\left(q^2;q\right)_{2n}}q^{n^2+2n}
=\frac{\left(q^2,q^{14},q^{16};q^{16}\right)_{\infty}}
{1-q^2}.
\end{align*}
Then, we  prove \eqref{S.68} as follows,
\begin{align*}
&\sum_{n=0}^{\infty}
\frac{\left(-q;q^2\right)_{n+1}\left(-q^4;q^4\right)_n}
{\left(-q^2;q^2\right)_{n+1}\left(q^2;q^2\right)_n\left(q^2;q^4\right)_{n+1}}q^{n^2+2n}
\\
&\quad=\frac{1}{\left(1-q\right)\left(1+q^2\right)}
\sum_{n=0}^{\infty}\frac{\left(-q^4;q^4\right)_n}
{\left(-q^4;q^2\right)_n\left(q^2;q\right)_{2n}}q^{n^2+2n}\\
&\quad=\frac{\left(-q;q^2\right)_{\infty}\left(q^2,q^{14},q^{16};q^{16}\right)_{\infty}}
{\left(q^2;q^2\right)_{\infty}}.
\end{align*}
\end{comment}
Putting $\left(\alpha,\beta\right)\to\left(q^2,q\right)$ in Theorem \ref{T4}, we get
\begin{align}
&\frac{\left(q,abq;q\right)_{\infty}}{\left(aq^2,bq^2;q\right)_{\infty}}
\sum_{n=0}^{\infty}\frac{\left(q/a,q/b;q\right)_n\left(q^2;q^2\right)_n}
{\left(q;q\right)_n\left(q;q\right)_{n+1}\left(q;q^2\right)_{n+1}}\left(abq\right)^n\nonumber\\
&\quad=\sum_{n=0}^{\infty}\left(1+q^{2n+1}\right)\frac{\left(q/a,q/b;q\right)_{2n}}{\left(aq^2,bq^2;q\right)_{2n}}
\left(ab\right)^{2n}q^{2n^2+3n}.\label{C15}
\end{align}
Then, we  can obtain \eqref{S.69} by replacing $q$ by $q^2$ and letting $\left(a,b\right)\to\left(0,-1/q\right)$ in \eqref{C15}.
\begin{comment}
\begin{align*}
&\sum_{n=0}^{\infty}
\frac{\left(-q;q^2\right)_{n+1}\left(q^4;q^4\right)_n}
{\left(q^2;q^2\right)_{n+1}\left(q^2;q^2\right)_n\left(q^2;q^4\right)_{n+1}}q^{n^2+2n}
\\
&\quad=\frac{1}{\left(1-q\right)\left(1-q^2\right)}
\sum_{n=0}^{\infty}\frac{\left(q^4;q^4\right)_n}
{\left(q^4;q^2\right)_n\left(q^2;q\right)_{2n}}q^{n^2+2n}\\
&\quad=\frac{\left(-q;q^2\right)_{\infty}\left(-q^2,-q^{14},q^{16};q^{16}\right)_{\infty}}
{\left(q^2;q^2\right)_{\infty}}.
\end{align*}
\end{comment}
\begin{comment}
We observe that
\begin{align}
\sum_{n=0}^{\infty}\frac{\left(-q^4;q^4\right)_{n-1}\left(-q;q^2\right)_nq^{n^2}}
{\left(q^2;q^4\right)_n\left(q^2;q^2\right)_n\left(-q^2;q^2\right)_{n-1}}
=\sum_{n=0}^{\infty}\frac{\left(\mathrm{i},-\mathrm{i};q^2\right)_n}
{\left(-1,q,q^2;q^2\right)_n}q^{n^2}.\label{equ7}
\end{align}
Replacing $q$ by $q^2$ and setting $\left(\alpha,\beta,a,b\right)\to\left(1,\mathrm{i},0,-q\right)$ in Theorem \ref{T4} and using the Jacobi triple product identity \eqref{Jacobi} with $\left(q,z\right)\to\left(q^{16},-q^6\right)$, we get
\begin{align*}
\frac{\left(q^2;q^2\right)_{\infty}}{\left(-q;q^2\right)_{\infty}}
\sum_{n=0}^{\infty}\frac{\left(\mathrm{i},-\mathrm{i};q^2\right)_n}{\left(-1,q,q^2;q^2\right)_n}q^{n^2}
=1+\sum_{n=1}^{\infty}\left(-1\right)^n\left(1+q^{4n}\right)q^{8n^2-2n}
=\left(q^6,q^{10},q^{16};q^{16}\right)_{\infty}.
\end{align*}
Combing the above identity with \eqref{equ7}, we obtain \eqref{S.71}.
\end{comment}
\begin{comment}
Taking $\left(\alpha,\beta,a,b\right)\to\left(q^2,q\mathrm{i},0,-1/q\right)$ into Theorem \ref{T4} and using the Jacobi triple product identity \eqref{Jacobi} with $\left(q,z\right)\to\left(q^8,-q^2\right)$ lead to
\begin{align*}
&\frac{\left(q^2;q\right)_{\infty}}{\left(-q;q\right)_{\infty}}
\sum_{n=0}^{\infty}\frac{\left(-q^2,q\mathrm{i},-q\mathrm{i};q\right)_n}
{\left(q,-q^2,q^{3/2},-q^{3/2};q\right)_n}q^{\frac{n^2+n}{2}}\nonumber\\
&\quad=
\sum_{n=0}^{\infty}\left(-1\right)^n\left(1-q^{4n+2}\right)
q^{4n^2+2n}\nonumber\\
&\quad=\left(q^2,q^6,q^8;q^8\right)_{\infty}.
\end{align*}
Therefore, we can complete the proof of \eqref{S.102} as follows,
\begin{align*}
&\sum_{n=0}^{\infty}\frac{\left(-q;q\right)_{n+1}\left(-q^2;q^2\right)_n
}{\left(q;q\right)_{2n+2}}q^{\frac{n^2+n}{2}}\nonumber\\
&\quad=\frac{1}{1-q^2}\sum_{n=0}^{\infty}\frac{\left(-q^2,q\mathrm{i},-q\mathrm{i};q\right)_n}
{\left(q,-q^2,q^{3/2},-q^{3/2};q\right)_n}q^{\frac{n^2+n}{2}}\nonumber\\
&\quad=
\frac{\left(-q;q\right)_{\infty}\left(q^2,q^6,q^8;q^8\right)_{\infty}}
{\left(1+q\right)\left(q;q\right)_{\infty}}.
\end{align*}

Taking $\left(\alpha,\beta,a,b\right)\to\left(1,\mathrm{i},0,-q\right)$ into Theorem \ref{T4} brings
\begin{align}
&\frac{\left(q;q\right)_{\infty}}{\left(-q;q\right)_{\infty}}\sum_{n=0}^{\infty}
\frac{\left(-1,\mathrm{i},-\mathrm{i};q\right)_n}{\left(-1,q,q^{1/2},-q^{1/2};q\right)_n}
q^{\frac{n\left(n+1\right)}{2}}\nonumber\\
&\quad=1+2\sum_{n=1}^{\infty}\left(-1\right)^nq^{4n^2}\nonumber\\
&\quad=\left(q^4,q^4,q^8;q^8\right)_{\infty}.\label{e10}
\end{align}
Then, we deduce that
\begin{align*}
&\frac{\left(q;q\right)_{\infty}}{\left(-q;q\right)_{\infty}}
\sum_{n=0}^{\infty}\frac{
\left(-q;q\right)_{n}\left(-q^2;q^2\right)_{n-1}
}{\left(q;q\right)_{2n}}q^{\frac{n\left(n+1\right)}{2}}\\
&\quad=\frac{\left(q;q\right)_{\infty}}{\left(-q;q\right)_{\infty}}
\left(\frac{1}{2}+\frac{1}{2}\sum_{n=0}^{\infty}
\frac{\left(-1,\mathrm{i},-\mathrm{i};q\right)_n}{\left(-1,q,q^{1/2},-q^{1/2};q\right)_n}
q^{\frac{n\left(n+1\right)}{2}}\right)\\
&\quad=\frac{\left(q;q\right)_{\infty}}
{2\left(-q;q\right)_{\infty}}+\frac{1}{2}\left(q^4,q^4,q^8;q^8\right)_{\infty}\\
&\quad=\left(-q^{16},-q^{16},q^{32};q^{32}\right)_{\infty}
-q\left(-q^{8},-q^{24},q^{32};q^{32}\right)_{\infty},
\end{align*}
where we have utilized
\end{comment}
\end{proof}
\section{Identities from two ${_5\phi_4}$ summations of Andrews}\label{sec7}
In section, we utilize Andrew's  two ${_5\phi_4}$ summation formulas to establish some $q$-series transformations.
%\subsection{Proofs of Slater's (73), (75), (90), (93), (114) and (116)}
\begin{comment}
\begin{theorem}\label{Thm5}
For $\max\left\{|ab|, |aq^6|, |bq^6|\right\}<1$, we have
\begin{align*}
&\frac{\left(q^{12},ab;q^6\right)_{\infty}}{\left(aq^6,bq^6;q^6\right)_{\infty}}
\sum_{n=0}^{\infty}
\frac{\left(q^6/a,q^6/b,q^4,q^8;q^6\right)_n}
{\left(q^6,-q^6,q^9,-q^9;q^6\right)_n}\left(ab\right)^n\\
&\quad=\sum_{n=0}^{\infty}\left(-1\right)^n\frac{1-q^{4n+2}}{1-q^2}
\frac{\left(q^6/a,q^6/b;q^6\right)_n}{\left(aq^6,bq^6;q^6\right)_n}a^nb^nq^{3n^2+n}.
\end{align*}
\end{theorem}
\begin{proof}
We begin by recalling the following strange $q$-series identity due to Andrews
\cite[Eq. (4.5)]{Andrews}:
\begin{align}
&{_5\phi_4}\left(\begin{array}{ccccc}
q^{-n},\,\alpha q^n,\,\alpha^{1/3}q^{1/3},\,\alpha^{1/3}q^{2/3},\,\alpha^{1/3}q\\
{\alpha}^{1/2}q,\,-\alpha^{1/2}q,\,\alpha^{1/2}q^{1/2},\,-\alpha^{1/2}q^{1/2}
\end{array};q,\,q
\right)\nonumber\\
&\quad=\frac{\left(1-\alpha\right)\left(1-\alpha^{1/3}q^{2n/3}\right)
\left(q;q\right)_n\left(\alpha^{1/3};q^{1/3}\right)_n}
{\left(1-\alpha^{1/3}\right)\left(1-\alpha q^{2n}\right)\left(\alpha;q\right)_n\left(q^{1/3};q^{1/3}\right)_n}\left(\alpha q\right)^{n/3}.
\label{Andrews2}
\end{align}
Then, we substitute $\left(q,\alpha, \beta,\gamma,c,d,h\right)
\to\left(q^6,q^6,q^4,q^8,q^9,-q^9,-q^6\right)$ and $\lambda=e$ into \eqref{Formula4} and employing \eqref{Andrews2} with $\left(q,\alpha\right)\to\left(q^6,q^6\right)$ to obtain the desired result.
\end{proof}
\end{comment}
\begin{theorem}\label{theorem2}
For $|\alpha ab/q|<1$, we have
\begin{align*}
&\frac{\left(\alpha q,\alpha ab/q;q\right)_{\infty}}{\left(\alpha a,\alpha b;q\right)_{\infty}}
\sum_{n=0}^{\infty}
\frac{\left(q/a,q/b;q\right)_n\left(\alpha^{1/3} q^{1/3};q^{1/3}\right)_{3n}}{\left(q;q\right)_n\left(\alpha q;q\right)_{2n}}\left(\alpha ab/q\right)^n
\\
&\quad=
\sum_{n=0}^{\infty}\left(-1\right)^n
\frac{1-\alpha^{\frac{1}{3}}q^{\frac{2n}{3}}}{1-\alpha^{\frac{1}{3}}}
\frac{\left(q/a,q/b;q\right)_n\left(\alpha^{1/3};q^{1/3}\right)_n}
{\left(\alpha a,\alpha b;q\right)_n\left(q^{1/3};q^{1/3}\right)_n}\alpha^{\frac{4n}{3}}a^nb^nq^{\frac{3n^2-7n}{6}}.
\end{align*}
\end{theorem}
\begin{proof}
We begin by recalling the following strange $q$-series identity due to Andrews
\cite[Eq. (4.5)]{Andrews}:
\begin{align}
&{_5\phi_4}\left(\begin{array}{ccccc}
q^{-n},\,\alpha q^n,\,\alpha^{1/3}q^{1/3},\,\alpha^{1/3}q^{2/3},\,\alpha^{1/3}q\\
{\alpha}^{1/2}q,\,-\alpha^{1/2}q,\,\alpha^{1/2}q^{1/2},\,-\alpha^{1/2}q^{1/2}
\end{array};q,\,q
\right)\nonumber\\
&\quad=\frac{\left(1-\alpha\right)\left(1-\alpha^{1/3}q^{2n/3}\right)
\left(q;q\right)_n\left(\alpha^{1/3};q^{1/3}\right)_n}
{\left(1-\alpha^{1/3}\right)\left(1-\alpha q^{2n}\right)\left(\alpha;q\right)_n\left(q^{1/3};q^{1/3}\right)_n}\left(\alpha q\right)^{n/3}.
\label{Andrews2}
\end{align}
Then, we substitute $\beta=\alpha^{1/3}q^{1/3}, \beta=\alpha^{1/3}q^{2/3}, \lambda=\alpha^{1/3}q, c=-d=\alpha^{1/2}q$ and 
$e=-h=\alpha^{1/2}q^{1/2}$ into \eqref{Formula4} and employ \eqref{Andrews2} to obtain the desired result.
\end{proof}
\begin{corollary}
We have
\begin{align}
&\sum_{n=0}^{\infty}\frac{\left(q;q\right)_{3n+1}}
{\left(q;q^3\right)_n\left(q^3;q^3\right)_{2n+1}}q^{3n^2+3n}
=\frac{\left(q,q^8,q^9;q^9\right)_{\infty}}{\left(q^3;q^3\right)_{\infty}}
,\tag{S.40}\label{S.40}\\
&\sum_{n=0}^{\infty}\frac{\left(q;q\right)_{3n+1}\left(-q^3;q^3\right)_n
}
{\left(q^3;q^3\right)_n\left(q^3;q^3\right)_{2n+1}}
q^{\frac{{3n^2+3n}}{2}}
=\frac{\left(-q^3;q^3\right)_{\infty}
\left(q,q^5,q^6;q^6\right)_{\infty}}{\left(q^3;q^3\right)_{\infty}}.\tag{\cite[Eq. (7.3)]{Bailey}}\label{new5}
\end{align}
\end{corollary}
\begin{proof}
Taking replacing  $q$ by $q^6$ and  fixing $\alpha=q^6$ in Theorem \ref{theorem2}, we deduce that,
for $\max\left\{|ab|, |aq^6|, |bq^6|\right\}<1$, 
\begin{align}
&\frac{\left(q^{12},ab;q^6\right)_{\infty}}{\left(aq^6,bq^6;q^6\right)_{\infty}}
\sum_{n=0}^{\infty}
\frac{\left(q^6/a,q^6/b,q^4,q^8;q^6\right)_n}
{\left(q^6,-q^6,q^9,-q^9;q^6\right)_n}\left(ab\right)^n\nonumber\\
&\quad=\sum_{n=0}^{\infty}\left(-1\right)^n\frac{1-q^{4n+2}}{1-q^2}
\frac{\left(q^6/a,q^6/b;q^6\right)_n}{\left(aq^6,bq^6;q^6\right)_n}a^nb^nq^{3n^2+n}.\label{C18}
\end{align}
Taking $\left(a,b\right)\to\left(0,0\right)$ into \eqref{C18}, we obtain
\begin{align}
&\left(q^{12};q^6\right)_{\infty}
\sum_{n=0}^{\infty}
\frac{\left(q^4,q^8;q^6\right)_n}{\left(q^{12},q^{18};q^{12}\right)_n}
q^{6n^2+6n}
%\nonumber\\
%&\quad=\frac{1}{1-q^2}\sum_{n=0}^{\infty}\left(-1\right)^n\left(1-q^{4n+2}\right)
%q^{9n^2+7n}\nonumber\\
%&\quad=\frac{1}{1-q^2}\sum_{n=-\infty}^{\infty}\left(-1\right)^nq^{9n^2+7n}\nonumber\\
%&\quad
=\frac{\left(q^2,q^{16},q^{18};q^{18}\right)_{\infty}}{1-q^2}.\label{B2}
\end{align}
Then, replacing $q^2$ by $q$, we get
\begin{align*}
\sum_{n=0}^{\infty}\frac{\left(q;q^3\right)_{n+1}\left(q^2;q^3\right)_n}
{\left(q^6;q^3\right)_{2n}}q^{3n^2+3n}
=\frac{\left(q,q^8,q^9;q^9\right)_{\infty}}{
\left(q^6;q^3\right)_{\infty}},
\end{align*}
which means \eqref{S.40} is true.

Moreover, inserting $\left(a,b\right)\to\left(0,-1\right)$ into \eqref{C18}, we obtain
\begin{align}
&\frac{\left(q^{12};q^6\right)_{\infty}}{\left(-q^6;q^6\right)_{\infty}}
\sum_{n=0}^{\infty}
\frac{\left(q^4,q^8;q^6\right)_n}
{\left(q^6;q^6\right)_n\left(q^{18};q^{12}\right)_n}q^{3n^2+3n}
%\nonumber\\&\quad=\frac{1}{1-q^2}\sum_{n=0}^{\infty}\left(-1\right)^n\left(1-q^{4n+2}\right)q^{6n^2+4n}\nonumber\\
%&\quad
=\frac{\left(q^2,q^{10},q^{12};q^{12}\right)_{\infty}}{1-q^2}.\label{B3}
\end{align}
Then, letting $q^2\to q$ in \eqref{B3}, we have
\begin{align*}
\sum_{n=0}^{\infty}\frac{\left(q;q^3\right)_{n+1}\left(q^2;q^3\right)_n}
{\left(q^3;q^3\right)_n\left(q^3;q^6\right)_{n+1}}q^{\frac{3n^2+3n}{2}}
=
\frac{\left(-q^3;q^3\right)_n\left(q,q^{5},q^{6};q^{6}\right)_{\infty}}
{\left(q^3;q^3\right)_{\infty}},
\end{align*}
which means \eqref{new5} is true.
\end{proof}

\begin{corollary}
We have
\begin{align}
\sum_{n=0}^{\infty}\frac{\left(q;q\right)_{3n}}{\left(q^3;q^3\right)_n\left(q^3;q^3\right)_{2n}}q^{3n^2}
&=\frac{\left(q^4,q^5,q^9;q^9\right)_{\infty}}{\left(q^3;q^3\right)_{\infty}},\tag{S.42}\label{S.42}\\
\sum_{n=0}^{\infty}\frac{\left(q;q\right)_{3n}\left(-q^3;q^3\right)_n}{\left(q^3;q^3\right)_n\left(q^3;q^3\right)_{2n}}q^{\frac{3n^2-3n}{2}}
&=\frac{\left(-q^3;q^3\right)_{\infty}}{\left(q^3;q^3\right)_{\infty}}
\left(\left(q,q^5,q^6;q^6\right)_{\infty}+\left(q^2,q^4,q^6;q^6\right)_{\infty}\right),\label{N12}\\
\sum_{n=0}^{\infty}\frac{\left(q^2;q^2\right)_{3n}}{\left(q^3;q^3\right)_{2n}\left(q^{12};q^{12}\right)_{n}}q^{3n^2}
&=\frac{\left(-q^3;q^6\right)_{\infty}\left(q^5,q^7,q^{12};q^{12}\right)_{\infty}}{\left(q^6;q^6\right)_{\infty}}
.\tag{\rm{\cite[Eq. (7.5)]{Bailey}}}\label{D1}
\end{align}
\end{corollary}
\begin{proof}
Replacing $q$ by $q^3$ and fixing $\alpha=1$ in Theorem \ref{theorem2}, we have
\begin{align}
&\frac{\left(q^3,ab/q^3;q^3\right)_{\infty}}{\left(a,b;q^3\right)_{\infty}}
\sum_{n=0}^{\infty}\frac{\left(q^3/a,q^3/b;q^3\right)_n\left(q;q\right)_{3n}}
{\left(q^3;q^3\right)_n\left(q^3;q^3\right)_{2n}}\left(ab/q^3\right)^n\nonumber\\
&\quad=
1+\sum_{n=0}^{\infty}\left(-1\right)^n\left(1+q^n\right)\frac{\left(q^3/a,q^3/b;q^3\right)_n}{\left(a,b;q^3\right)_n}
a^nb^nq^{\frac{3n^2-7n}{2}}.\label{C19}
\end{align}

Taking $\left(a,b\right)\to\left(0,0\right)$ and $\left(a,b\right)\to\left(0,-1\right)$  into \eqref{C19}, respectively, we arrive at \eqref{S.42} and \eqref{N12}. Moreover,
Replacing $q$ by $q^2$ and taking $\left(a,b\right)\to\left(0,-q^3\right)$ into \eqref{C19}, we arrive at \ref{D1}.
\end{proof}
\begin{theorem}\label{thm4}
For $\max\left\{|a|, |b|, |ab/q|\right\}<1$ and  $\omega=\mathrm{e^{2\pi\mathrm{i}/3}}$, we have
\begin{align*}
&\frac{\left(q,ab/q;q\right)_{\infty}}{\left(a,b;q\right)_{\infty}}
\sum_{n=0}^{\infty}
\frac{\left(q/a,q/b,\omega,\omega^2;q\right)_n}
{\left(q,-1,q^{1/2},-q^{1/2};q\right)_n}\left(ab/q\right)^n\\
&\quad=1+\sum_{n=1}^{\infty}\left(-1\right)^n\left(1+q^{3n}\right)
\frac{\left(q/a,q/b;q\right)_{3n}}
{\left(a,b;q\right)_{3n}}
a^{3n}b^{3n}q^{\frac{9n^2-9n}{2}}\\
&\qquad+\frac{q^2}{2}\sum_{n=1}^{\infty}\left(-1\right)^n\left(1+q^{3n-1}\right)
\frac{\left(q/a,q/b;q\right)_{3n-1}}
{\left(a,b;q\right)_{3n-1}}
a^{3n-1}b^{3n-1}q^{\frac{9n^2-15n}{2}}\\
&\qquad+\frac{1}{2q}\sum_{n=0}^{\infty}\left(-1\right)^n\left(1+q^{3n+1}\right)
\frac{\left(q/a,q/b;q\right)_{3n+1}}
{\left(a,b;q\right)_{3n+1}}
a^{3n+1}b^{3n+1}q^{\frac{9n^2-3n}{2}}
.
\end{align*}
\end{theorem}
\begin{proof}
Recall a $_5\phi_4$ summation formula first proved by Andrews \cite{Andrews}:
\begin{align}
&{_5\phi_4}\left(\begin{array}{ccccc}
q^{-n},\,\alpha q^n,\,\alpha^{1/3},\,\alpha^{1/3}\omega,\,\alpha^{1/3}\omega^2\\
\sqrt{\alpha},\,-\sqrt{\alpha},\,\sqrt{q\alpha},\,-\sqrt{q\alpha}
\end{array};q,\,q
\right)\nonumber\\
&\quad=\left\{
\begin{array}{ll}
0,&\quad n\equiv1,2\pmod3,\\
\frac{\left(q;q\right)_n\left(\alpha;q^3\right)_{n/3}}
{\left(\alpha;q\right)_{n}\left(q^3;q^3\right)_{n/3}}\alpha^{n/3},&\quad n\equiv0\pmod3.
\end{array}
\right.\label{Andrews1}
\end{align}
Letting $\alpha=1$ in \eqref{Andrews1}, we deduce that
\begin{align}
{_4\phi_3}\left(
\begin{array}{cccc}
q^{-n},\,q^n,\,\omega,\,\omega^2\\
-1,\,q^{1/2},\,-q^{1/2}
\end{array};q,\,q\right)
=\left\{
\begin{array}{ll}
-1/2,&\quad n\equiv1,\,2\pmod3,\\
1,&\quad n\equiv0\pmod3.
\end{array}\right.\label{A3}
\end{align}
Then, substituting $\left(\alpha, \beta,\gamma,c,d,h\right)
\to\left(1,\omega,\omega^2,q^{1/2},-q^{1/2},-1\right)$ and $\lambda=e$ into \eqref{Formula4}, we obtain the desired result.
\end{proof}
\begin{corollary}
We have
\begin{align}
&\sum_{n=0}^{\infty}\frac{\left(-q;q\right)_n\left(q^3;q^3\right)_{n-1}
}
{\left(q;q\right)_{n}\left(q;q\right)_{2n-1}}q^{\frac{n(n-1)}{2}}\nonumber\\
&\quad=\frac{\left(-q;q\right)_{\infty}}{\left(q;q\right)_{\infty}}
\left(
\left(q^{6},q^{12},q^{18};q^{18}\right)_{\infty}+
\left(q^{9},q^{9},q^{18};q^{18}\right)_{\infty}\right)
,\tag{S.73}\label{S.73}\\
&\sum_{n=0}^{\infty}\frac{\left(-q;q\right)_n\left(q^3;q^3\right)_{n-1}
}
{\left(q;q\right)_{n-1}\left(q;q\right)_{2n}}q^{\frac{n(n+1)}{2}}\nonumber\\
%&\quad=\left(q^9,q^9,q^{18};q^{18}\right)_{\infty}
%+q\left(q^3,q^{15},q^{18};q^{18}\right)_{\infty},\tag{\rm{\cite[Eq. (1.27)]{McLaughlin_2008}}}\label{C6}\\
&\quad=\frac{\left(-q;q\right)_{\infty}}{\left(q;q\right)_{\infty}}
\left(
\left(q^{9},q^{9},q^{18};q^{18}\right)_{\infty}-q
\left(q^{3},q^{15},q^{18};q^{18}\right)_{\infty}\right)
,\tag{S.75}\label{S.75}\\
&1+\sum_{n=1}^{\infty}\frac{\left(-1;q\right)_{n+1}\left(q^3;q^3\right)_{n-1}
}
{\left(q;q\right)_{n-1}\left(q;q\right)_{2n}}q^{\frac{n(n+1)}{2}}
=\frac{\left(-q;q\right)_{\infty}\left(q^{9},q^{9},q^{18};q^{18}\right)_{\infty}}{\left(q;q\right)_{\infty}}
,\tag{S.78}\label{S.78}\\
&\sum_{n=0}^{\infty}\frac{\left(q^3;q^3\right)_{n-1}}
{\left(q;q\right)_n\left(q;q\right)_{2n-1}}q^{n^2}
=\frac{\left(q^{12},q^{15},q^{27};q^{27}\right)_{\infty}}{\left(q;q\right)_{\infty}}
,\tag{S.93}\label{S.93}\\
&\sum_{n=0}^{\infty}\frac{\left(-q;q^2\right)_{n}
\left(q^6;q^6\right)_{n-1}}
{\left(q^2;q^2\right)_{2n-1}\left(q^2;q^2\right)_n}q^{n^2}
=\frac{\left(-q;q^2\right)_{\infty}\left(q^{15},q^{21},q^{36};q^{36}\right)_{\infty}}
{\left(q^2;q^2\right)_{\infty}}.\tag{S.114} \label{S.114}
\end{align}
\end{corollary}
\begin{proof}
We observe that
\begin{align}
&\sum_{n=0}^{\infty}\frac{\left(-q;q\right)_n\left(q^3;q^3\right)_{n-1}
}
{\left(q;q\right)_{n}\left(q;q\right)_{2n-1}}q^{\frac{n(n-1)}{2}}\nonumber\\
&\quad=1+\sum_{n=1}^{\infty}\frac{\left(-q;q\right)_n
\left(q\omega,q\omega^2;q\right)_{n-1}}
{\left(q,q^{1/2},-q^{1/2};q\right)_n\left(-q;q\right)_{n-1}}
q^{\frac{n\left(n-1\right)}{2}}\nonumber\\
&\quad=1+\frac{2}{3}\sum_{n=1}^{\infty}\frac{\left(-q,\omega,\omega^2;q\right)_{n}}
{\left(q,-1,q^{1/2},-q^{1/2};q\right)_n}
q^{\frac{n\left(n-1\right)}{2}}\nonumber\\
&\quad=\frac{1}{3}+\frac{2}{3}\sum_{n=0}^{\infty}\frac{\left(-q,\omega,\omega^2;q\right)_{n}}
{\left(q,-1,q^{1/2},-q^{1/2};q\right)_n}
q^{\frac{n\left(n-1\right)}{2}}.\label{A6}
\end{align}
Meanwhile, taking $\left(a,b\right)\to\left(0,-1\right)$ into Theorem \ref{thm4} results in
\begin{align}
&\frac{\left(q;q\right)_{\infty}}{\left(-1;q\right)_{\infty}}
\sum_{n=0}^{\infty}\frac{\left(-q,\omega,\omega^2;q\right)_n}
{\left(q,-1,q^{1/2},-q^{1/2};q\right)_n}q^{\frac{n\left(n-1\right)}{2}}\nonumber\\
&\quad=1+\frac{1}{2}\sum_{n=1}^{\infty}\left(-1\right)^n\left(1+q^{3n}\right)^2
q^{9n^2-3n}
+\frac{q^2}{4}\sum_{n=1}^{\infty}\left(-1\right)^n\left(1+q^{3n-1}\right)^2
q^{9n^2-9n}\nonumber\\
&\qquad+\frac{1}{4}\sum_{n=0}^{\infty}\left(-1\right)^n\left(1+q^{3n+1}\right)^2
q^{9n^2+3n}\nonumber\\
&\quad=\frac{1}{2}\sum_{n=-\infty}^{\infty}\left(-1\right)^nq^{9n^2}
+\frac{q}{2}\sum_{n=-\infty}^{\infty}\left(-1\right)^nq^{9n^2-6n}
+\frac{3}{4}\sum_{n=-\infty}^{\infty}\left(-1\right)^nq^{9n^2-3n}.\label{A7}
\end{align}
Then, by inserting \eqref{A7} into \eqref{A6}, we deduce that
\begin{align}
&\frac{\left(q;q\right)_{\infty}}{\left(-1;q\right)_{\infty}}
\sum_{n=0}^{\infty}\frac{\left(-q;q\right)_n\left(q^3;q^3\right)_{n-1}
}
{\left(q;q\right)_{n}\left(q;q\right)_{2n-1}}q^{\frac{n(n-1)}{2}}\nonumber\\
&\quad=\frac{1}{6}\frac{\left(q;q\right)_{\infty}}{\left(-1;q\right)_{\infty}}
+\frac{1}{3}\left(q^9,q^9,q^{18};q^{18}\right)_{\infty}
+\frac{q}{3}\left(q^3,q^{15},q^{18};q^{18}\right)_{\infty}
+\frac{1}{2}\left(q^6,q^{12},q^{18};q^{18}\right)_{\infty}
\nonumber\\
&\quad=\frac{1}{2}\left(q^9,q^9,q^{18};q^{18}\right)_{\infty}
+\frac{1}{2}\left(q^6,q^{12},q^{18};q^{18}\right)_{\infty},
\end{align}
where we have utilized the following fact:
\begin{align}
\frac{\left(q;q\right)_{\infty}}{\left(-q;q\right)_{\infty}}
%&=\left(q,q,q^2;q^2\right)_{\infty}\nonumber\\
&=\sum_{n=-\infty}^{\infty}\left(-1\right)^nq^{n^2}\nonumber\\
&=\sum_{n=-\infty}^{\infty}\left(-1\right)^nq^{9n^2}
-\sum_{n=-\infty}^{\infty}\left(-1\right)^nq^{\left(3n+1\right)^2}
-\sum_{n=-\infty}^{\infty}\left(-1\right)^nq^{\left(3n-1\right)^2}\nonumber\\
&=\left(q^9,q^9,q^{18};q^{18}\right)_{\infty}
-2q\left(q^3,q^{15},q^{18};q^{18}\right)_{\infty}.\label{A8}
\end{align}
Hence, \eqref{S.73} is true.

Next,  we take $\left(a,b\right)\to\left(0,-q\right)$ into Theorem \ref{thm4} to find that 
\begin{align}
&\frac{\left(q;q\right)_{\infty}}{\left(-q;q\right)_{\infty}}
\sum_{n=0}^{\infty}\frac{\left(\omega,\omega^2;q\right)_n}
{\left(q,q^{1/2},-q^{1/2};q\right)_n}
q^{\frac{n\left(n+1\right)}{2}}\nonumber\\
&\quad=1+2\sum_{n=1}^{\infty}\left(-1\right)^nq^{9n^2}
+\sum_{n=-\infty}^{\infty}\left(-1\right)^nq^{\left(3n+1\right)^2}\nonumber\\
&\quad=\left(q^9,q^9,q^{18};q^{18}\right)_{\infty}
+q\left(q^3,q^{15},q^{18};q^{18}\right)_{\infty}.\label{A4}
\end{align}
Also, we observe that
\begin{align}
&\frac{\left(q;q\right)_{\infty}}{\left(-q;q\right)_{\infty}}\sum_{n=0}^{\infty}\frac{\left(-q;q\right)_n\left(q^3;q^3\right)_{n-1}
}
{\left(q;q\right)_{n-1}\left(q;q\right)_{2n}}q^{\frac{n(n+1)}{2}}\nonumber\\
&\quad=\frac{\left(q;q\right)_{\infty}}{\left(-q;q\right)_{\infty}}+\frac{\left(q;q\right)_{\infty}}{\left(-q;q\right)_{\infty}}\sum_{n=1}^{\infty}\frac{\left(q\omega,q\omega^2;q\right)_{n-1}}
{\left(q,q^{1/2},-q^{1/2};q\right)_n}q^{\frac{n\left(n+1\right)}{2}}\nonumber\\
&\quad=\frac{\left(q;q\right)_{\infty}}{\left(-q;q\right)_{\infty}}+\frac{\left(q;q\right)_{\infty}}{3\left(-q;q\right)_{\infty}}\sum_{n=1}^{\infty}\frac{\left(\omega,\omega^2;q\right)_{n}}
{\left(q,q^{1/2},-q^{1/2};q\right)_n}q^{\frac{n\left(n+1\right)}{2}}\nonumber\\
&\quad=\frac{2\left(q;q\right)_{\infty}}{3\left(-q;q\right)_{\infty}}+\frac{\left(q;q\right)_{\infty}}{3\left(-q;q\right)_{\infty}}\sum_{n=0}^{\infty}
\frac{\left(\omega,\omega^2;q\right)_{n}}
{\left(q,q^{1/2},-q^{1/2};q\right)_n}q^{\frac{n\left(n+1\right)}{2}}\nonumber\\
&\quad=\frac{2}{3}\frac{\left(q;q\right)_{\infty}}{\left(-q;q\right)_{\infty}}
+\frac{1}{3}\left(\left(q^9,q^9,q^{18};q^{18}\right)_{\infty}
+q\left(q^3,q^{15},q^{18};q^{18}\right)_{\infty}
\right)\nonumber\\
&\quad=\left(q^9,q^9,q^{18};q^{18}\right)_{\infty}
-q\left(q^3,q^{15},q^{18};q^{18}\right)_{\infty},\nonumber
\end{align}
where in the last step we have used \eqref{A8}.
Therefore, we complete the proof of \eqref{S.75}.

Additionally, we can prove \eqref{S.78} as follows,
\begin{align*}
&1+\sum_{n=1}^{\infty}\frac{\left(q^3;q^3\right)_{n-1}\left(-1;q\right)_{n+1}}
{\left(q;q\right)_{2n}\left(-q;q\right)_{n-1}}q^{\frac{n(n+1)}{2}}\\
&\quad=
\frac{1}{3}+\frac{2}{3}\sum_{n=0}^{\infty}\frac{\left(\omega,\omega^2;q\right)_n}{\left(q,q^{1/2},-q^{1/2};q\right)_n}q^{\frac{n(n+1)}{2}}
\\
&\quad=\frac{\left(-q;q\right)_{\infty}\left(q^9,q^9,q^{18};q^{18}\right)_{\infty}}{\left(q;q\right)_{\infty}}.
\end{align*}



Taking $\left(a,b\right)\to\left(0,0\right)$ into Theorem \ref{thm4}, we immediately obtain
\begin{align}
&\left(q;q\right)_{\infty}\sum_{n=0}^{\infty}
\frac{\left(\omega,\omega^2;q\right)_n}
{\left(q,-1,q^{1/2},-q^{1/2};q\right)_n}q^{n^2}\nonumber\\
&\quad=1+\sum_{n=1}^{\infty}\left(-1\right)^n\left(1+q^{3n}\right)
q^{\frac{27}{2}n^2-\frac{3}{2}n}
+\frac{q}{2}\sum_{n=0}^{\infty}\left(-1\right)^n\left(1+q^{3n+1}\right)
q^{\frac{27}{2}n^2+\frac{15}{2}n}\nonumber\\
&\qquad+\frac{q^2}{2}\sum_{n=1}^{\infty}\left(-1\right)^n\left(1+q^{3n-1}\right)
q^{\frac{27}{2}n^2-\frac{21}{2}n}\nonumber\\
&\quad=\left(q^{12},q^{15},q^{27};q^{27}\right)_{\infty}
+\frac{q}{2}\left(q^{6},q^{21},q^{27};q^{27}\right)_{\infty}
+\frac{q^2}{2}\left(q^{3},q^{24},q^{27};q^{27}\right)_{\infty}.\label{A9}
\end{align}
Therefore, we can prove \eqref{S.93} as follows,
\begin{align}
&\left(q;q\right)_{\infty}\sum_{n=0}^{\infty}
\frac{\left(q^3;q^3\right)_{n-1}}
{\left(q;q\right)_n\left(q;q\right)_{2n-1}}q^{n^2}\nonumber\\
&\quad=\left(q;q\right)_{\infty}\left(1+\sum_{n=1}^{\infty}
\frac{\left(q,q\omega,q\omega^2;q\right)_{n-1}}
{\left(q;q\right)_n\left(q^{1/2},-q^{1/2};q\right)_{n}
\left(-q;q\right)_{n-1}}q^{n^2}\right)\nonumber\\
&\quad=\left(q;q\right)_{\infty}\left(1+\frac{2}{3}\sum_{n=1}^{\infty}
\frac{\left(\omega,\omega^2;q\right)_{n}}
{\left(q;q\right)_n\left(q^{1/2},-q^{1/2};q\right)_{n}
\left(-1;q\right)_{n}}q^{n^2}\right)\nonumber\\
&\quad=\left(q;q\right)_{\infty}\left(\frac{1}{3}+\frac{2}{3}\sum_{n=0}^{\infty}
\frac{\left(\omega,\omega^2;q\right)_{n}}
{\left(q;q\right)_n\left(q^{1/2},-q^{1/2};q\right)_{n}
\left(-1;q\right)_{n}}q^{n^2}\right)\nonumber\\
&\quad=\left(q^{12},q^{15},q^{27};q^{27}\right)_{\infty},
\end{align}
where we have used
\begin{align}
\left(q;q\right)_{\infty}%&=\left(q,q^2,q^3;q^3\right)_{\infty}\nonumber\\
&=\sum_{n=-\infty}^{\infty}\left(-1\right)^nq^{\frac{3n^2-n}{2}}\nonumber\\
&=\sum_{n=-\infty}^{\infty}\left(-1\right)^nq^{\frac{27n^2-3n}{2}}
-q\sum_{n=-\infty}^{\infty}\left(-1\right)^nq^{\frac{27n^2-15n}{2}}
-q^2\sum_{n=-\infty}^{\infty}\left(-1\right)^nq^{\frac{27n^2-21n}{2}}\nonumber\\
&=\left(q^{12},q^{15},q^{27};q^{27}\right)_{\infty}
-q\left(q^{6},q^{21},q^{27};q^{27}\right)_{\infty}
-q^2\left(q^{3},q^{24},q^{27};q^{27}\right)_{\infty}.\label{A10}
\end{align}

Replacing $q$ by $q^2$ and taking $\left(a,b\right)\to\left(0,-q\right)$ into Theorem \ref{thm4}, we deduce that
\begin{align}
&\frac{\left(q^2;q^2\right)_{\infty}}{\left(-q;q^2\right)_{\infty}}
\sum_{n=0}^{\infty}\frac{\left(-q,\omega,\omega^2;q^2\right)_n}
{\left(q^2,-1,q,-q;q^2\right)_n}q^{n^2}
\nonumber\\
&\quad=1+\sum_{n=0}^{\infty}\left(-1\right)^n\left(1+q^{6n}\right)q^{18n^2-3n}
+\frac{q}{2}\sum_{n=0}^{\infty}\left(-1\right)^n\left(1+q^{6n+2}\right)q^{18n^2+9n}
\nonumber\\
&\qquad+\frac{q^3}{2}\sum_{n=0}^{\infty}\left(-1\right)^n
\left(1+q^{6n-2}\right)q^{18n^2-15n}\nonumber\\
&\quad=\left(q^{15},q^{21},q^{36};q^{36}\right)_{\infty}
+\frac{q}{2}\left(q^{9},q^{27},q^{36};q^{36}\right)_{\infty}
+\frac{q^3}{2}\left(q^{3},q^{33},q^{36};q^{36}\right)_{\infty}.\label{A11}
\end{align}
In terms of \eqref{A11}, we see that
\begin{align*}
&\frac{\left(q^2;q^2\right)_{\infty}}{\left(-q;q^2\right)_{\infty}}
\sum_{n=0}^{\infty}\frac{\left(-q;q^2\right)_n\left(q^6;q^6\right)_{n-1}}
{\left(q^2;q^2\right)_{2n-1}\left(q^2;q^2\right)_n}q^{n^2}\\
&\quad=
\frac{\left(q^2;q^2\right)_{\infty}}{\left(-q;q^2\right)_{\infty}}
\left(1+\sum_{n=1}^{\infty}\frac{\left(-q;q^2\right)_n
\left(q^2\omega,q^2\omega^2;q^6\right)_{n-1}}
{\left(q,-q,q^2;q^2\right)_{n}\left(-q^2;q^2\right)_{n-1}}q^{n^2}
\right)\\
&\quad=
\frac{\left(q^2;q^2\right)_{\infty}}{\left(-q;q^2\right)_{\infty}}
\left(1+\frac{2}{3}\sum_{n=1}^{\infty}\frac{\left(-q;q^2\right)_n
\left(\omega,\omega^2;q^6\right)_{n}}
{\left(-1,q,-q,q^2;q^2\right)_{n}}q^{n^2}
\right)\\
&\quad=
\frac{\left(q^2;q^2\right)_{\infty}}{\left(-q;q^2\right)_{\infty}}
\left(\frac{1}{3}+\frac{2}{3}\sum_{n=0}^{\infty}\frac{\left(-q;q^2\right)_n
\left(\omega,\omega^2;q^6\right)_{n}}
{\left(-1,q,-q,q^2;q^2\right)_{n}}q^{n^2}
\right)\\
&\quad=\frac{\left(q^2;q^2\right)_{\infty}}{3\left(-q;q^2\right)_{\infty}}
+\frac{2}{3}\left(q^{15},q^{21},q^{36};q^{36}\right)_{\infty}
+\frac{q}{3}\left(q^{9},q^{27},q^{36};q^{36}\right)_{\infty}
+\frac{q^3}{3}\left(q^{3},q^{33},q^{36};q^{36}\right)_{\infty}\\
&\quad=\left(q^{15},q^{21},q^{36};q^{36}\right)_{\infty},
\end{align*}
where we have used
\begin{align*}
\frac{\left(q^2;q^2\right)_{\infty}}{\left(-q;q^2\right)_{\infty}}
&=\sum_{n=-\infty}^{\infty}\left(-1\right)^nq^{2n^2-n}\\
&=\sum_{n=-\infty}^{\infty}\left(-1\right)^nq^{18n^2-3n}
-q\sum_{n=-\infty}^{\infty}\left(-1\right)^nq^{18n^2+9n}
-q^3\sum_{n=-\infty}^{\infty}\left(-1\right)^nq^{18n^2+15n}
\\
&=\left(q^{15},q^{21},q^{36};q^{36}\right)_{\infty}
-q\left(q^{9},q^{27},q^{36};q^{36}\right)_{\infty}
-q^3\left(q^{3},q^{33},q^{36};q^{36}\right)_{\infty}.
\end{align*}
Therefore, we complete the proof of \eqref{S.114}.
\end{proof}
\begin{theorem}\label{T1}
For $\max\left\{|aq^3|,|bq^3|,|abq^2|\right\}<1$, we have
\begin{align*}
&\frac{\left(q^3,abq^2;q\right)_{\infty}}{\left(aq^3,bq^3,q\right)_{\infty}}
\sum_{n=0}^{\infty}
\frac{\left(q/a,q/b;q\right)_n\left(q^3;q^3\right)_n}
{\left(q;q\right)_n\left(q^3;q\right)_{2n}}a^nb^nq^{2n}\\
&\quad=
\sum_{n=0}^{\infty}\left(-1\right)^n\left(1-q^{6n+3}\right)
\frac{\left(q/a,q/b;q\right)_{3n}}{\left(aq^3,bq^3;q\right)_{3n}}a^{3n}b^{3n}
q^{\frac{9n^2+15n}{2}}.
\end{align*}
\end{theorem}
\begin{proof}
Taking $\alpha=q^3$ into \eqref{Andrews1}, we get
\begin{align}
&{_5\phi_4}\left(\begin{array}{ccccc}
q^{-n},\,q^{n+3},\,q,\,q\omega,\,q\omega^2\\
q^{3/2},\,-q^{3/2},\,q^2,\,-q^2
\end{array};q,\,q
\right)\nonumber\\
&\quad=\left\{
\begin{array}{ll}
0,\,&n\equiv1,2\pmod3,\\
\frac{\left(1-q\right)\left(1-q^2\right)}
{\left(1-q^{3n+1}\right)\left(1-q^{3n+2}\right)}q^n,\, &n\equiv0\pmod3.
\end{array}
\right.\label{A12}
\end{align}
Setting $\left(\alpha,\beta,\gamma,\lambda,c,d,e,h\right)
=\left(q^3,q,q\omega,q\omega^2,q^{3/2},-q^{3/2},q^2,-q^2\right)$ in \eqref{Formula4} and using \eqref{A12}, we immediately obtain Theorem \ref{T1}.
\end{proof}
\begin{corollary}
We have
\begin{align}
\sum_{n=0}^{\infty}\frac{\left(-q;q\right)_{n+1}\left(q^3;q^3\right)_{n}}
{\left(q;q\right)_n\left(q;q\right)_{2n+2}}q^{\frac{n^2+3n}{2}}
&=\frac{\left(-q;q\right)_{\infty}\left(q^{3},q^{15},q^{18};q^{18}\right)_{\infty}}{\left(q;q\right)_{\infty}}
,\tag{S.76}\label{S.76}
\\
\sum_{n=0}^{\infty}\frac{\left(q^3;q^3\right)_{n}}
{\left(q;q\right)_n\left(q;q\right)_{2n+2}}q^{n^2+3n}
&=\frac{\left(q^{3},q^{24},q^{27};q^{27}\right)_{\infty}}{\left(q;q\right)_{\infty}}
,\tag{S.90}\label{S.90}
\\
\sum_{n=0}^{\infty}\frac{\left(-q;q^2\right)_{n+1}
\left(q^6;q^6\right)_n}
{\left(q^2;q^2\right)_{2n+2}\left(q^2;q^2\right)_n}q^{n^2+4n}
&=\frac{\left(-q;q^2\right)_{\infty}\left(q^3,q^{33},q^{36};q^{36}\right)_{\infty}}
{\left(q^2;q^2\right)_{\infty}}.\tag{S.116}\label{S.116}
\end{align}
\end{corollary}
\begin{proof}
Taking $\left(a,b\right)\to\left(0,-1/q\right)$ and $\left(a,b\right)\to\left(0,0\right)$, respectively,  into Theorem \ref{T1}, we obtain \eqref{S.76} and \eqref{S.90} after simplication.

Replacing $q$ by $q^2$ and taking $\left(a,b\right)\to\left(0,-1/q\right)$ into Theorem \ref{T1}, we get \eqref{S.116}.
\end{proof}
\begin{section}{Conclusion}
In this paper,  by reproving  $75$ identities of Rogers-Ramanujan type proposed by Slater and  uncovering several novel identities within this class, we show  that Liu's formula \eqref{Formula3} can imply many Rogers-Ramanujan type identities.  The readers who have interests can use  Liu's formula \eqref{Formula3} to discover more Rogers-Ramanujan type identities.

Below in Table $\ref{table1}$, we list the theorems and the corresponding  Rogers-Ramanujan type identities proposed by Slater.
\begin{table}[thbp]
\caption{Theorems and corresponding Rogers-Ramanujan type identities}\label{table1}
\centering
\begin{tabularx}{\textwidth}{lX}
  %\hline
  \toprule
  % after \\: \hline or \cline{col1-col2} \cline{col3-col4} ...
  {Theorem \ref{thm9}} & $\eqref{S.2}$, $\eqref{S.3}$, $\eqref{S.5}$, $\eqref{S.7}$, $\eqref{S.8}$, $\eqref{S.9}$, $\eqref{S.12}$, $\eqref{S.13}$, $\eqref{S.14}$, $\eqref{S.18}$, $\eqref{S.24}$, $\eqref{S.25}$, $\eqref{S.28}$, $\eqref{S.29}$, $\eqref{S.30}$, $\eqref{S.34}$, $\eqref{S.36}$, $\eqref{S.38}$, $\eqref{S.39}$, $\eqref{S.47}$, $\eqref{S.50}$, $\eqref{S.51}$, $\eqref{S.52}$, $\eqref{S.64}$, $\eqref{S.84}$, $\eqref{S.86}$ \\
  Theorem \ref{T7} & $\eqref{S.45}$, $\eqref{S.60}$ \\
  Theorem \ref{T9} & $\eqref{S.96}$ \\
  Theorem \ref{thm2} & $\eqref{S.11}$, $\eqref{S.22}$, $\eqref{S.58}$, $\eqref{S.72}$, $\eqref{S.87}$ \\
  Theorem \ref{thm11} & $\eqref{S.16}$, $\eqref{S.17}$, $\eqref{S.20}$, $\eqref{S.31}$, $\eqref{S.32}$, $\eqref{S.94}$, $\eqref{S.99}$ \\
  Theorem \ref{T5} & $\eqref{S.48}$, $\eqref{S.65}$, $\eqref{S.67}$, $\eqref{S.70}$, $\eqref{S.127}$ \\
  Theorem \ref{T3} & $\eqref{S.21}$, $\eqref{S.33}$ \\
  Theorem \ref{T6} & $\eqref{S.53}$, $\eqref{S.55}$, $\eqref{S.57}$ \\
  %Theorem \ref{T8} & $\eqref{S.58}$, $\eqref{S.72}$ \\
  %Theorem \ref{thm6} & $\eqref{S.61}$, $\eqref{S.98}$ \\
  Theorem \ref{T4} & $\eqref{S.35}$, $\eqref{S.43}$, $\eqref{S.49}$, $\eqref{S.54}$, $\eqref{S.59}$, $\eqref{S.61}$, $\eqref{S.68}$, $\eqref{S.69}$, $\eqref{S.71}$, $\eqref{S.98}$,  $\eqref{S.104}$ \\
  Theorem \ref{theorem2} & $\eqref{S.40}$, $\eqref{S.42}$\\
  Theorem \ref{thm4} & $\eqref{S.73}$, $\eqref{S.75}$, $\eqref{S.78}$, $\eqref{S.93}$, $\eqref{S.114}$\\
  Theorem \ref{T1} & $\eqref{S.76}$, $\eqref{S.90}$, $\eqref{S.116}$\\
  \bottomrule
\end{tabularx}
\end{table}
Astute readers may be 
aware that we have skipped  $55$ identities of Rogers-Ramanujan type proposed by Slater.  We omit the proof of $\left(\mathrm{S}.1\right)$ since it is a special case of  the Jacobi triple product identity \eqref{Jacobi}.  Challenges arise when we attempt to apply the preceding methodology to the remaining $54$ identities. The primary difficulties lie in evaluating certain terminated  $_5\phi_4$, $_4\phi_3$ and $_3\phi_2$ series. For instance, for Slater's identity (S.6), we take $\left(\alpha,\beta,a,b,c,d\right)\to\left(1,-1,0,0,q,0\right)$ into \eqref{Formula1} to deduce that
\begin{align}
\left(q;q\right)_{\infty}\sum_{n=0}^{\infty}
\frac{\left(-1;q\right)_n}{\left(q;q\right)^2_n}q^{n^2}
=1+\sum_{n=1}^{\infty}\left(-1\right)^n\left(1+q^n\right)q^{\frac{3n^2-n}{2}}
{_3\phi_2}\left(\begin{array}{ccc}
q^{-n},q^n,-1\\
q,0
\end{array};q,\,q\right).\label{B1}
\end{align}
However, we cannot find a useful method to evaluate the terminated
 $_3\phi_2$ series on the right hand-side of \eqref{B1}.  Settling the difficulties in calculating these kinds of terminated sums would be an interesting thing.
\end{section}
\section{Acknowledgment}
This work is supported by the National Natural Science Foundation of China (Grant 12371328).
The authors  would like to  express gratitude to Professor Zhiguo Liu for his helpful comments during the writing of this manuscript.
\begin{thebibliography}{99}
\bibitem{Andrews_1973}G.E. Andrews, On the $q$-analog of Kummer's theorem and applications.
Duke Math. J. \textbf{40} (1973), 525--528.
\bibitem{Andrews1}G.E. Andrews, On $q$-analogues of the Watson and Whipple summations, SIAM J. Math. Anal. \textbf{7} (1976), 332--336.

\bibitem{Andrews}G.E. Andrews, Connection coefficient problems and partitions, in: D. Ray-Chaudhuri (Ed.), Proc. Sympos. Pure Math., \textbf{34} (1979), 1--24.
    \bibitem{Andrews2}G.E. Andrews, Bressoud polynomials, Rogers-Ramanujan type identities, and applications, 
    Ramanujan J., \textbf{41} (2016), 287--304.
 \bibitem{Ramanujan}G.E. Andrews and B.C. Berndt, Ramanujan's Lost Notebook, Part \uppercase\expandafter{\romannumeral2}, Springer, 2009.
     \bibitem{Bailey}W.N. Bailey, Identities of the Rogers-Ramanujan type, Proc. London Math. Soc. \textbf{50} (1948), 1--10.
     \bibitem{Bowman}D. Bowman, J. McLaughlin and A.V. Sills, Some more identities of Rogers-Ramanujan type, Ramanujan J., \textbf{18} (2009), 307--325.
\bibitem{Wang_Chen}D. Chen and L. Wang, Representations of mock theta functions, Adv. Math., \textbf{365} (2020), Art. 107307.
\bibitem{Gasper_and_Rahman}
G. Gasper and M. Rahman,
Basic Hypergeometric Series, second edition,
Encyclopedia of Mathematics and Its Applications 96, Cambridge
University Press, Cambridge, 2004.
\bibitem{Gessel}I. Gessel and D. Stanton, Applications of $q$-Lagrange inversion to basic hypergeometric series, Trans. Amer. Math. Soc., \textbf{277} (1983), 173--201.
%\bibitem{Hardy}G.H. Hardy, Lectures by Godfrey H. Hardy on te mathematical work of Ramanujan, Edwards Brothers, Ann Arbor, Michigan, 1937, Fall Term 1936. Notes taken by Marshall Hall at the Institute For Advanced Study, Princeton, NJ.
\bibitem{Harsh}H.V. Harsh, Y.S. Kim, M.A. Rakha and A.K. Rathie, A study of $q$-contiguous function relations, Commun. Korean Math. Soc., \textbf{31} (2016), 65--94.
  % \bibitem{Liu_2002}Z.-G. Liu, An expansion formula for $q$-series and applications, Ramanujan J., \textbf{6} (2002), 429--447.
    \bibitem{Liu2}Z.-G. Liu, A $q$-series expansion formula and the Askey-Wilson polynomials, Ramanujan J., \textbf{30} (2013), 193--210.
\bibitem{Liu_1}Z.-G. Liu, On the $q$-derivative and $q$-series expansions, Int. J. Number Theory, \textbf{9} (2013), 2069--2089.
    \bibitem{Liu_4}Z.-G. Liu, On the $q$-partial differential equations and $q$-series, the Legacy of Srinivasa Ramanujan, Ramanujan Math. Soc. Lect. Notes Ser., vol. 20, Ramanujan Math. Soc., Mysore, 2013, 213--250.
        \bibitem{MacMahon}P.A. MacMahon, Combinatory Analysis, vol. \MakeUppercase{\romannumeral 2}, Cambridge University Press, 1918.
             \bibitem{McLaughlin-2008}J. McLaughlin, A.V. Sills and P. Zimmer, Rogers-Ramanujan-Slater type identities, Electron. J. Combin. \textbf{15} (2008),  \#DS15, 59 pp.
        \bibitem{McLaughlin_2008}J. McLaughlin and A.V. Sills, Ramanujan-Slater type identities related to the moduli $18$ and $24$, J. Math. Anal. Appl., \textbf{244} (2008), 765--777.
        \bibitem{McLaughlin}J. McLaughlin, A.V. Sills and P. Zimmer, Rogers-Ramanujan computer searches, J. Symbolic Comput. \textbf{44} (2009), 5077--5091.
 \bibitem{Rogers}L.J. Rogers, Second memoir on the expansion of certain infinite products, Proc. London Math. Soc., \textbf{25} (1894), 318--343.
     \bibitem{Schur}I. Schur, Ein Beitrag zur additiven Zahlentheorie und zur Theorie der kettenbr\"{u}che, S.-B. Preuss. Akad. Wiss. Phys. Math. Klasse (1917), 302--321.
\bibitem{Slater}L.J. Slater, Further identities of the Rogers-Ramanujan type, Proc. London Math. Soc., \textbf{54} (1952), 147--167.
    \bibitem{Hari}H.M. Srivastava, M.P. Chaudhary and F.K. Wakene, A family of theta-function identities based upon $q$-binomial theorem and Heine's transformations, Montes Taurus J. Pure Appl. Math., \textbf{2} (2020), 1--6.
    \bibitem{Verma}A. Verma and V.K. Jain, Transformations between basic hypergeometric series on different bases and identities of Rogers-Ramanujan type, J. Math Anal. Appl., \textbf{76} (1980), 230--269.
        \bibitem{WangChun}C. Wang and S. Chern, Some $q$-transformation formulas and Hecke type identities, Int. J. Number Theory, \textbf{15} (2019), 1349--1367.
        \bibitem{Wang}L. Wang, New proofs of Ramanujan's identities on false theta functions, Ramanujan J., \textbf{50} (2019), 423--431.
         \bibitem{Wang_Yee}L. Wang and A.J. Yee, Some Hecke-Rogers type identities, Adv. Math., \textbf{349} (2019), 733--748.
\end{thebibliography}
\end{document}
